% !TeX program = lualatex
\documentclass[11pt,oneside,openany]{book}
\usepackage[paperwidth=184mm,paperheight=254mm,inner=22mm,outer=22mm,top=23mm,bottom=24mm,headheight=16pt,footskip=12mm]{geometry}
\usepackage{amsmath,amsthm}
\usepackage{libertinus}
\AtBeginDocument{\renewcommand{\Box}{\mathord{\mdlgwhtsquare}}}
\usepackage[final]{microtype}
\usepackage{xcolor,tikz,graphicx,array,booktabs,tabularx,enumitem,longtable,multicol,needspace}
\usetikzlibrary{arrows.meta,positioning,calc,decorations.pathreplacing}
\usepackage{titlesec,fancyhdr}
\definecolor{ink}{HTML}{182E40}\definecolor{pine}{HTML}{126362}\definecolor{lightpine}{HTML}{E3F1EC}\definecolor{brick}{HTML}{984628}\definecolor{slate}{HTML}{435B68}
\usepackage[unicode,pdfdisplaydoctitle,bookmarksnumbered,bookmarksopen=false,colorlinks,linkcolor=pine,citecolor=pine,urlcolor=pine]{hyperref}
\usepackage{bookmark}
\hypersetup{pdftitle={The Mathematics of Fields and Interfaces},pdfsubject={Topology, geometry, analysis and quantum algebras: a companion to Lorentzian Sewing},bookmarksdepth=2}
\setcounter{tocdepth}{1}\setcounter{secnumdepth}{2}
\setlength{\emergencystretch}{2.2em}\raggedbottom
\setlist[enumerate]{leftmargin=1.8em,itemsep=.2em,topsep=.45em}\setlist[itemize]{leftmargin=1.45em,itemsep=.2em,topsep=.45em}
\titleformat{\chapter}[display]{\normalfont\huge\bfseries\color{ink}}{\sffamily\small\bfseries\color{pine}\MakeUppercase{\chaptertitlename}\enspace\thechapter}{6pt}{}
\titlespacing*{\chapter}{0pt}{10pt}{17pt}
\titleformat{\section}{\Large\bfseries\color{ink}}{\thesection}{.6em}{}
\titleformat{\subsection}{\large\bfseries\color{pine}}{\thesubsection}{.55em}{}
\titlespacing*{\section}{0pt}{2.2ex plus .55ex minus .2ex}{.8ex}
\pagestyle{fancy}\fancyhf{}\fancyhead[L]{\small\sffamily\color{slate}The Mathematics of Fields and Interfaces}\fancyhead[R]{\small\sffamily\color{slate}\nouppercase{\leftmark}}\fancyfoot[R]{\thepage}\renewcommand{\headrulewidth}{.3pt}
\renewcommand{\chaptermark}[1]{\markboth{Chapter \thechapter}{}}\fancypagestyle{plain}{\fancyhf{}\fancyfoot[R]{\thepage}\renewcommand{\headrulewidth}{0pt}}
\theoremstyle{plain}\newtheorem{theorem}{Theorem}[chapter]\newtheorem{proposition}[theorem]{Proposition}\newtheorem{lemma}[theorem]{Lemma}\newtheorem{corollary}[theorem]{Corollary}
\theoremstyle{definition}\newtheorem{definition}[theorem]{Definition}\newtheorem{example}[theorem]{Example}\newtheorem{exercise}{Exercise}[chapter]\newtheorem{remark}[theorem]{Remark}
\providecommand{\dd}{\mathrm{d}}\providecommand{\supp}{\operatorname{supp}}\providecommand{\WF}{\operatorname{WF}}
\begin{document}
\frontmatter
\hypersetup{pageanchor=false}
\begin{titlepage}\thispagestyle{empty}\setlength{\parindent}{0pt}
\vspace*{18mm}{\large\sffamily\color{pine}MATHEMATICAL PHYSICS\par}\vspace{9mm}
{\fontsize{35}{40}\selectfont\bfseries\color{ink}The Mathematics of\\Fields and Interfaces\par}\vspace{7mm}
{\Large Topology, geometry, analysis and quantum algebras\par}\vspace{5mm}{\large\itshape A mathematical companion to Lorentzian Sewing\par}
\vfill
\begin{center}\begin{tikzpicture}[x=7.7cm,y=1.3cm,>=Stealth,font=\small]
\draw[->,slate!65] (-.04,0)--(1.09,0) node[right] {$X/L$};
\draw[->,slate!65] (0,-1.02)--(0,1.62) node[above] {$2mK_H$};
\draw[dotted,slate!65] (1,-.88)--(1,1.44);
\draw[very thick,pine,domain=0:1,samples=64] plot (\x,{(exp(-2*\x)+exp(-2*(1-\x)))/(1-exp(-2))}) node[right] {$H=+1$};
\draw[very thick,brick,domain=0:1,samples=64] plot (\x,{(exp(-2*\x)-exp(-2*(1-\x)))/(1+exp(-2))}) node[right] {$H=-1$};
\node[anchor=north west,text=slate] at (.05,-.98) {\sffamily Static response on a circle, $mL=2$};
\end{tikzpicture}\end{center}\vfill
{\small\sffamily\color{slate}Graduate text\enspace/\enspace English edition\par}
\end{titlepage}
\hypersetup{pageanchor=true}
% !TeX root = main.tex
\chapter*{Preface}\addcontentsline{toc}{chapter}{Preface}\markboth{Preface}{}
A solution of the wave equation is sometimes written as a smooth function, sometimes as a point in a phase space, sometimes as a label for a quantum observable. These descriptions cannot simply be substituted for one another. A weak solution may not even have a value at a particular point. A label for an observable can be changed by the equation of motion without changing the observable. A commutator does not assign that observable a vacuum expectation value. How does one pass from each object to the next, and what survives the passage?

This book develops the mathematics of those passages. It can be read independently of \emph{Lorentzian Sewing: Geometry, Dynamics, and Observables}. The reader is assumed to know sets and functions, elementary linear algebra over the real and complex numbers, and calculus in Euclidean space, including change of variables and integration by parts. Familiarity with waves and an action principle will make some applications recognizable, but neither a course in differential geometry nor one in functional analysis is assumed. The inverse and implicit function theorems in Euclidean space, local existence and smooth dependence for smooth ordinary differential equations, and the elementary results of Lebesgue integration are available as analysis inputs. Where we require a deeper theorem we state the version being used and identify its source. A proof in a flat example is not a proof of that theorem on an arbitrary curved space.

Each change of mathematical object answers a question. Charts give local formulas for a form; integration must return a number independent of those charts. Field values live in different fibers; a connection makes their derivatives meaningful. When smoothness is weakened, even restriction to a boundary becomes a theorem with a sharp threshold. Passing from the classical bracket to a quantum algebra leaves two further questions open: which states are positive, and which singular fields may be multiplied?

Each part develops its own arguments before applying them to interfaces. The exercises ask, among other things, for counterexamples: a definition is often understood for the first time upon seeing which of its hypotheses one cannot drop. The last chapter gives a second way into the book for a reader with experience: work through one field problem, and follow the references back at the first operation one cannot justify. A guide at the back locates the mathematics needed for each chapter of \emph{Lorentzian Sewing}.

\paragraph{Conventions.} A dimension counts spacetime coordinates unless another space is named. The metric signature is $(-,+,\ldots,+)$. The letter $\dd$ differentiates on a base manifold; $\delta$ differentiates with respect to fields. A compactly supported field variation need not vanish at an artificial internal cut. We will distinguish ordinary exterior boundaries, where a variational problem has to be specified, from a cut introduced only to describe the same object using smaller regions.
\tableofcontents
\mainmatter
\part{Spaces and smooth geometry}
% !TeX root = ../main.tex
\chapter{Topological spaces and continuous maps}\label{f:topology}

A point moving steadily around a dial returns to its starting position without a jump. Its numerical reading, expressed as a fraction of a turn in $[0,1)$, does jump: readings close to $1$ are followed by $0$. The difference is already visible without a law of motion. Every small arc about the starting position contains readings from both ends of the interval, whereas a small interval about $0$ in $[0,1)$ need contain none near $1$. To decide whether a map is continuous, we must know what counts as a neighborhood on each side of the map.

\section{Open sets and neighborhoods}

In $\mathbb R^d$, a condition on a point is called open if, whenever it holds at $x$, it continues to hold throughout some ball $B(x,r)=\{y:\|y-x\|<r\}$ of positive radius. A strict inequality such as $x_1>0$ in the plane gives an open condition. If two open conditions hold, the smaller of the two radii works for both. There need not be a single positive radius that works for infinitely many conditions: on the line, the intersection of the intervals $(-1/n,1/n)$ over all positive integers $n$ is just $\{0\}$. On the other hand, to satisfy at least one of any collection of open conditions it suffices to stay inside the ball for one that holds. The following axioms retain these rules without making a choice of distance.

\begin{definition}\label{f:top-open}
A \emph{topology} on a set $X$ is a collection $\mathcal T$ of subsets of $X$ containing $\varnothing$ and $X$, closed under arbitrary unions and finite intersections. The pair $(X,\mathcal T)$ is a \emph{topological space}; members of $\mathcal T$ are \emph{open}, and their complements in $X$ are \emph{closed}. A \emph{neighborhood} of $x\in X$ is a set containing an open set that contains $x$. A neighborhood need not itself be open.
\end{definition}

The ball condition above defines the \emph{usual topology} of Euclidean space. It is not the only topology on the underlying set. In the \emph{discrete topology} every subset is open; specifying a single point is an open condition. At the opposite extreme, in the \emph{indiscrete topology} only the empty set and the whole space are open. Nor do ``open'' and ``closed'' divide subsets into two mutually exclusive kinds. In every topology the whole space is both. In the usual line $[0,1)$ is neither: neither it nor its complement contains an open interval about each of its points.

Given $A\subseteq X$, its \emph{interior} $\operatorname{int}_X A$ is the union of all open sets in $X$ contained in $A$. Its \emph{closure} $\overline A$ in $X$ is the intersection of all closed sets in $X$ containing $A$. By taking complements in the topology axioms, arbitrary intersections of closed sets are closed. Thus these constructions really give the largest open set inside $A$ and the smallest closed set containing $A$.

\begin{proposition}[Neighborhood tests]\label{f:top-closure}
A point lies in $\operatorname{int}_X A$ precisely when it has a neighborhood contained in $A$. It lies in $\overline A$ precisely when every neighborhood of it meets $A$.
\end{proposition}
\begin{proof}
The first assertion is the union defining the interior, together with the definition of neighborhood. For the second, a point fails to lie in $\overline A$ just when some closed set containing $A$ misses the point. The complement of that closed set is an open neighborhood disjoint from $A$. Conversely, the complement of any such open neighborhood is a closed set containing $A$.
\end{proof}

For instance, the rational numbers have empty interior in the usual line, since every interval of positive length contains an irrational number. Their closure is all of $\mathbb R$, since every such interval also contains a rational number. A subset whose closure is the whole space is called \emph{dense}; density does not require the subset to contain a single open interval.

\section{Bases and limits}

To specify the usual topology of the line it is unnecessary to list all open sets. It suffices to give the open intervals: each point of an open set lies in some interval contained in the set. Such intervals are called basic open sets. The test for a proposed collection of basic sets works even before a topology has been chosen.

\begin{proposition}[Basis criterion]\label{f:top-basis}
Suppose a collection $\mathcal B$ covers $X$, and whenever $x\in B\cap C$ with $B,C\in\mathcal B$, there is a $D\in\mathcal B$ with $x\in D\subseteq B\cap C$. Then the unions of members of $\mathcal B$ form a topology on $X$. A collection of open sets in a given topology is a \emph{basis} for it exactly when every point of every open set lies in a member contained in that open set.
\end{proposition}
\begin{proof}
The empty union is empty, the union of the cover is $X$, and a union of unions is again a union of members. If $x$ belongs to two such unions, choose a member from each containing $x$. The hypothesis supplies a basic set about $x$ inside their intersection. Taking the union of these smaller sets proves that the intersection is open. This also explains the last criterion: the indicated basic sets recover each open set as their union.
\end{proof}

For a positive integer $d$, even the boxes in $\mathbb R^d$ of the form $\prod_{j=1}^d(a_j,b_j)$ with \emph{rational} endpoints form a basis. Each box is open: the minimum of the distances from a point in it to its faces is positive. Conversely, inside any ball about $x$ of radius $r$, one can put a rational box about $x$ by choosing
\[
 x_j-r/\sqrt d<a_j<x_j<b_j<x_j+r/\sqrt d
\]
in each coordinate. Then the squared distance to $x$ of any point in the box is less than $d(r/\sqrt d)^2=r^2$. There are only countably many such boxes, since each is described by a fixed finite number of rational endpoints. Here and below, countable collections include finite ones.

A space with a countable basis is \emph{second countable}. There is also a pointwise notion. A \emph{local basis} at $x$ is a collection of open neighborhoods of $x$ such that every neighborhood of $x$ contains one of them. A space with a countable local basis at each point is \emph{first countable}. Second countability implies first countability: retain the basis elements containing the point. The converse fails on an uncountable discrete set. Each point has its singleton as a local basis element, but a basis for the whole space must include every singleton.

A sequence $(x_n)$ \emph{converges} to $x$ if every neighborhood of $x$ contains $x_n$ for all sufficiently large $n$. In Euclidean space this says exactly that $\|x_n-x\|\to0$, by testing the balls about $x$. A space is \emph{Hausdorff} if every pair of distinct points has disjoint open neighborhoods. Euclidean space is Hausdorff: use balls of radius one third of the distance between the two points. A sequence in a Hausdorff space cannot have two different limits, because its tail cannot lie in two disjoint neighborhoods. At the opposite extreme, in an indiscrete space every sequence converges to every point.

\begin{proposition}\label{f:top-sequence-closure}
A limit of a sequence in $A$ lies in $\overline A$. Conversely, if $x\in\overline A$ has a countable local basis, there is a sequence in $A$ converging to $x$.
\end{proposition}
\begin{proof}
Every neighborhood of a limit meets the sequence and hence meets $A$. For the converse, write the local basis as $U_1,U_2,\ldots$, allowing repetitions. Choose $a_n\in A\cap U_1\cap\cdots\cap U_n$, which is possible because the finite intersection of the $U_j$ is an open neighborhood of $x$. Any neighborhood of $x$ contains some $U_j$; then it contains $a_n$ for all $n\geq j$.
\end{proof}

The converse is genuinely conditional. Give an uncountable set $X$ the topology consisting of $\varnothing$ and the sets with countable complement. This is a topology: finite unions of countable complements are countable, and any union containing one nonempty open set already has countable complement. For $p\in X$, every neighborhood of $p$ meets $A=X\setminus\{p\}$. Yet the range of any sequence in $A$ is countable; its complement is an open neighborhood of $p$ missing \emph{every} term. Here closure cannot be recovered from sequences alone.

\section{Continuity and homeomorphisms}

Consider a function $f:X\to Y$ and an open set $V\subseteq Y$ of acceptable outputs. The inputs producing those outputs form the \emph{inverse image} $f^{-1}(V)=\{x\in X:f(x)\in V\}$. This notation does not require $f$ to have an inverse function. If each acceptable input has a neighborhood of acceptable inputs, their set is open. This gives a definition in which distances no longer appear.

\begin{definition}\label{f:top-continuous-def}
A map $f:X\to Y$ between topological spaces is \emph{continuous} if $f^{-1}(V)$ is open in $X$ for every open $V$ in $Y$.
\end{definition}

\begin{proposition}[Tests for continuity]\label{f:top-continuity}
Continuity is equivalent to the following local condition: for every $x\in X$ and every open $V$ containing $f(x)$ there is an open $U$ containing $x$ with $f(U)\subseteq V$. In the inverse-image definition it is enough to test members of a basis for $Y$. One may instead test that the inverse image of each closed set is closed.
\end{proposition}
\begin{proof}
For a continuous map take $U=f^{-1}(V)$. Conversely, the local condition writes $f^{-1}(V)$ as a union of open neighborhoods of its points. The basis test works because inverse images preserve arbitrary unions. The closed-set test follows from $f^{-1}(Y\setminus C)=X\setminus f^{-1}(C)$.
\end{proof}

In Euclidean spaces, take $V=B(f(x),\varepsilon)$. A suitable $U$ contains a ball $B(x,\delta)$, so $\|z-x\|<\delta$ implies $\|f(z)-f(x)\|<\varepsilon$: this is the calculus test. A map that sends open sets to open sets is called an \emph{open map}, a different condition. For instance, $t\mapsto t^2$ on the line is continuous but sends $(-1,1)$ to $[0,1)$, which is not open.

Continuity preserves limits of sequences: if $x_n\to x$, the terms eventually lie in $f^{-1}(V)$ for each open neighborhood $V$ of $f(x)$. On a first countable domain this is also a test for continuity. Indeed, if the local condition fails at $x$ for some $V$, enumerate a local basis $U_1,U_2,\ldots$ there. No intersection $U_1\cap\cdots\cap U_n$ is contained in $f^{-1}(V)$; choose $x_n$ in that intersection but outside $f^{-1}(V)$. Then $x_n\to x$ while $f(x_n)$ never belongs to $V$.

If $f:X\to Y$ and $g:Y\to Z$ are continuous, so is $g\circ f$, because $(g\circ f)^{-1}(W)=f^{-1}(g^{-1}(W))$. Identity maps are continuous. Thus two spaces can serve as interchangeable models if there is a bijection between them that is continuous \emph{in both directions}. Such a bijection is called a \emph{homeomorphism}. For a bijection $f$ this requires, in addition to continuity, that $f$ send open sets to open sets: those images are the inverse images under the inverse function. The dial will give an example where only one direction works.

\section{Subspaces, products, and disjoint unions}

A curve drawn in the plane inherits neighborhoods by cutting plane neighborhoods down to the curve. More generally, the \emph{subspace topology} on $A\subseteq X$ has as its open sets the intersections $A\cap U$ with $U$ open in $X$. The topology axioms follow by intersecting unions and finite intersections with $A$. In $A=[0,1)$, for instance, $[0,1/2)$ is open because it is $A\cap(-1,1/2)$, not because it is open in the line. A subspace of a Hausdorff space is Hausdorff: intersect separating neighborhoods with the subspace.

A different construction records a pair of values. The \emph{product topology} on $X\times Y$ has basis the rectangles $U\times V$ with $U$ open in $X$ and $V$ open in $Y$. These cover the product and the intersection of two is another rectangle of this kind. For two copies of the usual line this gives the usual plane topology by the box criterion above.

In contrast, when a value belongs to \emph{one} of a family of alternatives, we form the \emph{disjoint union} $D=\coprod_{i\in I}X_i$. Its elements are tagged pairs $(i,x)$ with $x\in X_i$, so even originally identical copies remain distinct. With $j_i(x)=(i,x)$, declare $O\subseteq D$ open if $j_i^{-1}(O)$ is open in $X_i$ for every $i$. This is a topology since inverse images preserve unions and intersections. Each single copy is open, and its complement, being the union of the other copies, is open too.

\begin{proposition}[Mapping properties]\label{f:top-universal}
If $k:Z\to X$ takes values in $A\subseteq X$, there is a unique $h:Z\to A$ with $k=i\circ h$, where $i:A\hookrightarrow X$ is inclusion. The map $h$ is continuous exactly when $k$ is. A pair of maps $u:Z\to X$, $v:Z\to Y$ defines a unique map into the product with those coordinates, and it is continuous exactly when both coordinates are continuous. Maps $g_i:X_i\to Z$ define a unique map $g:D\to Z$ with $g\circ j_i=g_i$, and it is continuous exactly when every $g_i$ is continuous.
\end{proposition}
\begin{proof}
For the subspace, $h(z)$ must be $k(z)$ considered as a point of $A$. Then $h^{-1}(A\cap U)=k^{-1}(U)$, which tests all subspace opens.

For the product the only possible map is $h(z)=(u(z),v(z))$. Its inverse image of a basic rectangle is
\[
 h^{-1}(U\times V)=u^{-1}(U)\cap v^{-1}(V),
\]
which is open if the coordinates are continuous. Conversely the projections are continuous, as seen by taking rectangles $U\times Y$ and $X\times V$; compose with them to recover the coordinates.

On the disjoint union set $g(i,x)=g_i(x)$. For an open $W\subseteq Z$, the part of $g^{-1}(W)$ in copy $i$ pulls back to $g_i^{-1}(W)$, exactly the test in the definition of the topology of $D$.
\end{proof}

\begin{example}[One turn around the circle]\label{f:top-turn}
Give $S^1=\{(u,v)\in\mathbb R^2:u^2+v^2=1\}$ the subspace topology. The map
\[
 p(t)=(\cos(2\pi t),\sin(2\pi t))
\]
from the line into $S^1$ is continuous: its two real coordinates are continuous and its image lies on the circle. Its restriction to $[0,1)$ is a continuous bijection. But for $t_n=1-1/n$, the points $p(t_n)$ converge to $p(0)$ while the $t_n$ do not converge to $0$ in $[0,1)$: eventually they lie outside $[0,1/2)$. The inverse of this bijection is not continuous.
\end{example}

The disjoint union allows completely independent choices of maps on its copies. For pieces of a \emph{single} space, choices must agree where they overlap, and the topology of the pieces matters as well.

\begin{proposition}[Pasting maps]\label{f:top-pasting}
Let $X$ be covered by subspaces $A_i$, and suppose continuous maps $f_i:A_i\to Y$ agree on pairwise intersections. They define a unique map $f:X\to Y$. It is continuous if all $A_i$ are open, or if the cover is finite and all $A_i$ are closed.
\end{proposition}
\begin{proof}
Agreement gives a well-defined value $f(x)=f_i(x)$ in any piece containing $x$. In the open-cover case, for open $V\subseteq Y$ each $f_i^{-1}(V)$ is open in $A_i$ and hence in $X$, so their union $f^{-1}(V)$ is open. In the finite closed-cover case apply the same argument to a closed subset of $Y$: each inverse image in a piece is closed in $X$, and now their union is finite.
\end{proof}

\section{Exercises}

\begin{exercise}\label{f:top-ex-relative}
Take $Y=[0,2)$ and $A=[0,1)\cup(3/2,2)$. Find the interior and the closure of $A$ first in $\mathbb R$, then in $Y$. Which endpoint changes its status in each calculation?
\end{exercise}
\begin{exercise}\label{f:top-ex-indicator}
Give $T=\{0,1\}$ the topology $\{\varnothing,\{1\},T\}$. For $A\subseteq X$, let $\chi_A:X\to T$ take value $1$ on $A$ and $0$ elsewhere. Characterize continuity of $\chi_A$. How does the answer change if $T$ is discrete?
\end{exercise}
\begin{exercise}\label{f:top-ex-basis}
Use the intervals $[a,b)$, $a<b$, as a proposed basis on $\mathbb R$. Prove that they give a topology strictly finer than the usual one (that is, with strictly more open sets). In this topology which of the sequences $1/n$ and $-1/n$ converges to $0$?
\end{exercise}
\begin{exercise}\label{f:top-ex-diagonal}
Show that $X$ is Hausdorff if and only if the diagonal $\Delta=\{(x,x):x\in X\}$ is closed in $X\times X$. Deduce that for continuous $f,g:Z\to X$ with $X$ Hausdorff, the set on which $f=g$ is closed in $Z$.
\end{exercise}

\subsection*{Solutions and hints}
\noindent\emph{Exercise~\ref{f:top-ex-relative}.}
In the line the answers are $(0,1)\cup(3/2,2)$ and $[0,1]\cup[3/2,2]$. In $Y$ the set $A$ is itself open, and its closure is $[0,1]\cup[3/2,2)$. Neighborhoods of $0$ have lost their negative side; $2$ is not a candidate point for closure in $Y$.

\smallskip\noindent\emph{Exercise~\ref{f:top-ex-indicator}.}
The only nontrivial open set to test is $\{1\}$, whose inverse image is $A$. Thus continuity says precisely that $A$ is open. With the discrete topology one must also test $\{0\}$, so $A$ must be both open and closed.

\smallskip\noindent\emph{Exercise~\ref{f:top-ex-basis}.}
Two proposed basic intervals meeting at $x$ contain $[x,c)$, where $c$ is the smaller right endpoint. Every usual open interval contains such a neighborhood of each of its points, but $[0,1)$ is not open in the usual topology. Every basic neighborhood of $0$ contains $[0,\varepsilon)$ for some $\varepsilon>0$, hence eventually contains $1/n$. The sequence $-1/n$ never enters the neighborhood $[0,1)$.

\smallskip\noindent\emph{Exercise~\ref{f:top-ex-diagonal}.}
A rectangle $U\times V$ misses $\Delta$ exactly when $U\cap V=\varnothing$. Apply the rectangle basis at each pair $(x,y)$ with $x\ne y$ to see that openness of the complement is precisely the Hausdorff condition. The set in the second assertion is the inverse image of $\Delta$ under $z\mapsto(f(z),g(z))$.
% !TeX root = ../main.tex
\chapter{Compactness, connectedness and identifications}\label{f:quotient}

Example~\ref{f:top-turn} gave a continuous bijection onto the circle whose inverse was not continuous. There is a better way to parametrize a turn: keep \emph{both} endpoints of $[0,1]$ and regard them as one point. To prove that this gives the right topology, we shall need a way to recognize a homeomorphism without constructing its inverse. The relevant property of the closed interval is not its length but a consequence of completeness: whenever open sets cover it, finitely many of those very sets already suffice.

\section{Finite subcovers}

\begin{definition}\label{f:q-compact}
An \emph{open cover} of $K\subseteq X$ is a family of open sets $U_i$ in $X$ with $K\subseteq\bigcup_{i\in I}U_i$. The subset $K$ is \emph{compact} if every open cover has a \emph{finite subcover}: there exist finitely many indices $i_1,\ldots,i_m$ from the given family with $K\subseteq U_{i_1}\cup\cdots\cup U_{i_m}$. For a space take $K=X$.
\end{definition}

This is a property of the subspace $K$: intersecting the cover members with $K$ gives a cover by subspace opens, and every such open is the intersection with an ambient open. The quantifier on covers is important. A single cover without a finite subcover disproves compactness. The intervals $[0,1-1/n)$, $n\geq2$, are open in $[0,1)$ and cover it; any finite selection stops short of $1$ and misses part of $[0,1)$.

\begin{example}[Closed boxes]\label{f:q-box}
A closed bounded box in $\mathbb R^d$ is compact. Suppose a particular open cover of it has no finite subcover. Bisect every side; some resulting box has no finite subcover from the same family, for otherwise we could combine finitely many finite covers. Repeat inside that box. In each coordinate the resulting nested closed intervals have lengths tending to zero and hence a common point by completeness of the real line. These coordinates give a point $x$ in every selected box. One cover member contains $x$, so it contains a ball about $x$. Eventually the diameter of a selected box is smaller than the radius of that ball. The entire box is then covered by that single member, a contradiction.
\end{example}

\begin{theorem}\label{f:q-compact-hausdorff}
A closed subset of a compact space is compact. The continuous image of a compact space is compact. In a Hausdorff space every compact subset is closed. Consequently a continuous map from a compact space to a Hausdorff space sends closed sets to closed sets; if it is bijective, it is a homeomorphism.
\end{theorem}
\begin{proof}
For a closed subset $F$ of a compact $K$, adjoin $K\setminus F$ to any cover of $F$ by opens in $K$. A finite subcover of $K$ supplies a finite cover of $F$ after the added set is discarded. For a continuous map $f:K\to Y$, pull back an open cover of $f(K)$. The inverse images cover $K$; the indices of a finite subcover give a finite subcover of the image.

Now let $K\subseteq Y$ be nonempty and compact, $Y$ Hausdorff, and $y\notin K$. For each $z\in K$ choose disjoint open neighborhoods $U_z$ of $z$ and $V_z$ of $y$. Choose $U_{z_1},\ldots,U_{z_m}$ covering $K$. Then
\[
 U=\bigcup_{j=1}^m U_{z_j},\qquad V=\bigcap_{j=1}^m V_{z_j}
\]
are disjoint open neighborhoods of $K$ and $y$, respectively. Thus every point outside $K$ is interior to its complement. The empty set is already closed. Finally, if $f:X\to Y$ with $X$ compact, $Y$ Hausdorff and $f$ continuous, each closed $F\subseteq X$ is compact, so $f(F)$ is compact and hence closed. If $f$ is bijective this is the closed-set test for continuity of its inverse.
\end{proof}

In $\mathbb R^d$ it follows that compactness is equivalent to being closed and bounded (the Heine--Borel theorem). A closed bounded set is closed in a large compact box. Conversely, a compact set is closed by the Hausdorff property, and finitely many of the balls $B(0,n)$ must cover it, so it is bounded. The closedness conclusion is not intrinsic to compactness alone: in a two-point indiscrete space either singleton is compact but neither is closed. On the other hand the usual circle \emph{is} compact, since the continuous map $p(t)=(\cos(2\pi t),\sin(2\pi t))$ takes the compact interval $[0,1]$ onto it.

\section{Connected sets, paths, and components}

A continuous real function on an interval cannot pass from a negative value to a positive one without taking the value zero. To express the underlying obstruction without having to order the points of a space, consider a proposed division of its points into two parts. If membership in either part is an open condition, the division can be detected locally everywhere, without any point at which the two parts meet.

\begin{definition}\label{f:q-connected}
A \emph{separation} of a space $X$ is a pair of disjoint nonempty open sets $U,V$ whose union is $X$. A space admitting no separation is \emph{connected}. Equivalently, $X$ is connected if and only if its only subsets that are both open and closed are $\varnothing$ and $X$. A subset is connected if it is connected in its subspace topology.
\end{definition}

\begin{proposition}\label{f:q-connected-image}
The connected subsets of the usual line are precisely the intervals, counting singletons and the empty set as intervals. The continuous image of a connected space is connected.
\end{proposition}
\begin{proof}
If a subset contains $a<b$ but misses a number $c$ between them, the open rays on either side of $c$ separate it. Conversely, a separation of an interval would give a real function on it equal to $0$ on one part and $1$ on the other. The function is continuous by pasting its constant restrictions on the two open parts. This contradicts the intermediate value theorem on the closed segment between points in the two parts. For the image assertion, a separation of the image, using its subspace topology, pulls back to disjoint opens covering the domain. Both are nonempty because both parts contain image points.
\end{proof}

\begin{proposition}[Connected components]\label{f:q-components}
Through each point $x$ passes a unique maximal connected subset, its \emph{connected component}. The components partition the space and are closed.
\end{proposition}
\begin{proof}
Take the union $C_x$ of all connected subsets containing $x$, including $\{x\}$. In a proposed separation of the union, every one of those subsets meets the part containing $x$. If one also met the other part, its intersections with the two parts would separate it. Thus $C_x$ is connected and is the largest possible choice. If two such components meet, their union is connected by the same common-point argument, so maximality makes them equal. Finally the closure of a connected set $C$ is connected. Indeed every nonempty relatively open subset of $\overline C$ meets $C$ by the neighborhood test for closure. A separation of $\overline C$ would therefore restrict to a separation of $C$. Taking the closure cannot enlarge a maximal connected set.
\end{proof}

A component need not be open. For example, an irrational number between two distinct rationals separates those points by disjoint rays in $\mathbb Q$. Thus all components of $\mathbb Q$ are singletons, none open in its usual subspace topology.

A \emph{path} from $x$ to $y$ in $X$ is a continuous map $\gamma:[0,1]\to X$ with $\gamma(0)=x$ and $\gamma(1)=y$. A space is \emph{path connected} if every pair can be joined by a path. Being joined by a path is an equivalence relation: use constant paths, reversal $t\mapsto\gamma(1-t)$, and concatenation. For paths $\gamma,\eta$ with $\gamma(1)=\eta(0)$, use $\gamma(2t)$ on $[0,1/2]$ and $\eta(2t-1)$ on $[1/2,1]$; continuity follows from finite closed pasting (Proposition~\ref{f:top-pasting}). The equivalence classes are the \emph{path components}. Each lies in a connected component since every path image is connected. In particular path connected implies connected, but the converse fails even for compact subsets of the plane (Exercise~\ref{f:q-ex-sine}). On $S^1$ an actual path between $p(a)$ and $p(b)$ is $t\mapsto p((1-t)a+tb)$.

A space is \emph{locally path connected} if every neighborhood of each point contains a path connected open neighborhood of that point. This is a condition about paths inside arbitrarily small sets, not about reaching distant points.

\begin{proposition}\label{f:q-local-path}
In a locally path connected space, the path components are open and coincide with the connected components.
\end{proposition}
\begin{proof}
Around each point of a path component take a path connected open neighborhood. All points of the neighborhood belong to the same path component, so the component is open. Its complement is also open as a union of other path components. A connected component therefore cannot meet two path components, while each path component already lies in a connected component.
\end{proof}

\section{Compact neighborhoods and countable covers}

The line is not compact, but its open intervals can be chosen to have compact closures, as small as desired. The distinction between a compact \emph{space} and a compact \emph{neighborhood} makes it possible to use finite-cover arguments even in an unbounded space. We call a space \emph{locally compact} if each point has a compact neighborhood, that is, a compact subset containing an open neighborhood of the point. In a Hausdorff space this one-neighborhood condition provides arbitrarily small ones.

\begin{proposition}[Shrinking a neighborhood]\label{f:q-shrinking}
In a locally compact Hausdorff space, given an open $O$ and $x\in O$ there is an open $V$ such that $x\in V\subseteq\overline V\subseteq O$ and $\overline V$ is compact. Conversely this property implies local compactness.
\end{proposition}
\begin{proof}
Choose a compact neighborhood $K$ and an open $N$ with $x\in N\subseteq K$. The set $F=K\setminus O$ is closed in $K$ and hence compact. By the separation argument in Theorem~\ref{f:q-compact-hausdorff}, there are disjoint open sets $P$ containing $x$ and $W$ containing $F$ (take $W=\varnothing$ if $F=\varnothing$). Let $V=N\cap P$. Since $K$ is closed and $W$ is open, $\overline V\subseteq K\setminus W\subseteq O$. The closure is closed in the compact $K$, hence compact. Conversely $\overline V$ itself is a compact neighborhood of $x$.
\end{proof}

A different restriction concerns the size of a cover. A space is \emph{Lindel\"of} if every open cover has a countable subcover, finite subcovers included. The existence of a compact neighborhood at each point does not imply this covering property. An uncountable discrete space is locally compact and Hausdorff, yet its cover by singletons has no countable subcover.

\begin{proposition}\label{f:q-lindelof}
Every second countable space is Lindel\"of. Every subspace of a second countable space is second countable and therefore also Lindel\"of.
\end{proposition}
\begin{proof}
Let $\mathcal B$ be a countable basis and fix an open cover. For each nonempty $B\in\mathcal B$ lying in at least one cover member, select one such member. The selected family is at most countable. It covers: any point lies in an original cover member, hence in a basic neighborhood inside that member, and some selected member contains that entire neighborhood. For $A\subseteq X$, intersect the original basis elements with $A$ to obtain a countable subspace basis.
\end{proof}

\begin{corollary}\label{f:q-compact-basis}
A second countable locally compact Hausdorff space has a countable basis of open sets with compact closure. Given $x\in O$ with $O$ open, a member can be chosen about $x$ whose closure is contained in $O$.
\end{corollary}
\begin{proof}
From a countable basis retain only elements with compact closure. Use Proposition~\ref{f:q-shrinking} to take $x\in V\subseteq\overline V\subseteq O$ with compact closure, then take a basis element $x\in B\subseteq V$. Its closure is closed in $\overline V$ and contained in $O$, so $B$ was retained.
\end{proof}

\section{Identifying points}

Let an equivalence relation $\sim$ on $X$ specify which points are to be regarded as the same. Each point $x$ belongs to the class $[x]=\{x'\in X:x'\sim x\}$; by reflexivity, symmetry and transitivity, these classes partition $X$. Write $X/{\sim}$ for the set of classes and $q(x)=[x]$ for the projection. A condition on the classes selects a set of original points: it selects every point of each class satisfying the condition. This suggests testing openness before the identification.

\begin{definition}\label{f:q-quotient-map}
A surjection $q:X\to Y$ is a \emph{quotient map} if $V\subseteq Y$ is open \emph{if and only if} $q^{-1}(V)$ is open in $X$. Given just a surjection from a topological space to a set, this rule defines the \emph{quotient topology} on that set, in particular on $X/{\sim}$.
\end{definition}

It defines a topology because inverse images commute with unions and finite intersections. It is the finest topology on $Y$ for which $q$ is continuous: continuity forbids declaring any other sets open. A continuous surjection need not be quotient, as the bijection in Example~\ref{f:top-turn} shows. A bijective quotient map has continuous inverse by the very definition.

\begin{theorem}[Descent through a quotient]\label{f:q-descent}
For a quotient map $q:X\to Y$, a map $f:X\to Z$ has the form $f=\bar f\circ q$ for a unique map $\bar f:Y\to Z$ exactly when $f$ is constant on each fiber $q^{-1}(\{y\})$. In that case $\bar f$ is continuous if and only if $f$ is continuous.
\end{theorem}
\begin{proof}
A composite through $q$ takes the same value on points with the same $q$-value. Conversely set $\bar f(y)=f(x)$ for any $x$ with $q(x)=y$. Surjectivity supplies such an $x$, constancy makes the value independent of $x$, and the formula forces uniqueness. For open $W\subseteq Z$ one has
\[
 q^{-1}(\bar f^{-1}(W))=(\bar f\circ q)^{-1}(W)=f^{-1}(W).
\]
If $f$ is continuous, the right-hand side is open, so the quotient condition makes $\bar f^{-1}(W)$ open. The other implication follows by composition since $q$ is continuous.
\end{proof}

\begin{example}[The circle by identification]\label{f:q-circle}
Identify just $0$ and $1$ in $[0,1]$ and give the class set $Y$ the quotient topology. The map $p:[0,1]\to S^1$ has exactly these fibers: two parameters in the interval give the same cosine and sine precisely when their difference is an integer. By descent the map
\[
 \bar p:Y\longrightarrow S^1,\qquad [t]\longmapsto(\cos(2\pi t),\sin(2\pi t)),
\]
is a continuous bijection. Its domain is compact as the continuous image of $[0,1]$ under $q$, and $S^1$ is Hausdorff as a subspace of the plane. Theorem~\ref{f:q-compact-hausdorff} makes $\bar p$ a homeomorphism. In particular, a continuous map on the closed interval prescribes a continuous map on the circle precisely when its two endpoint values agree.
\end{example}

The image of an open set under a quotient map need not be open. In the last example the full inverse image of $q([0,1/4))$ is $[0,1/4)\cup\{1\}$, not open at $1$. For any quotient map, $q^{-1}(q(A))$ is called the \emph{saturation} of $A$: it adds every point sharing a fiber with a point already in $A$. A set is \emph{saturated} when it is a union of entire fibers. The open sets downstairs correspond exactly to the saturated open sets upstairs, by taking inverse image in one direction and image in the other.

Many identifications arise from reversible transformations, such as integer translations of the line. Suppose a group $G$ acts on $X$ by homeomorphisms: each $g$ gives a homeomorphism $x\mapsto gx$, with $1x=x$ and $g(hx)=(gh)x$. Declaring $x$ equivalent to $gx$ gives as classes the \emph{orbits} $Gx=\{gx:g\in G\}$, since a move can be undone by its inverse and two successive moves compose. The \emph{orbit space} $X/G$ is the set of orbits with the quotient topology. In this instance the projection is open: the saturation of an open $U$ is $\bigcup_{g\in G}gU$, a union of opens. Openness of the projection does not by itself separate the orbits. For that one must test neighborhoods containing whole fibers.

\begin{proposition}\label{f:q-separated-fibers}
For a quotient map $q:X\to Y$, the space $Y$ is Hausdorff exactly when any two distinct fibers have disjoint saturated open neighborhoods in $X$.
\end{proposition}
\begin{proof}
Disjoint open neighborhoods of two points pull back to such neighborhoods of their fibers. In the other direction, images of saturated open sets are open by the quotient condition, and disjoint saturated sets have disjoint images.
\end{proof}

\begin{example}[The line with two origins]\label{f:q-double-origin}
Start with two disjoint copies $\mathbb R_1,\mathbb R_2$ of the line and identify $(t,1)$ with $(t,2)$ for every $t\ne0$, leaving the two zeros distinct. The original disjoint union is Hausdorff. Even the individual fibers are closed there, being singletons or pairs of points, so each singleton in the quotient is closed. Nevertheless every open neighborhood of the first origin pulls back to a set containing an interval about zero in the first copy. A neighborhood of the second does the same in the second copy. Choose a small nonzero $t$ belonging to both coordinate intervals. It represents a single point of the quotient lying in both neighborhoods. The quotient is not Hausdorff.

Each $q(\mathbb R_i)$ is open: its preimage is that entire copy together with the other copy minus its origin. Moreover $q$ maps the $i$th copy homeomorphically onto $q(\mathbb R_i)$. The inverse there is the ordinary coordinate $[t,j]\mapsto t$, which is continuous by descent since it is continuous on both original copies. Thus every point has an open neighborhood homeomorphic to a line, including each of the two points that cannot be separated.
\end{example}

An open neighborhood $U$ described by a homeomorphism $\varphi:U\to V$ onto an open set of $\mathbb R^d$ is a \emph{coordinate neighborhood}; $\varphi$ is a \emph{chart}. For the circle, projection $(u,v)\mapsto u$ on the upper semicircle $v>0$ has the continuous inverse $s\mapsto(s,\sqrt{1-s^2})$ on $(-1,1)$. Rotating this neighborhood gives charts about the remaining points. In any chart, inverse images of small open balls give path connected neighborhoods: join two points of a ball by the straight segment and carry it back by $\varphi^{-1}$. A closed ball centered at $\varphi(x)$ of positive radius and contained in $V$ likewise gives a compact neighborhood of $x$ by continuity of $\varphi^{-1}$.

The doubled origin shows why one must still ask separately for the Hausdorff property. Countability is a global requirement too: an uncountable disjoint union of lines has charts everywhere but its cover by copies has no countable subcover. In a second countable space, by contrast, Proposition~\ref{f:q-lindelof} selects a countable subcover from any cover by charts.

\section{Exercises}

\begin{exercise}\label{f:q-ex-extrema}
Let $K$ be nonempty and compact and let $f:K\to\mathbb R$ be continuous. Show that $f$ attains a maximum and a minimum. What goes wrong for $f(t)=t$ on $(0,1)$?
\end{exercise}
\begin{exercise}\label{f:q-ex-rationals}
The rational line is Hausdorff and second countable. Show that in its usual topology it has no compact neighborhood of any point. Use the inclusion into the real line to test compactness.
\end{exercise}
\begin{exercise}\label{f:q-ex-orbits}
Compare the actions on $\mathbb R$ by translation by integers and by rationals. Show that $\mathbb R/\mathbb Z$ is homeomorphic to the circle, whereas $\mathbb R/\mathbb Q$ has more than one point and the indiscrete topology.
\end{exercise}
\begin{exercise}\label{f:q-ex-sine}
Let $C=\{(u,\sin(1/u)):0<u\leq1\}$ and take its closure $K$ in the plane. Describe $K$ and show it is compact and connected. Prove that no path in $K$ joins $(0,0)$ to a point of $C$.
\end{exercise}

\subsection*{Solutions and hints}
\noindent\emph{Exercise~\ref{f:q-ex-extrema}.}
The image is a nonempty compact subset of the line, so it is bounded and closed. Its supremum and infimum are finite and belong to its closure by the defining property of each bound. Thus they are in the image. For the open interval the bounds still exist but are not attained: the image is not closed.

\smallskip\noindent\emph{Exercise~\ref{f:q-ex-rationals}.}
A compact neighborhood $K$ of a rational $r$ would contain $(r-\varepsilon,r+\varepsilon)\cap\mathbb Q$ for some $\varepsilon>0$. Its image in $\mathbb R$ would be compact and hence closed. But every irrational in that interval is in the real closure of this image and none is in the image.

\smallskip\noindent\emph{Exercise~\ref{f:q-ex-orbits}.}
For integer translations the circle parametrization $p$ descends to a continuous bijection, and the orbit space is compact because it is already the image of $[0,1]$ under projection. Apply the compact-to-Hausdorff theorem. For rational translations each orbit is dense, so any nonempty open set meets every orbit and its saturation is all of $\mathbb R$. Only the empty set and the whole space are saturated open. The orbits of $0$ and $\sqrt2$ are distinct.

\smallskip\noindent\emph{Exercise~\ref{f:q-ex-sine}.}
The graph lies in the closed box $[0,1]\times[-1,1]$. The only points added are $\{0\}\times[-1,1]$: for $-1\leq v\leq1$ take $a$ with $\sin a=v$ and $u_n=1/(a+2\pi n)$ for sufficiently large $n$. Then $(u_n,\sin(1/u_n))\to(0,v)$. Away from $u=0$ the graph adds no points by continuity. Its closure is closed and bounded, and it is connected since $C$ is the continuous image of an interval. If a path from $(0,0)$ ends in $C$, write its coordinates as $(u(t),v(t))$ and let $b$ be the last time with $u(b)=0$; the zero set is closed in $[0,1]$ and $b<1$. For every $0<\delta\leq1-b$, the coordinate $u$ takes all values from $0$ to $u(b+\delta)>0$ on $[b,b+\delta]$. Among the positive ones there are values with $\sin(1/u)=1$ and values with $\sin(1/u)=-1$. Thus $v$ takes both values arbitrarily close to $b$, contradicting continuity there.
% !TeX root = ../main.tex
\chapter{Smooth manifolds and differentiation}\label{f:smooth}
A circle has small neighborhoods homeomorphic to intervals. That statement lets us test continuity on the circle, but it does not yet say what the velocity of a moving point is. Even on an interval one can change the coordinate $x$ to $x^3$: no open set detects the change, yet a particle described by the inverse change has no finite derivative at the origin. To differentiate we must specify which coordinates may replace one another. We begin with spaces without boundary. What is special about a boundary point will be easier to recognize once differentiation at an ordinary point is understood.

\section{Charts, changes of variables and smooth maps}
A topological manifold of dimension $n$ is a Hausdorff, second-countable space locally homeomorphic to open subsets of $\mathbb R^n$. A \emph{chart} is such a homeomorphism $x:U\to x(U)$, not just an assignment of $n$ numbers. Given charts $x:U\to x(U)$ and $y:V\to y(V)$, the change $y\circ x^{-1}$ is defined on $x(U\cap V)$; it is not generally defined on all of $x(U)$. The charts are smoothly compatible if this map and its inverse are smooth in the sense of Euclidean calculus. An atlas consists of pairwise compatible charts covering the space. Including every chart compatible with this atlas gives its \emph{smooth structure}. By the Euclidean chain rule, two charts both compatible with a covering atlas are compatible with each other: near each common point insert a chart of the atlas between the two changes. Thus there is no further choice in passing to the maximal atlas. A smooth manifold is a topological manifold equipped with this extra structure.

\begin{example}[Reading an atlas on the sphere]\label{f:stereo}
Give $S^2=\{(X,Y,Z):X^2+Y^2+Z^2=1\}$ its subspace topology. The lines through the north pole give coordinates $(u,v)=(X,Y)/(1-Z)$ on its complement. With $r^2=u^2+v^2$ the inverse there is
\[
 (X,Y,Z)=\frac{(2u,2v,r^2-1)}{1+r^2}.
\]
Indeed the vector on the right has squared length one and $1-Z=2/(1+r^2)$. At the south pole the chart gives $(0,0)$; it is the \emph{north} pole which has no coordinates here. The chart from the other pole is $(X,Y)/(1+Z)$. On the overlap the change is $(u,v)\mapsto (u,v)/r^2$, a smooth involution of the punctured plane. These two charts cover the sphere. No global coordinate is required. In fact no single planar chart can cover $S^2$: its image would be a nonempty compact open subset of $\mathbb R^2$, which does not exist. The last argument is topological; the compatibility calculation is smooth.
\end{example}

For a map $F:M\to N$, assume first that $F$ is continuous. In charts near $p$ and $F(p)$ its representative is $y\circ F\circ x^{-1}$, restricting the domain to $x(U\cap F^{-1}(V))$ as needed. The map is \emph{smooth} when these representatives are smooth. Testing in one pair about each point suffices: changing either chart only pre- or post-composes with a smooth map. Equivalently one may ask directly for such local smooth representatives on open neighborhoods; this implies continuity. Smooth maps compose, and a smooth bijection with smooth inverse is a diffeomorphism. Neither bijectivity nor a merely topological inverse implies this last property.

We use the following Euclidean theorem (the inverse function theorem): a $C^1$ map between open subsets of $\mathbb R^n$ with invertible Jacobian at a point restricts on neighborhoods to a $C^1$ diffeomorphism; if the map is smooth so is the inverse. It is a \emph{local} assertion. The exponential parametrization $t\mapsto(\cos t,\sin t)$ of the circle is locally invertible as a map of one-manifolds, but is not globally injective. The distinction will matter for both coordinates on a hypersurface and coordinates obtained by following geodesics.

\section{Germs, velocities and the tangent space}
A curve $c:(-a,a)\to M$ through $p=c(0)$ differentiates every smooth real function $h$ near $p$: it gives $(h\circ c)'(0)$. This operation cannot distinguish two functions which agree on a neighborhood of $p$, however different they are far away. Their equivalence class under local equality is a \emph{germ} at $p$; the real smooth germs form an algebra under pointwise addition and multiplication on common neighborhoods. Evaluation $h\mapsto h(p)$ is well defined on it.

\begin{definition}
A tangent vector at $p$ is a real-linear map $X$ from the algebra of smooth germs at $p$ to $\mathbb R$ such that $X(hk)=h(p)X(k)+k(p)X(h)$. The space of these maps is $T_pM$. An operation obeying this product rule is a derivation at $p$.
\end{definition}
No coordinates occur in this definition. Why should it give exactly $n$ numbers? In a chart write $x_0=x(p)$. After shrinking to a convex coordinate ball, the fundamental theorem of calculus along line segments gives
\[
 h(x)-h(x_0)=\sum_{j=1}^n(x^j-x_0^j)h_j(x),\qquad
 h_j(x)=\int_0^1\partial_j h(x_0+t(x-x_0))\,dt .
\]
Here $h$ in coordinates denotes $h\circ x^{-1}$; $h_j$ is smooth and $h_j(x_0)=\partial_j h(x_0)$. Since $X(1)=2X(1)$, $X$ kills constants. Applying it to the displayed expression gives
\begin{equation}
 X(h)=\sum_j X(x^j)\,\partial_j h(x_0).\label{f:tangent-coords}
\end{equation}
Conversely this formula with arbitrary coefficients is a derivation. The $n$ coordinate derivatives $\partial/\partial x^j|_p$ form a basis; the coefficients of the same vector in a new chart are $X(y^a)=\sum_j(\partial y^a/\partial x^j)X(x^j)$. In particular these numbers do not specify a new point of $M$. They specify a linear probing operation at one particular point.

The velocity of $c$ is precisely the derivation described at the start of the section. Every derivation is realized by a curve: in a chart take the straight line through $x_0$ with coordinates of velocity $X(x^j)$, then apply $x^{-1}$ for sufficiently small positive and negative parameters. Different curves give the same tangent vector exactly when their first coordinate derivatives agree in one, hence every, chart. Unlike the definition by derivations, this last proof uses the fact that $x_0$ lies in an \emph{open} Euclidean set. On a closed half-line, a two-sided curve through the endpoint cannot have velocity into either half at that instant. This will not cause the dimension of the tangent space to drop.

\section{Differentials, covectors and rank}
A smooth map $F:M\to N$ pushes a tangent vector forward by precomposing what the vector differentiates. For a test germ $h$ at $F(p)$ define
\[
 (dF_p X)(h)=X(h\circ F).
\]
The product rule for $X$ proves that this is a derivation at $F(p)$. In coordinates its matrix has entries $\partial (y^a\circ F)/\partial x^j$ at $p$. The germ formula yields $d(G\circ F)_p=dG_{F(p)}\,dF_p$ by associativity of composition. In particular $dF_p$ is invertible for a diffeomorphism. The converse in equal dimensions holds locally by the inverse function theorem. The \emph{rank} of $F$ at $p$ is the dimension of the image of this linear map: the chain rule proves it cannot be changed by choosing different coordinates.

A covector at $p$ belongs to the dual space $T_p^*M=\operatorname{Hom}(T_pM,\mathbb R)$. For a real function $f$, $df_p(X)=X(f)$. In a chart the $dx^i$ are dual to the coordinate derivatives. The operation on covectors induced by $F$ goes in the other direction:
\[
 F^*\eta=\eta\circ dF_p\ \in T_p^*M,\qquad \eta\in T_{F(p)}^*N .
\]
For example, take $F(u,v)=(u,uv)$ from the plane to a plane with coordinates $x,y$. Then $F^*dy=v\,du+u\,dv$, and $dF(\partial_v)=u\,\partial_y$. At $u=0$ the map annihilates one tangent direction but still pulls back any covector. One cannot in general push a covector forward: it would have to assign a value to every target vector although $dF$ might miss some, and give one answer for all preimages although some vectors may be collapsed. This distinction is why restriction to a submanifold will require no choice of normal direction.

\section{Vector fields, brackets and flows}
The disjoint union of the spaces $T_pM$, with their points of attachment retained, is the tangent bundle $TM$. On a coordinate patch describe it by $(x,v)$ with $v\in\mathbb R^n$; the transition is $(x,v)\mapsto(y(x),D y(x)v)$. It is smooth with smooth inverse. This is our first example of a space built from local vector spaces, to be studied in Chapter~\ref{f:bundles}. A vector field $X$ assigns to each $p$ a vector in $T_pM$ with smooth coordinate coefficients. It differentiates a function to another function: $(Xf)(p)=X_p(f)$. The product rule holds at each point.

For two smooth vector fields define the bracket $[X,Y]f=X(Yf)-Y(Xf)$. The second derivatives of $f$ cancel because they are symmetric, giving
\[
 [X,Y]^j=\sum_i\big(X^i\partial_iY^j-Y^i\partial_iX^j\big).
\]
It is a vector field, independent of coordinates by its original definition. The commutator of operators satisfies the Jacobi identity: expanding the three double commutators $[X,[Y,Z]]$ and its cyclic permutations cancels each ordered product of three operators twice with opposing signs. The bracket is real bilinear, but $[X,fY]=(Xf)Y+f[X,Y]$; in particular it is not a bilinear operation on tangent vectors at just one point.

A local flow $\Phi^X_t(p)$ solves $\partial_t\Phi_t^X(p)=X_{\Phi_t^X(p)}$, $\Phi_0^X(p)=p$. In charts this is a system of ordinary differential equations. Local existence, uniqueness and smooth dependence on the initial point glue the local solutions; uniqueness also gives $\Phi_t^X\circ\Phi_s^X=\Phi_{t+s}^X$ wherever both sides are defined near the trajectory. Each fixed-time map there has inverse $\Phi_{-t}^X$. A vector field is complete if these maps are defined for every real time and starting point. The equation $\dot x=x^2$ on the whole real line has solution $x(t)=x_0/(1-tx_0)$, so a smooth field need not be complete.

The bracket has an observable geometric effect even in the plane. Take $X=\partial_x$, $Y=x\partial_y$. Flow first by $X$ for time $s$, then by $Y$ for time $t$, then undo the first and the second flows in this same order. The four points are $(x+s,y)$, $(x+s,y+(x+s)t)$, $(x,y+(x+s)t)$ and $(x,y+st)$. The displacement is $st\,\partial_y$, agreeing here exactly with $st[X,Y]$. Thus mixed changes of position can measure a derivative of a vector field that inspection at the starting point alone misses.

\Needspace{.58\textheight}
\section*{Exercises}
\begin{exercise}The underlying real line is coordinatized in one atlas by $x$ and in another by $u=x^3$. Show that the identity from the first structure to the second is smooth while its inverse is not. Are the two manifolds nevertheless diffeomorphic?\end{exercise}
\noindent\emph{Solution.} The identity is represented by $x\mapsto x^3$, with inverse not differentiable at zero. A different point map, $x\mapsto\sqrt[3]{x}$ in the underlying real numbers, has coordinate representative $u(\sqrt[3]{x})=x$, a smooth map with smooth inverse. An identification of point sets and an abstract isomorphism of smooth manifolds are different assertions.
\begin{exercise}In the sphere atlas of Example~\ref{f:stereo}, differentiate the transition at $(u,v)=(2,0)$ and find the components of the old vector $\partial_u+\partial_v$ in the new coordinates. Is the transition orientation preserving?\end{exercise}
\noindent\emph{Check.} Its matrix there is $\operatorname{diag}(-1/4,1/4)$. At a general nonzero point inversion reverses the radial direction and preserves the transverse one, both with scale $1/r^2$. Its determinant is $-1/r^4$. A convention for oriented charts will have to account for this sign.
\begin{exercise}The map $i(t)=(\cos t,\sin t)$ takes $\mathbb R$ to the plane. Apply its differential and the pullback of the plane covector field $-y\,dx+x\,dy$ to $\partial_t$. Explain why this does not assign a globally defined real angle to the circle.\end{exercise}
\noindent\emph{Solution.} $di(\partial_t)=(-\sin t,\cos t)$, on which the covector has value one, so the pullback is $dt$. But $i(t+2\pi)=i(t)$ while $t+2\pi\ne t$; the real function $t$ does not descend. Its derivative need not share its obstruction to descent.
\begin{exercise}Prove that every smooth vector field on a compact manifold without boundary is complete. Why does the same assertion fail for a vector field on a closed interval when the endpoints are part of the space?\end{exercise}
\noindent\emph{Hint.} At each point, smooth ODE theory gives a common positive existence time for initial points in a smaller neighborhood. Use a finite cover, take a minimum, and restart the solution. On $[0,1]$ the constant field $+\partial_x$ cannot be followed forward from $1$ inside the interval.
% !TeX root = ../main.tex
\chapter{Submanifolds, boundaries and gluing}\label{f:submanifolds}
A subset may be specified by listing its points, giving a parametrization, or setting functions equal to zero. These descriptions need not display the same information: $y=0$ and $y^2=0$ give the same line, but differentiating the second equation along it gives no constraint at all. We first learn how to recognize a smooth subset from an equation. After that, an inequality will let us keep \emph{one side} of the subset, and a coordinate construction will let us attach another side again.

\section{Parametrizations and regular level sets}
A smooth map $i:N\to M$ is an \emph{immersion} if $di_p$ is injective at every point, a map $F:M\to Q$ is a \emph{submersion} if $dF_p$ is surjective everywhere, and a smooth embedding is an immersion that is a homeomorphism onto its image with the subspace topology. One can see the local meaning of these conditions with only the inverse function theorem.

\begin{theorem}[Local rank theorem]\label{f:rank}
If the rank of $F:M^n\to Q^m$ is constantly $r$ on a neighborhood of an interior point, then suitable coordinates near the point and its image put $F$ in the form $(s,t)\mapsto(s,0)$, with $s$ having $r$ entries.
\end{theorem}
\begin{proof}
Reorder source and target coordinates so the minor $\partial(F^1,\ldots,F^r)/\partial(x^1,\ldots,x^r)$ at the point is invertible. Replace the first $r$ source coordinates by $(F^1,\ldots,F^r)$ and keep the others. The new coordinate map has invertible derivative, hence a smooth local inverse. Now $F(s,t)=(s,G(s,t))$. A column from differentiating in $t$ has zero top block. If its lower block were nonzero, that column would be independent of the first $r$ columns. Thus $\partial_tG=0$ throughout a smaller product neighborhood, and $G(s,t)=g(s)$ there. The target coordinate change $(y,z)\mapsto(y,z-g(y))$ with its evident inverse gives the required form. Translations may be used to center the charts. Maximal rank at a point persists on a neighborhood; rank smaller than the maximum at a single point need not.
\end{proof}
An \emph{embedded submanifold} of dimension $k$ has, at each point, an ambient chart in which the subset is exactly the slice $x^{k+1}=\cdots=x^n=0$. Such charts give it its own smooth structure and make its inclusion an embedding. Conversely the local rank theorem gives slice charts near each \emph{source} point of an immersion; for an embedding the subspace topology lets us shrink the ambient neighborhood so that no other part of the image enters this slice chart. An immersion alone need not have that property. For example, $t\mapsto(\sin t,\sin 2t)$ on the circle has derivative $(\cos t,2\cos 2t)$ never zero, but the parameters $0$ and $\pi$ give a crossing at the origin.

\begin{corollary}[Regular level sets]\label{f:regular-value}
For $F:M^n\to\mathbb R^r$, suppose $dF_p$ is onto at each $p\in F^{-1}(a)$. Then the nonempty level set $N=F^{-1}(a)$ is an embedded submanifold of codimension $r$, and under the inclusion $T_pN=\ker dF_p$.
\end{corollary}
\begin{proof}
Use the rank theorem at each point of the level. In the new coordinates $F$ is projection onto $r$ coordinates (up to target coordinate change). Fixing their values leaves $n-r$ free coordinates; velocities in these and only these coordinates are annihilated by $dF_p$. Restriction of the ambient charts supplies the subspace topology. No condition on distant levels is required.
\end{proof}
For the sphere of radius $R>0$ the differential of $X\mapsto X\cdot X$ at a point of the sphere is $v\mapsto2X\cdot v$, which is nonzero. Thus its tangent vectors satisfy $X\cdot v=0$. At radius zero there is no such regular-value argument: the set is a point but the kernel is the entire ambient space. On the parabola $y=x^2$ the defining function $\rho=y-x^2$ has $d\rho=dy-2x\,dx\ne0$. The coordinates $(x,\rho)$ show both the curve and its two sides. The function $\rho^2$ has the same zero set and neither a useful sign nor a useful differential there.

\section{Directions along and across a subset}\label{f:normal}
The codimension of $N^k\subset M^n$ counts transverse linear directions, not a choice of $n-k$ distinguished arrows. Define the \emph{normal space} as $T_pM/T_pN$ and its dual, the \emph{conormal space}, as
\[
 N^*_pN=\{\xi\in T_p^*M:\xi(v)=0\text{ for every }v\in T_pN\} .
\]
A functional on the quotient is exactly a functional constant on addition of tangential vectors; this proves the identification without a metric. For the regular level of $(F^1,\ldots,F^r)$, the independent $dF^j_p$ span this annihilator by the corollary and dimension count. These spaces vary smoothly in slice charts and give normal and conormal bundles, examples of the construction in Chapter~\ref{f:bundles}. The pullback of each defining differential \emph{to} $N$ is zero, although it is useful as an ambient covector \emph{along} $N$.

Two embedded submanifolds $A,B$ meet transversely when $T_pA+T_pB=T_pM$ at every common point. Near such a point make $B$ the regular level of a submersion $F$; transversality is precisely surjectivity of $d(F|_A)$. The regular-level theorem applied on $A$ makes $A\cap B$ an embedded submanifold of the expected codimension and gives $T_p(A\cap B)=T_pA\cap T_pB$. If one already has a submanifold intersection with this last equality, one calls the intersection \emph{clean}; the sum need not exhaust $T_pM$. A plane in three-space taken twice meets itself cleanly, not transversely. The $x$-axis and the parabola $y=x^2$ meet at a point with equal tangent lines, so they do not even meet cleanly. This is why the shapes of the sets, not just the number of constraints in a list, enter intersection arguments.

\section{What changes at a boundary}
A smooth manifold with boundary uses charts onto relatively open parts of $[0,\infty)\times\mathbb R^{n-1}$ as well as interior charts. More generally a manifold with corners has local models $[0,\infty)^k\times\mathbb R^{n-k}$. A function or coordinate change is smooth up to the walls if locally it is the restriction of a smooth map on an open Euclidean neighborhood; a change must have a smooth inverse in the same sense. Partial derivatives of extensions at a wall are unique: any two extensions agree on the nearby open interior, hence so do their derivatives by continuity up to the wall. The derivation argument of Chapter~\ref{f:smooth} therefore still gives $T_pM$ dimension $n$.

Near a point where exactly $k$ nonnegative coordinates vanish, velocities of curves starting there for $t\geq0$ fill the cone $v^i\geq0$ for those $k$ coordinates. Any such velocity is realized by a ray in a small model chart. The maximal linear subspace inside that cone has dimension $n-k$; the differential of a change of charts, and that of its inverse, map the inward cones into one another. Thus $k$, the \emph{depth}, is intrinsic. The loci of fixed depth have smooth pieces called strata. We use regions with specified embedded faces: each local wall is part of an indicated closed boundary hypersurface, and different incident walls belong to different faces. This excludes identifications between local walls of the same face when counting corners. A stratum of codimension two belongs to two faces in a quadrant; it is not itself a third boundary hypersurface.

A collar of an embedded boundary face $F$ identifies a neighborhood with $[0,\epsilon)\times F$ and sends $(0,p)$ to $p$. For the unit disk one explicit collar is $(r,p)\mapsto (1-r)p$, $p\in S^1$, $0\le r<1/2$; its inverse reads off radius and direction. Ordinary smooth manifolds with boundary admit collars (the collar theorem, see \cite{LeeSM}); the local mechanism is flow from the boundary along a smooth inward transverse vector field. Near one point the differential at time zero maps boundary tangent directions identically and the time direction to the chosen transverse vector. The inverse function theorem in an extended chart then gives a local product. Global control of width, and corners requiring simultaneous products for several walls, are additional matters. In the next construction the compatible product collars are \emph{data}; they are not inferred from the existence of a single face diffeomorphism.

\section{A gluing theorem with checkable hypotheses}\label{f:glue-space}
Suppose $F_-$ in $R_-$ is matched to $F_+$ in $R_+$ by a diffeomorphism $f$. Begin with the quotient of $R_-\sqcup R_+$ identifying only matched \emph{face points}. Let $c_\pm$ be the inward collars and $y$ a coordinate chart on $F_-$. On the two halves of a neighborhood use respectively
\[
 [c_-(r,p)]\longmapsto(-r,y(p)),\qquad
 [c_+(r,f(p))]\longmapsto(r,y(p)).
\]
The formulas agree at zero. For a target pair $(s,z)$ the inverse uses $p=y^{-1}(z)$, choosing the first collar with parameter $-s$ when $s\le0$ and the second with parameter $s$ when $s\ge0$. Each half is a closed subset \emph{of the open coordinate box}; pasting continuous maps on these two closed sets makes the inverse continuous. Continuity of the forward map follows by the quotient test. This argument constructs a chart crossing the face and does not identify any positive-$r$ points on the two collars.

\begin{figure}[htbp]\centering
\begin{tikzpicture}[x=.8cm,y=.66cm,>=Stealth]
\fill[lightpine] (-3,0) rectangle (0,2);\fill[brick!12] (0,0) rectangle (3,2);
\draw[ink] (-3,0) rectangle (3,2);\draw[pine,very thick] (0,0)--(0,2);
\draw[->,ink] (-2.6,1)--(-.5,1);\draw[->,ink] (.5,1)--(2.6,1);
\node at (-1.5,1.5) {$s=-r_-$};\node at (1.5,1.5) {$s=r_+$};\node[below] at (0,0) {$s=0$};
\begin{scope}[shift={(6.4,1)}]\fill[lightpine] (-1,-1) rectangle (0,0);\fill[lightpine] (0,0) rectangle (1,1);\fill[brick!12] (-1,0) rectangle (0,1);\fill[brick!12] (0,-1) rectangle (1,0);\draw[ink] (-1,-1) rectangle (1,1);\draw[pine,thick] (0,-1)--(0,1) (-1,0)--(1,0);\node[below] at (0,-1) {four sectors};\end{scope}
\end{tikzpicture}\caption{Left: a single coordinate increases to the right on both sides. Right: a neighborhood where two internal cuts cross needs four regional sectors. Their interiors consist of different points; only the common subfaces are paired.}\label{f:chart-picture}\end{figure}

To make this local construction global, one must still address the topology and the overlap maps. For a useful finite version assume finitely many regions and pairings of closed faces. At each equivalence class require a saturated neighborhood covered before gluing by one box for each sector of a product $(-\epsilon,\epsilon)^k\times[0,\epsilon)^l\times V$, where $V$ is open Euclidean. The nonnegative regional coordinates of the first $k$ factors receive all possible choices of signs; the next $l$ represent exterior walls never paired. Require the face pairings to be exactly the common subfaces with these coordinates, and the changes on overlapping assembled boxes, including inverses, to extend smoothly in the common coordinates. Distinct proposed signed changes which rescale only one side, for example $s$ on the left and $2s$ on the right, do not pass this test.

\begin{proposition}\label{f:sector-theorem}
Under these assumptions the quotient is a Hausdorff, second-countable manifold with corners. The point corresponding to the center of the displayed product has depth $l$.
\end{proposition}
\begin{proof}
Each box gives a homeomorphism by the same finite closed pasting, and the overlap condition makes the charts smoothly compatible. Every class has a bounded finite number of representatives, at most $b=2^n$ in dimension $n$. To saturate a closed set before taking the quotient, apply all words of length at most $b-1$ in the face pairings and their inverses: remove loops from any longer chain of representatives. A word is a homeomorphism between closed subsets, by restricting at each composition to the preimage of the next closed domain. The resulting saturation of a closed set is closed as a finite union. Two distinct finite classes have disjoint open neighborhoods in the disjoint union, which is Hausdorff. In each such open $U$ keep only the points whose whole class lies in $U$, that is, remove the saturation of its closed complement. The resulting saturated open neighborhoods still contain the respective classes and separate their images downstairs. Finally any open cover of the quotient pulls back to an open cover of a second-countable space and hence has a countable subcover. Apply this to the cover by assembled charts, each of which has a countable basis. The union of the selected local bases is a basis. Only the $l$ exterior coordinates are nonnegative in the assembled product, proving the assertion about depth.
\end{proof}

\section*{Exercises}
\begin{exercise}Intersect the unit sphere in $\mathbb R^3$ with $Z=a$. For which $a$ is the intersection transverse and nonempty? Compare the tangent space of the actual intersection with the intersection of tangent spaces when $a=1$.\end{exercise}
\noindent\emph{Solution.} The normals $(X,Y,Z)$ and $(0,0,1)$ are independent when $|a|<1$. At $a=1$ the intersection is a point; both tangent planes there are horizontal, so their intersection is two-dimensional. For $|a|>1$ the sets do not meet.
\begin{exercise}Give the graph of a smooth map $g:U\subset\mathbb R^p\to\mathbb R^q$ both as an embedding and as a regular level set. Find its conormals without using a metric.\end{exercise}
\noindent\emph{Check.} The embedding is $x\mapsto(x,g(x))$, inverted by projection on its image. The submersion is $(x,y)\mapsto y-g(x)$ because its $y$-block is the identity. The conormals are linear combinations of $dy^a-\sum_i\partial_i g^a\,dx^i$.
\begin{exercise}Join two quadrants along their entire vertical edges, leaving the horizontal edges exterior. Write down a chart at the image of the two origins. Repeat with three half-lines sharing their endpoint and explain why a manifold chart fails there.\end{exercise}
\noindent\emph{Solution.} For the quadrants use $(\pm r,\rho)$ with $\rho\ge0$: the result is a half-plane at a depth-one point. Removing the origin of a small connected neighborhood of the three joined half-lines leaves three components, whereas deleting a point from an interval or half-interval leaves at most two.
\begin{exercise}Let $F=\{\rho=0\}$ be a regular hypersurface. If $\tilde\rho$ is another defining function near a point, prove that $\tilde\rho=a\rho$ there with $a$ smooth and nonzero after shrinking. When do the two functions select the same side?\end{exercise}
\noindent\emph{Hint.} In coordinates $(y,s=\rho)$, subtract $\tilde\rho(y,0)=0$ and integrate $\partial_s\tilde\rho(y,ts)$ for $0\le t\le1$. At $s=0$ its value is nonzero by regularity. The same retained inequality corresponds to $a>0$.
% !TeX root = ../main.tex
\chapter{Differential forms and integration}\label{f:forms}
Euclidean change of variables multiplies a volume element by a determinant. To integrate something not confined to a single chart we must answer two questions: which geometric object transforms by that determinant, and how can the integrals over overlapping charts be assembled without counting anything twice? These questions lead separately to alternating forms and to partitions of unity. They meet in Stokes' theorem. No metric will be needed.

\section{From bilinear to alternating functions}\label{f:tensor-product}
What turns a bilinear rule into an ordinary linear map? For vector spaces $U,V$ over the same field $\mathbb K$, start with formal finite linear combinations of symbols $[u,v]$. Divide by the subspace spanned by the relations making the symbols linear in each slot. For instance, $[a u+b u',v]-a[u,v]-b[u',v]$ is set to zero, as is the corresponding expression in the second slot. The resulting vector space is $U\otimes V$; the class of a symbol is $u\otimes v$. A bilinear map $B:U\times V\to W$ now gives a linear map $U\otimes V\to W$ by $u\otimes v\mapsto B(u,v)$: the relations are killed exactly by bilinearity. Conversely every such linear map gives a bilinear rule, and is determined by that rule since the symbols span. This is the universal property of the tensor product.

If $e_i,f_j$ are bases, their pairwise tensors span by expansion in the two slots. They are independent: the bilinear rules $(u,v)\mapsto e^i(u)f^j(v)$ formed from the dual basis covectors pick out individual coefficients. In particular, for finite-dimensional $U$, the map $U^*\otimes V\to\operatorname{Hom}(U,V)$ given by $\alpha\otimes v\mapsto (x\mapsto\alpha(x)v)$ is an isomorphism: the basis tensors give precisely the elementary matrices. A vector-valued one-form will use this identification in each fiber. Taking $U$ real, the space $U_{\mathbb C}=U\otimes_{\mathbb R}\mathbb C$ allows complex coefficients; its elements are uniquely $u+i v$ with $u,v\in U$. This operation adds no inner product. Iterating the tensor product allows arbitrary multilinear rules; alternating rules are a special case with fewer independent coefficients.

For a finite-dimensional real vector space $V$, a covariant $k$-tensor is a multilinear map $V^k\to\mathbb R$. An alternating one changes sign when two arguments are exchanged; equivalently it vanishes whenever two arguments agree. The space of these maps is denoted $\Lambda^k V^*$, with $\Lambda^0 V^*=\mathbb R$. The exterior product of covectors $\alpha_1,\ldots,\alpha_k$ is the alternating tensor defined by
\[
 (\alpha_1\wedge\cdots\wedge\alpha_k)(v_1,\ldots,v_k)=\det[\alpha_i(v_j)] .
\]
If $e^1,\ldots,e^n$ is a dual basis, evaluating on the increasing $k$-tuples of primal basis vectors shows that the increasing wedges of the $e^i$ form a basis. Define the wedge of elementary wedges by concatenation and extend bilinearly. The determinant formula, or expansion on ordered basis vectors, shows independence of the chosen basis. Moving $k$ slots past $l$ slots changes the sign by $(-1)^{kl}$. In particular there are no nonzero $k$-forms for $k>n$.

A differential $k$-form assigns an alternating tensor to each tangent space. Its coefficients in the coordinate basis must be smooth. In that basis the expansion is unique:
\[
 \eta=\sum_{i_1<\dots<i_k}\eta_{i_1\ldots i_k}\,dx^{i_1}\wedge\cdots\wedge dx^{i_k}.
\]
The pullback by $F:M\to N$ evaluates $\eta$ on $dF_p v_1,\ldots,dF_p v_k$; it commutes with wedges and reverses the order of composition of maps. A top-degree form written $a(y)\,dy^1\wedge\cdots\wedge dy^n$ has pullback coefficient $(a\circ F)\det DF$ in equal dimensions. For a fiber-independent example, take $(r,\theta)\mapsto(r\cos\theta,r\sin\theta)$: $dx\wedge dy$ pulls back to $r\,dr\wedge d\theta$. The formula exists even at $r=0$; it is the use of polar coordinates as a \emph{chart} there which fails.

\section{Exterior differentiation and periods}\label{f:exterior}
Define $d$ on a coordinate expansion by
\[
 d\!\left(\sum_I a_I\,dx^I\right)=\sum_{j,I}\partial_j a_I\,dx^j\wedge dx^I .
\]
This operator restricts on functions to the differential of Chapter~\ref{f:smooth}. Equality of mixed partial derivatives and antisymmetry give $d^2=0$. The ordinary product rule gives the graded rule $d(\eta\wedge\beta)=d\eta\wedge\beta+(-1)^k\eta\wedge d\beta$ for a $k$-form $\eta$. These formulas also prove that $d$ is independent of the chart. Indeed $F^*(dy^i)=d(y^i\circ F)$ by the chain rule, and $d$ of this last one-form vanishes in source coordinates. Applying the graded rule to each pulled-back coefficient times basis wedge gives $dF^*\eta=F^*d\eta$ first for maps between Euclidean open sets and then in manifold charts.

Contraction $\iota_X$ inserts $X$ into the first slot of a form. It is a graded derivation of degree $-1$, as may be checked first on a wedge of covectors by expanding a determinant along its first column. One can also differentiate a form by moving points: set $\mathcal L_X\eta=(d/dt)|_0(\Phi_t^X)^*\eta$. Pullback keeps this a form over each \emph{original} point so the derivative is defined. On functions $\mathcal L_X f=Xf$ and on $df$ it gives $d(Xf)$. Both $\mathcal L_X$ and $d\iota_X+\iota_Xd$ are degree-zero derivations, agree on functions and their differentials by these formulas, and hence agree on every local form. This proves Cartan's formula. It relates infinitesimal symmetry to differentiation and boundary terms without requiring coordinates adapted to the symmetry.

A form is \emph{closed} if $d\eta=0$ and \emph{exact} if $\eta=d\beta$. On the punctured plane consider
\[
 \eta=\frac{-y\,dx+x\,dy}{x^2+y^2} .
\]
Direct differentiation gives zero. Its pullback to $t\mapsto(r\cos t,r\sin t)$ with fixed $r>0$ is $dt$, whose integral along one revolution is $2\pi$. An exact one-form integrates along a path to the endpoint difference of its potential by the one-variable fundamental theorem, and cannot give this answer on a closed path. Locally one may choose a branch of the angle, but a global real angle does not exist. The quotient of closed $k$-forms by exact ones is the $k$th de Rham cohomology; we have exhibited a nonzero class without invoking a classification of surfaces. This period will have an interpretation in parallel transport.

\section{Localizing a smooth integral}
The support of a function or form is the closure of the set where it is nonzero. A family of subsets is locally finite if each point has a neighborhood meeting only finitely many members. A smooth partition of unity subordinate to an open cover is a locally finite family of nonnegative smooth functions whose sum is one, with each support contained in a member of the cover. It is the supports, not just the sets where the functions are positive, that must lie in the desired patches. This lets a local piece be extended by zero without introducing a kink.

\begin{theorem}[Smooth partitions of unity]\label{f:partition}
A Hausdorff second-countable smooth manifold, also with the corner charts of Chapter~\ref{f:submanifolds}, admits a smooth partition of unity subordinate to every open cover.
\end{theorem}
\begin{proof}
In a coordinate patch choose two concentric balls with the closure of the smaller inside the larger and the closure of the larger inside the patch. There is a smooth function equal to one on the smaller ball, supported in the larger, with values in $[0,1]$: a one-variable cutoff can be made from $b(t)=e^{-1/t}$ for $t>0$, $0$ otherwise, using $b(t)/(b(t)+b(1-t))$ on a rescaled squared radius. Use an intermediate radius to put the support strictly inside the larger ball. At a corner restrict the ambient cutoff to the coordinate model. Extend by zero outside the patch.

Here is why enough such choices can be made without losing local finiteness. Coordinate neighborhoods of compact closure form a basis; by second countability select countably many $V_i$ covering the space. There is an exhaustion by compact sets $K_j$ with $K_j\subset\operatorname{int}K_{j+1}$. Inductively cover the compact set $K_j\cup\overline V_{j+1}$ by finitely many relatively compact coordinate neighborhoods; the union of their closures can be taken for $K_{j+1}$. Start with $K_1=\overline V_1$. For each compact shell $K_j\setminus\operatorname{int}K_{j-1}$ choose pairs of balls as above whose larger closures stay in members of the prescribed cover and in $\operatorname{int}K_{j+1}\setminus K_{j-2}$ (early sets with nonpositive indices are empty). Finitely many smaller balls cover the shell. The resulting family is locally finite: inside $\operatorname{int}K_N$ no ball from shells with $j\ge N+2$ can occur. The sum $b$ of the chosen cutoff functions is smooth and positive everywhere since all shells are covered. Dividing each cutoff by $b$ gives the desired partition.
\end{proof}

\phantomsection\label{f:orientation-density}An orientation consists of a continuously consistent choice of sign for ordered tangent bases. Equivalently one can choose charts with positive transition determinants. On a connected oriented manifold there are two global choices. A smooth nowhere-zero form of top degree determines one; conversely local oriented coordinate volume forms can be combined by a partition of unity, since at each point a positive combination of positive volumes stays positive. A \emph{density} of top degree instead changes by the absolute determinant, so can be integrated without orientation. A volume form and the associated positive density are not interchangeable in formulas that keep track of boundary signs.

For a compactly supported top form $\omega$ on an oriented manifold, take an oriented chart cover and a subordinate partition $\chi_i$. Integrate each $\chi_i\omega$ in its chart by the ordinary Euclidean integral of its coefficient, and sum. Only finitely many supports meet the compact support of $\omega$. Given another partition $\psi_j$, insert $1=\sum_j\psi_j$ in each first integral: each $\chi_i\psi_j\omega$ is supported in the intersection. Euclidean change of variables with positive Jacobian equates its two chart integrals. Sum first in $i$ and then in $j$ to get the other answer. This proves that $\int_M\omega$ is an operation on the form, not on its description. The same argument with absolute Jacobians defines integration of densities. Integrable forms of noncompact support require convergence; compact support suffices for the theorem below.

\section{Stokes' theorem and oriented flux}
At a face write the region as $\rho\le0$ with $d\rho\ne0$. A tangent vector to the ambient manifold is outward transverse when $d\rho(N)>0$. The boundary orientation makes $(e_1,\ldots,e_{n-1})$ positive on the face exactly when $(N,e_1,\ldots,e_{n-1})$ is positive in the ambient space. Any other outward transverse vector differs from a positive multiple of $N$ by something tangent, which does not change this test. If two faces meet, exchanging the order in which their outward vectors are inserted reverses the orientation of their common codimension-two corner. Thus $\partial\partial M=0$ as an oriented formal sum, not as a statement that the set of corner points is empty.

\begin{theorem}[Stokes]\label{f:stokes}
For an oriented $n$-manifold $M$ with embedded faces, and a smooth $(n-1)$-form $\eta$ compactly supported up to the boundary,
\[
 \int_M d\eta=\sum_F\int_F j_F^*\eta ,
\]
where each codimension-one face has the outward-first orientation. Use a locally finite face decomposition so the expression on a compact support is finite.
\end{theorem}
\begin{proof}
By the partition construction, $d\eta=\sum_i d(\chi_i\eta)$ (the derivatives of the partition add to zero). It suffices to work with one supported chart form. In an oriented Euclidean orthant chart express it uniquely as
\[
 \eta=\sum_{j=1}^n(-1)^{j-1}a_j\,dx^1\wedge\cdots\wedge\widehat{dx^j}\wedge\cdots\wedge dx^n .
\]
The hat omits a factor. Then $d\eta=(\sum_j\partial_j a_j)\,dx^1\wedge\cdots\wedge dx^n$. Integrating in the $j$ direction over the whole line gives zero by compact support. In a direction restricted to $x^j\ge0$ it gives $-a_j|_{x^j=0}$. On that face all terms of $j_F^*\eta$ containing $dx^j$ vanish; inserting the outward vector $-\partial_j$ first gives exactly the same negative coefficient relative to the remaining coordinate integral. Intersections of faces have measure zero in these integrals. Addition over the pieces proves the formula, with orientation signs transformed in both integrals if the chart reverses the model orientation.
\end{proof}

Given a positive volume form $\epsilon$, normalize $N$ by $\nu(N)=1$ for $\nu=d\rho$ and write $\epsilon_\nu=j^*\iota_N\epsilon$. Then any vector $V$ splits as $\nu(V)N+T$ with $T$ tangent to the face. Alternation kills $j^*\iota_T\epsilon$; hence the outward flux form is $j^*\iota_V\epsilon=\nu(V)\epsilon_\nu$. Rescaling $\nu$ by a positive function rescales $\epsilon_\nu$ inversely. Instead of integrating the form with the boundary orientation one can integrate the outward density $\nu(V)|\epsilon_\nu|$ with a positive measure. On two matched faces the outward densities should then be \emph{added}: no extra minus sign is waiting in a Jacobian. We will use this convention when a variational derivative returns a flux through a boundary.

\section*{Exercises}
\begin{exercise}Orient the unit sphere by its outward normal in $\mathbb R^3$. Pull back its area form in the north-pole chart of Example~\ref{f:stereo}. Use the corresponding density to evaluate the total area.\end{exercise}
\noindent\emph{Check.} For the inverse chart $P(u,v)$ the signed coefficient is $P\cdot(P_u\mathbin{\times}P_v)=-4/(1+r^2)^2$: this planar coordinate order gives the \emph{opposite} orientation. The density is positive. The omitted pole has measure zero in its other chart, and the integral is $2\pi\int_0^\infty4r(1+r^2)^{-2}dr=4\pi$.
\begin{exercise}On the cylinder $S^1\times\mathbb R$ let the height be $z$ and let $d\theta$ denote the pullback of the angular form. Is $dz\wedge d\theta$ exact? Is $d\theta$ exact?\end{exercise}
\noindent\emph{Solution.} $d(z\,d\theta)=dz\wedge d\theta$. A circle at fixed height has angular period $2\pi$; exactness of the two-form does not imply exactness of either factor.
\begin{exercise}In the $(t,x)$ plane oriented by $dt\wedge dx$, retain $x\le h(t)$. Compute the outward flux per unit $|dt|$ for $V=V^t\partial_t+V^x\partial_x$ and check a flow which carries the boundary along with it.\end{exercise}
\noindent\emph{Solution.} Take $\nu=dx-h'(t)dt$ and $N=\partial_x$, so $j^*\iota_N(dt\wedge dx)=-dt$. The outward coefficient is $V^x-h'V^t$, zero if $V^x=h'V^t$. Substituting directly into $\iota_V(dt\wedge dx)=V^t dx-V^x dt$ gives the same oriented form.
\begin{exercise}Construct a smooth vector field on a manifold that agrees with a prescribed smooth vector field on a neighborhood of a point but vanishes outside one chosen coordinate patch containing it. Why would simply setting the coordinate formula to zero on the other patches fail?\end{exercise}
\noindent\emph{Hint.} Multiply the local field by a cutoff identically one on a smaller neighborhood and supported strictly inside the patch. Smoothness near the boundary of its support, not a new rule for transformations of vectors, is the issue.

\part{Geometry of fields and spacetime}
% !TeX root = ../main.tex
\chapter{Vector bundles and symmetry}\label{f:bundles}

A tangent plane to the sphere can be drawn inside $\mathbb R^3$, but two components suffice to specify any vector in it. The difficulty is not the dimension count. It is choosing the two axes as the point moves. A vector bundle lets us use axes on small open sets without assuming that one choice works everywhere. The changes of axes will do more than convert components: we can start with the changes alone and construct the space to which the vectors belong.

\section{Local linear coordinates}

\begin{definition}
A smooth vector bundle of rank $r$ over $M$ consists of a smooth projection $\pi:E\to M$, an $r$-dimensional vector space $E_x=\pi^{-1}(x)$ over each $x$, and local trivializations
\[
 \tau_i:E|_{U_i}:=\pi^{-1}(U_i)\longrightarrow U_i\times\mathbb K^r,
 \qquad \pi=\operatorname{pr}_1\circ\tau_i .
\]
Here $\mathbb K=\mathbb R$ or $\mathbb C$, the $U_i$ cover $M$, and each $\tau_i$ is a diffeomorphism that is linear on fibers in its second coordinate. A smooth section on $U$ is a smooth map $s:U\to E$ with $\pi(s(x))=x$; write $\Gamma(E)$ for the sections on $M$.
\end{definition}

If one replaces $\mathbb K^r$ by a smooth manifold $W$ and drops fiberwise linearity, the same local product condition defines a smooth \emph{fiber bundle}. A section still makes sense, even though two of its values in one fiber cannot in general be added.

For a vector bundle, a trivialization is equally a \emph{frame} $e_i(x):\mathbb K^r\to E_x$, the inverse fiber isomorphism. Its columns are sections, and every section has a unique expression $s=e_i z_i$ there, with $z_i$ a column of smooth functions. Sections can be added and multiplied by functions because the operations are in a single fiber at each point. In particular there is always a zero section. A global frame is a stronger assertion: it gives the isomorphism $(x,v)\mapsto e(x)v$ from the product bundle onto $E$.

A morphism $B:E\to E'$ over $f:M\to N$ is smooth, satisfies $\pi'B=f\pi$, and restricts to a linear map $E_x\to E'_{f(x)}$. For instance $df:TM\to TN$ is such a morphism; it need not be invertible. It is often easier to do the comparison over just one base. Form the \emph{pullback}
\[
 f^*E'=\{(x,w)\in M\times E':\ f(x)=\pi'(w)\},\qquad (x,w)\longmapsto x .
\]
If $e'$ is a frame on $V\subset N$, then $(x,z)\mapsto(x,e'(f(x))z)$ gives a product over $f^{-1}(V)$, even when $f$ is not an immersion. Sending $B(v)$ to $(\pi(v),B(v))$ converts the morphism above into a morphism $E\to f^*E'$ over the identity; projection onto the second factor reverses this construction. Thus we can pull back a section $s'$ to $x\mapsto(x,s'(f(x)))$ without trying to identify tangent spaces at different points. In particular, the velocity of a parametrized curve $c$ is a section of $c^*TM$ over its parameter interval.

\section{Constructing a bundle from changes of frame}

Write $e_j=e_i g_{ij}$ on $U_i\cap U_j$. The same vector then has components $z_i=g_{ij}z_j$: the first index names the destination of a column. For $TM$ in coordinate frames, $g_{ij}$ is the Jacobian of the $i$ coordinates with respect to the $j$ coordinates. Rewriting one frame through two others gives
\begin{equation}
 g_{ii}=1,\qquad g_{ij}g_{jk}=g_{ik}\quad\text{on }U_i\cap U_j\cap U_k .
 \label{f:bundle-cocycle}
\end{equation}
These identities, called the cocycle identities, imply $g_{ji}=g_{ij}^{-1}$.

\begin{theorem}[Bundle gluing]\label{f:bundle-gluing}
Let $M$ be a Hausdorff second-countable smooth manifold. Smooth maps $g_{ij}:U_i\cap U_j\to\mathrm{GL}(r,\mathbb K)$ satisfying \eqref{f:bundle-cocycle} determine a smooth vector bundle. With its stated frames it is unique up to a unique frame-preserving isomorphism. Its sections are exactly the families of smooth columns satisfying $z_i=g_{ij}z_j$.
\end{theorem}
\begin{proof}
In the disjoint union of the products $U_i\times\mathbb K^r$ put $(x,w)_j\sim(x,g_{ij}(x)w)_i$. The cocycle identities give reflexivity, symmetry and transitivity, so no additional identifications arise from chains. Give the set of classes $E$ the quotient topology. Projection to $x$ is continuous. Over $U_i$ the formula
\[
 [x,w]_j\longmapsto (x,g_{ij}(x)w)
\]
is independent of the representative since $g_{ij}g_{jk}=g_{ik}$. Its composite with the quotient map is continuous on the open saturated set over $U_i$, and its inverse is $(x,z)\mapsto[x,z]_i$. Hence these are product homeomorphisms. Points over distinct base points separate downstairs; points over the same point separate in one open product chart. A countable subcover of $M$ gives countably many second-countable products, whose bases together form a basis for $E$. The chart changes are smooth and linear, so the same charts supply the smooth structure and the fiber operations.

For any other realization with frames $\widetilde e_i$, the only possible map preserving them sends $[x,z]_i$ to $\widetilde e_i(x)z$. It is well-defined and a linear isomorphism in every product chart. Finally $x\mapsto[x,z_i(x)]_i$ gives a smooth section precisely when its values agree on overlaps, as claimed.
\end{proof}

The construction also tests whether local maps describe a single morphism. Choose source frames $e_i$ and target frames $e'_i$ on corresponding patches (refine the source cover so their pullbacks exist). Write $B(e_i z)=e'_i(f(x))B_i(x)z$. Computing $B(e_j z)$ in the two ways gives
\begin{equation}
 B_i g_{ij}=(g'_{ij}\circ f)B_j .\label{f:morphism-overlap}
\end{equation}
Conversely this equality lets the local maps descend. In particular two bundles over the same base are isomorphic if invertible $B_i$ obey this equation on a common refinement. A passive change $e_i\mapsto e_i k_i$ in just one bundle rewrites its transitions as $k_i^{-1}g_{ij}k_j$; it has moved no vector.

\begin{exercise}[Dual coordinates]
Determine the transitions on $E^*$ if covectors too are written as columns. Explain why the result is not in general $g_{ij}$, even for a real bundle.

\emph{Solution.} The linear pairing in the $i$ frame is $\lambda_i^{\mathsf T}z_i$. Equality with $\lambda_j^{\mathsf T}z_j$ for all $z_j$ forces $\lambda_i=(g_{ij}^{-1})^{\mathsf T}\lambda_j$. The two transitions agree for real orthogonal matrices, not for an arbitrary change of basis. The complex linear dual also uses transpose, not conjugate transpose.
\end{exercise}

\section{Groups that move the frames}\label{f:group-action}

A Lie group $G$ is a smooth manifold with smooth multiplication and inversion. A smooth left action on $Y$ is a smooth map $(g,y)\mapsto g y$ with $1y=y$ and $(gh)y=g(hy)$. An orbit $G y$ collects the points reachable from $y$; the stabilizer $G_y=\{g:g y=y\}$ collects the moves invisible at that point. Two elements $g,h$ reach the same point precisely when $h^{-1}g\in G_y$. This gives a bijection from the cosets $G/G_y$ onto the orbit; it does not by itself prove the orbit is an embedded submanifold. A right action instead satisfies $(y g)h=y(gh)$, as does composition of frames with matrices on their right.

The infinitesimal motions form the Lie algebra $\mathfrak g=T_1G$. For $A\in T_1G$ take the vector field $A^L(g)=(dL_g)_1A$, where $L_g(h)=gh$. Its bracket with $B^L$ is again left invariant because brackets commute with pushforward by a diffeomorphism. Define $[A,B]=[A^L,B^L](1)$. For a matrix group $A^L(g)=gA$; differentiating $gB$ in direction $gA$ and subtracting the reversed derivative gives $g(AB-BA)$. The curve through $1$ with velocity field $A^L$ is denoted $t\mapsto\exp(tA)$: uniqueness and translation of the local flow give $\exp((t+s)A)=\exp(tA)\exp(sA)$, which extends it to all $t$ \cite{LeeSM}.

A linear representation is a smooth homomorphism $\rho:G\to\mathrm{GL}(W)$ on a finite-dimensional vector space. The derivative $\rho_*=(d\rho)_1$ sends brackets to commutators: differentiating $\rho\circ L_g=L_{\rho(g)}\circ\rho$ makes the left-invariant fields $\rho$-related, and their brackets are then related as well. For example $\mathrm{SO}(2)$ acts on the oriented plane; rotation through $t$ differentiates at $0$ to the operator $(u,v)\mapsto(-v,u)$. For a discrete group the Lie algebra is zero, even though its group elements may act nontrivially. Conjugation by $g$ in $G$ also differentiates at $1$, giving the linear map $\operatorname{Ad}_g:\mathfrak g\to\mathfrak g$; for matrices this is $A\mapsto gAg^{-1}$.

In general a representation of a Lie algebra is a linear map into $\operatorname{End}(W)$ carrying brackets to commutators. Differentiating a group representation gives one, but a Lie algebra representation need not come from the chosen group. For $G=S^1$, its Lie algebra $i\mathbb R$ is abelian. Multiplication by $i\lambda t$ on $\mathbb C$ defines a Lie algebra action of $it$ for every real $\lambda$. A corresponding group action would have to satisfy $\rho(e^{i\theta})=e^{i\lambda\theta}$: differentiating the one-parameter group law gives the scalar ODE with this initial-value solution. It is single-valued on $S^1$ exactly when $\lambda\in\mathbb Z$. The obstruction is a full turn, which no derivative at the identity detects.

\section{Principal fibers and associated fields}

A principal $G$-bundle $P\to M$ carries a smooth right action that is free and transitive on each fiber, with local products in which the action is ordinary right multiplication on $G$. Each fiber resembles $G$ but has no specified identity. A local section $e$ supplies one: every point over $x$ is uniquely $e(x)g$. For example the frames $p:\mathbb K^r\to E_x$ form a principal $\mathrm{GL}(r,\mathbb K)$-bundle, with $p g=p\circ g$. A gauge transformation is an equivariant automorphism $U:P\to P$ over the identity. Writing $U(e_i)=e_i u_i$, equivariance yields $e_j u_j=U(e_i g_{ij})=e_i u_i g_{ij}$, so $u_j=g_{ij}^{-1}u_i g_{ij}$. Conversely such local functions define $U(e_i a)=e_i u_i a$. Unlike a change of notation, $U$ acts on points of $P$.

Given a smooth left action on $W$, form the associated bundle by the rule
\begin{equation}
 P\times_G W=(P\times W)/{\sim},\qquad [p g,w]=[p,g w] .\label{f:associated-bundle}
\end{equation}
In the $e_i$ frame the class of $(e_i a,w)$ has coordinate $a w$, independent of how the class was written. The chart changes are $w_i=g_{ij}w_j$, and are linear when the action is a representation. Taking the frame bundle and its defining representation recovers $E$ by $[p,w]\mapsto p(w)$. A gauge transformation acts by $[p,w]\mapsto[U(p),w]$; it therefore has component matrix $\rho(u_i)$, not its inverse. The conjugation action on $G$ itself gives a bundle of groups whose sections are exactly the families $u_i$ just described.

A choice of extra structure often reduces the group of frames. For a real bundle, pull the Euclidean inner product back by each local frame coordinate map to obtain a positive form $b_i$ on $E|_{U_i}$. Theorem~\ref{f:partition} gives functions $\chi_i$ for which $b=\sum_i\chi_i b_i$ is globally smooth: each term extends by zero off its patch, and the sum is locally finite. For $v\ne0$ at least one positive-weight summand has $b_i(v,v)>0$, so $b$ is a fiber metric. Applying Gram--Schmidt to a smooth local frame now gives a smooth orthonormal frame. Between two of these the transition lies in $\mathrm O(r)$; the orthonormal frames thus make a principal $\mathrm O(r)$-bundle inside the full frame bundle. The construction gives a metric, not a preferred one.

\begin{example}[Frames on the round sphere]
The map $R\mapsto R e_3$ sends $\mathrm{SO}(3)$ onto $S^2$. The stabilizer of $e_3$ is $\mathrm{SO}(2)$ acting in the first two coordinates, so this subgroup acts freely and transitively on the right on each fiber. To see the local products, choose a fixed vector not parallel to $x$ and normalize its projection onto $x^\perp$, obtaining $u(x)$ on a neighborhood. The columns $(u(x),x\times u(x),x)$ make a smooth local section. The associated plane bundle is $TS^2$: the map $[R,(a,b)]\mapsto R(a,b,0)$ lands in $(R e_3)^\perp$ and is a linear isomorphism in the displayed frame. Here the same rotations act on two different spaces: on frames from the right and on plane coordinates from the left.
\end{example}

\begin{exercise}[Sections without a choice of frame]
Show that sections of $P\times_G W$ are equivalent to smooth functions $\varphi:P\to W$ satisfying $\varphi(p g)=g^{-1}\varphi(p)$.

\emph{Hint.} For fixed $p$ there is a unique $w$ with $s(\pi(p))=[p,w]$. Change $p$ to $p g$ in this equation. Conversely $[p,\varphi(p)]$ is independent of the point chosen in the principal fiber; test smoothness using a local section.
\end{exercise}

\section{A circle of unoriented lines}

\begin{example}[The tautological real line]\label{f:circle-line}
The space of unoriented lines through $0$ in $\mathbb R^2$ is a circle. Parametrize it by $\ell(\theta)=\mathbb R v(\theta)$ with $v(\theta)=(\cos(\theta/2),\sin(\theta/2))$ and $\theta$ taken modulo $2\pi$. Over a line put the line itself. On any arc with a lifted angle, $v(\theta)$ is a unit frame. Over the whole circle, $(\theta,z)\mapsto(\ell(\theta),z v(\theta))$ identifies the total space with the quotient by $(\theta,z)\sim(\theta+2\pi,-z)$. A section lifts to a function satisfying $z(\theta+2\pi)=-z(\theta)$. Such a real function must vanish somewhere, by the intermediate value theorem. For example the projection of $(1,0)$ onto each line gives the smooth section with coefficient $\cos(\theta/2)$. The obstruction is to a global frame, not to the existence of sections.
\end{example}

Cover the circle by two connected arcs whose intersection has two components, and on each arc choose a unit frame. Record a sign $a$ for crossing from the first arc to the second on one component and $b$ for returning on the other. Their product $ba$ compares a continuously followed unit vector with its starting value after a circuit. Now fix the patch frames and compare two possible sets of transition data. An isomorphism acts on the two patches by $z_i\mapsto u_i z_i$, with a single constant sign on each connected arc, not independent choices at its two ends. Commutation with the first crossing means $u_2a=a'u_1$, and likewise $u_1b=b'u_2$ at the return. Thus the allowed pair acts on the data by
\begin{equation}
 (a,b)\longmapsto (u_2 a u_1^{-1},\,u_1 b u_2^{-1}).\label{f:sign-groupoid}
\end{equation}
The product is invariant. Taking $u_1=1$, $u_2=a^{-1}$ gives the representative $(1,ba)$, so there are exactly two orbits. The positive product gives a global unit frame; the negative product gives Example~\ref{f:circle-line}. In both cases the stabilizer is the diagonal pair of signs. Keeping objects $(a,b)$ and arrows $(u_1,u_2)$ with their sources and targets gives the \emph{action groupoid}: arrows compose by multiplying the frame maps and every arrow has an inverse. One can retain just the two representatives $(1,+1)$ and $(1,-1)$ and all arrows between them. This smaller groupoid is equivalent to the original: every object is isomorphic to a representative and no arrows between representatives were discarded. It consists of two one-object groupoids, each with automorphism group $\{+1,-1\}$, not just two points with identity arrows. This calculation holds the base fixed and uses the discrete structure group of unit frames.

\begin{exercise}[Unwinding a bundle]
Pull the line of Example~\ref{f:circle-line} back by the degree-$m$ map $e^{i\theta}\mapsto e^{im\theta}$, $m\in\mathbb Z$. For which $m$ is the pullback trivial? What happens to the original line if its real fibers are complexified?

\emph{Solution.} A lifted frame $v(m\theta)$ returns with sign $(-1)^m$. For even $m$ it closes to a global frame; for odd $m$ every real coefficient of a section is antiperiodic and has a zero. On complexification, $e^{i\theta/2}v(\theta)$ is a periodic nowhere-zero complex frame: both factors change sign. The original two-sheeted bundle of real unit frames has not thereby acquired a section.
\end{exercise}
% !TeX root = ../main.tex
\chapter{Connections, curvature and geodesics}\label{f:connections}

In the line bundle of Example~\ref{f:circle-line}, a vector can be followed along a path by keeping its coefficient constant in a local unit frame. This works because two such frames differ by a locally constant sign. For a general frame change $z_+=h z_-$, differentiating produces a term $(\mathrm dh)z_-$. A connection compensates for this term. Applied to the tangent bundle, the same construction will specify what it means for a curve to follow its own velocity without turning.

\section{Differentiation independent of the frame}

A connection on $E$ is a $\mathbb K$-linear derivative $\nabla:\Gamma(E)\to\Gamma(T^*M\otimes E)$ satisfying $\nabla(fs)=\mathrm df\otimes s+f\nabla s$. The output assigns to $X_x\in T_xM$ a vector in $E_x$; evaluation on a vector field is written $\nabla_Xs$. In particular $\nabla_{fX}s=f\nabla_Xs$. The rule is local: if $s$ vanishes near $x$, choose a bump function equal to one near $x$ with support where $s=0$ and differentiate its product with $s$. Thus we can compute in a local frame $e$. Writing the derivatives of the frame columns as $\nabla e_j=\sum_i e_i A^i{}_j$ gives
\[
 \nabla(e z)=e(\mathrm dz+A z).
\]
Every matrix of one-forms $A$ defines such an operator on a product patch. Connections exist globally: take the zero matrices in local frames and sum the resulting local derivatives with a subordinate partition of unity. The derivative of $f$ in the Leibniz rule gets coefficient $\sum_i\chi_i=1$. The difference of two connections has no $\mathrm df$ term and is an $\operatorname{End}(E)$-valued one-form; adding any such form to a connection gives another one.

Suppose $j_-:S\to M_-$ and $j_+:S\to M_+$ are smooth, and $H:j_-^*E_-\to j_+^*E_+$ is a bundle isomorphism over $S$. In pulled-back frames let $z_+=h z_-$ be its action on columns over $S$. A pulled-back connection has local matrix $j_\pm^* A_\pm$, evaluated on $X\in T_yS$ using $(dj_\pm)_y X$. To say that $H$ intertwines the two derivatives is to require
\[
 \mathrm d(h z_-)+(j_+^*A_+)h z_-=h(\mathrm dz_-+(j_-^*A_-)z_-)
\]
for every local column. Cancelling its derivative and multiplying by $h^{-1}$ gives
\begin{equation}
 j_+^*A_+=h(j_-^*A_-)h^{-1}-(\mathrm dh)h^{-1}.\label{f:connection-change}
\end{equation}
For two descriptions of the same section on an open overlap both base maps are inclusions. This proves that the pullback prescription is itself independent of frames, since pullback respects products and commutes with $\mathrm d$. For a face diffeomorphism $f:F_-\to F_+$ with inclusions $\iota_\pm$, the maps are $j_-=\iota_-$ and $j_+=\iota_+\circ f$. Now \eqref{f:connection-change} compares only tangential derivatives. A frame map defined on the face has supplied no transverse derivative of $h$.

\begin{exercise}[Comparing two base coordinates]
On $s>0$ take $j_-(s)=s$, $j_+(s)=s^2=x$, $h=e^{ics}$ and $A_-=i s\,\mathrm ds$ on a complex line, with $c$ real. Find the target potential in its own coordinate.

\emph{Solution.} The pulled-back target potential must be $i(s-c)\,\mathrm ds$. Since $\mathrm dx=2s\,\mathrm ds$, the target potential is $\tfrac i2(1-c/\sqrt{x})\,\mathrm dx$. As a check, $(\mathrm d+A_-)e^{-is^2/2}=0$. On the comparison interval the transformed column is $e^{ics-is^2/2}$; its logarithmic derivative cancels the pulled-back potential.
\end{exercise}

\section{Parallel transport and horizontal frames}

Along a smooth path $\gamma:[0,1]\to M$, a parallel column solves
\begin{equation}
 \dot z(t)=-A_{\gamma(t)}(\dot\gamma(t))z(t).\label{f:parallel-ode}
\end{equation}
A linear ODE with smooth coefficients has a unique solution on each compact parameter subinterval in a chart. Subdivide the path into finitely many such intervals and use the transition matrices at their endpoints. Equation~\eqref{f:connection-change} makes the result independent of the subdivision and frames, giving transport $P_\gamma:E_{\gamma(0)}\to E_{\gamma(1)}$. Reversal gives its inverse and concatenation gives composition by uniqueness. Piecewise smooth paths work by concatenation. If the components change by $h$, the transport matrix becomes $h(\gamma(1))P_\gamma h(\gamma(0))^{-1}$; for a loop these are conjugate matrices of the same endomorphism. How should covectors travel? Demand that $\lambda(t)(s(t))$ stay constant whenever $\lambda$ and $s$ travel in parallel. In dual frames $\dot z=-A(\dot\gamma)z$ then forces $\dot\lambda=A(\dot\gamma)^{\mathsf T}\lambda$. Thus the dual connection has matrix $-A^{\mathsf T}$ and its transport is the dual of $P_\gamma^{-1}$, a map of covectors from the initial point to the final one.

There is a corresponding derivative on frames \cite{KN}. On a right principal $G$-bundle a connection form $\omega$ has values in $\mathfrak g$; it returns $\xi$ on the vertical velocity of $p\exp(t\xi)$ at zero and satisfies $R_a^*\omega=\operatorname{Ad}_{a^{-1}}\omega$. Its kernel is a complement to the fiber tangents, the \emph{horizontal} subspace. In the product given by a section $e$, put $A=e^*\omega$. Differentiating $p=e(x)g$ first in $x$, then in $g$, gives
\[
 \omega=\operatorname{Ad}_{g^{-1}}A+g^{-1}\mathrm dg .
\]
Here $g^{-1}\mathrm dg$ translates a tangent vector to $G$ back to the identity by left multiplication. Conversely this formula constructs $\omega$ from potentials with the principal change-of-section rule $(e k)^*\omega=\operatorname{Ad}_{k^{-1}}A+k^{-1}\mathrm dk$. For the frame bundle take $G=\mathrm{GL}(r,\mathbb K)$: the rule is precisely \eqref{f:connection-change} with $h=k^{-1}$, so every vector bundle connection determines a principal frame connection. For matrix $G$ a horizontal lift satisfies $\dot g=-A(\dot\gamma)g$. Hence in the associated representation $\rho$, transporting $[e(\gamma(0)),w]$ transports its column by $\rho(g(t))$, with $g(0)=1$. Its local derivative has matrix $\rho_* A$. Not every connection on an associated vector bundle must respect a prescribed smaller frame group.

\section{What curvature detects}

\begin{definition}[Curvature]\label{f:curvature}
For vector fields $X,Y$ let $F(X,Y)s=\nabla_X\nabla_Ys-\nabla_Y\nabla_Xs-\nabla_{[X,Y]}s$. The last term subtracts the commutator already present in differentiating functions. Expanding in a frame cancels all derivatives of the column and leaves
\[
 F(X,Y)=(\mathrm dA)(X,Y)+[A(X),A(Y)],\qquad F=\mathrm dA+A\wedge A .
\]
Here $\mathrm dA(X,Y)=X(A(Y))-Y(A(X))-A([X,Y])$. Thus $F$ is an $\operatorname{End}(E)$-valued two-form, and the intertwining law for $\nabla$ gives $F_+=h F_-h^{-1}$ on a common base.
\end{definition}

Transport around all loops based at $p$ gives the \emph{holonomy group} in $\mathrm{GL}(E_p)$. To see curvature in this group, take a coordinate square of side $\epsilon$ based at $p$, traversed first in direction $i$ and then $j$. Near $p$ change components by $h(x)=1+\sum_k A_k(p)(x^k-x^k(p))$. Equation~\eqref{f:connection-change} makes the new $A$ vanish at $p$, so it is $O(\epsilon)$ on the square. Integrating the parallel ODE once around the edges gives $1-\oint A+O(\epsilon^4)$: the first correction to a constant column is $O(\epsilon^2)$ and inserting it in the integral costs a further $O(\epsilon^2)$. By Stokes applied entry by entry, $\oint A=\epsilon^2 F_{ij}(p)+O(\epsilon^3)$, since $A\wedge A$ vanishes at $p$. The frame at $p$ itself was unchanged. Thus curvature, with a minus sign for our parallel equation, is the leading transport per oriented area.

Conversely, if $F=0$, then transport is unchanged by a smooth deformation of paths with fixed endpoints. To see this, pull back the connection to a parameter square $H(t,u)$ and transport the same initial vector in the $t$ direction for each $u$, obtaining $s(t,u)$. The two parameter fields commute, so
\[
 \nabla_t\nabla_u s=\nabla_u\nabla_t s+(H^*F)(\partial_t,\partial_u)s=0 .
\]
Since $\nabla_u s(0,u)=0$, uniqueness of parallel transport gives $\nabla_u s(1,u)=0$ as well. The final vectors are in one fixed fiber, hence equal. In a coordinate ball this constructs a parallel frame by transporting a basis out from its center; path independence makes its columns parallel in every direction. It does not eliminate holonomy along noncontractible loops. On a product complex line over a circle, $A=i\alpha\,\mathrm d\theta$ in a global unit frame gives $F=0$ but transport $e^{-2\pi i\alpha}$ on one positive circuit. Here $\mathrm d\theta$ is global although the angle is not. In contrast the real line of Example~\ref{f:circle-line} has transport $-1$ with $A=0$ in its unit frames: the transition supplies the sign.

\begin{exercise}[The sign of an infinitesimal loop]
On a trivial plane bundle take $A=P\,\mathrm dx+Q\,\mathrm dy$ for constant matrices. Traverse a rectangle with sides $a,b$, starting along $+x$, then $+y$. Find the term of order $ab$ in its transport.

\emph{Hint.} The four factors in their order of action give $e^{bQ}e^{aP}e^{-bQ}e^{-aP}$. Their product through degree two is $1-ab[P,Q]$; the minus sign comes from \eqref{f:parallel-ode}. In particular constant matrix coefficients need not give zero curvature.
\end{exercise}

\section{The connection selected by a metric}

An affine connection is a connection on $TM$. Its torsion is $T(X,Y)=\nabla_XY-\nabla_YX-[X,Y]$, a tensor: the derivatives of functions from the two first-order operations cancel. A pseudo-Riemannian metric $g$ is a smooth symmetric nondegenerate bilinear form on tangent spaces, of fixed signature. Nondegeneracy identifies vectors with covectors by $v^\flat=g(v,\cdot)$; write $\lambda^\sharp$ for the inverse identification, $g^{-1}$ for the resulting pairing on covectors, and $\operatorname{grad} f=(\mathrm df)^\sharp$. Compatibility with $g$ means $Xg(Y,Z)=g(\nabla_XY,Z)+g(Y,\nabla_XZ)$.

\Needspace{8\baselineskip}
\begin{theorem}[Fundamental theorem of pseudo-Riemannian geometry]
There is a unique torsion-free connection compatible with $g$, the Levi--Civita connection.\label{f:levi-civita} It is specified by the Koszul formula
\begin{align*}
 2g(\nabla_XY,Z)={}&Xg(Y,Z)+Yg(Z,X)-Zg(X,Y)\\
 &+g([X,Y],Z)-g([Y,Z],X)+g([Z,X],Y).
\end{align*}
\end{theorem}
\begin{proof}
Write the compatibility equation for $(X,Y,Z)$ and its two cyclic permutations. Add the first two and subtract the third. Replacing $\nabla_YX$ by $\nabla_XY-[X,Y]$, and likewise for the cyclic pairs, leaves the displayed equality. Thus nondegeneracy gives uniqueness. Conversely call its right-hand side $K(X,Y,Z)$. The bracket product rule makes $K$ linear over functions in $X$ and $Z$; for example the terms $Y(f)g(Z,X)-Z(f)g(X,Y)$ in the metric derivatives of $K(fX,Y,Z)$ cancel the bracket terms $-Y(f)g(X,Z)+Z(f)g(X,Y)$. In the middle slot the extra term is $2X(f)g(Y,Z)$, so raising the last slot of $K/2$ defines a connection. Adding $K(X,Y,Z)$ to $K(X,Z,Y)$ gives $2Xg(Y,Z)$, and subtracting $K(Y,X,Z)$ from $K(X,Y,Z)$ gives $2g([X,Y],Z)$. These are compatibility and zero torsion, respectively.
\end{proof}

Write $\nabla_{\partial_i}\partial_j=\Gamma^k{}_{ij}\partial_k$, summing repeated indices. Since coordinate brackets vanish, the Koszul formula gives
\[
 \Gamma^k{}_{ij}=\tfrac12 g^{km}(\partial_i g_{jm}+\partial_j g_{im}-\partial_m g_{ij}).
\]
Our Riemann curvature is the curvature in Definition~\ref{f:curvature} for this connection: $R(\partial_i,\partial_j)\partial_\ell=R^k{}_{\ell i j}\partial_k$, where
\[
 R^k{}_{\ell i j}=\partial_i\Gamma^k{}_{j\ell}-\partial_j\Gamma^k{}_{i\ell}
 +\Gamma^k{}_{im}\Gamma^m{}_{j\ell}-\Gamma^k{}_{jm}\Gamma^m{}_{i\ell} .
\]
For instance on the unit sphere, away from the poles, $g=\mathrm du^2+\sin^2 u\,\mathrm d\phi^2$. Here $\Gamma^u{}_{\phi\phi}=-\sin u\cos u$ and $\Gamma^\phi{}_{u\phi}=\Gamma^\phi{}_{\phi u}=\cot u$ are the only nonzero types. Thus $R^u{}_{\phi u\phi}=\partial_u(-\sin u\cos u)-(-\sin u\cos u)\cot u=\sin^2u$. The two contractions used for a scalar field are $\operatorname{Ric}_{\ell j}=R^i{}_{\ell i j}$ and $R=g^{\ell j}\operatorname{Ric}_{\ell j}$, called the Ricci tensor and the scalar curvature. On this sphere the nonzero diagonal Ricci entries are $1$ and $\sin^2 u$; contracting with the inverse metric gives $R=2$. The curvature convention is thus positive on the round sphere.

\section{One-point and two-point exponential coordinates}

A geodesic is a curve with parallel velocity, $\nabla_{\dot\gamma}\dot\gamma=0$, with its parameter called affine. In coordinates it satisfies $\ddot x^k+\Gamma^k{}_{ij}(x)\dot x^i\dot x^j=0$. Compatibility implies that $g(\dot\gamma,\dot\gamma)$ is constant, including when it is zero. Given $p$ and $v\in T_pM$, the second-order ODE has a unique local solution $\gamma_{p,v}$ with these initial data. Define $\exp_p(v)=\gamma_{p,v}(1)$ where the solution exists up to time one. Rescaling time gives $\exp_p(tv)=\gamma_{p,v}(t)$ near $t=0$; hence $(d\exp_p)_0$ is the identity under $T_0(T_pM)\simeq T_pM$. Smooth dependence and the inverse function theorem give a diffeomorphism from a star-shaped neighborhood of zero onto a \emph{normal neighborhood of $p$}. Straight radial lines in these coordinates are geodesics through $p$. This says nothing yet about a pair of points neither of which is $p$.

\begin{exercise}[What do geodesics forget?]
Suppose two affine connections have the same affinely parametrized geodesics for every initial velocity. What can they still disagree on?

\emph{Solution.} Their difference $C(X,Y)=\widetilde\nabla_XY-\nabla_XY$ is a tensor. On the common geodesic starting with $v$ one has $C(v,v)=0$. Applying this to $v+w$ shows $C(v,w)=-C(w,v)$. Conversely any such alternating $C$ drops out of the geodesic equation but changes the torsion by $2C$. Among torsion-free connections the affinely parametrized geodesics determine the connection.
\end{exercise}

\begin{theorem}[Convex normal neighborhoods]\label{f:convex-normal}
On a smooth pseudo-Riemannian manifold without boundary, each point has arbitrarily small open neighborhoods $U$ normal for every $p\in U$, in which every pair $p,q$ is joined by a unique geodesic segment lying entirely in $U$ (with endpoints at parameters $0,1$). Its initial velocity $v=\exp_p^{-1}(q)$ depends smoothly on $(p,q)\in U\times U$.
\end{theorem}

We use the convex neighborhood theorem for smooth connections here \cite[Ch.~3]{ON}; see also \cite{LeeRM} in the positive-definite case. Smooth dependence of the short branch follows using the endpoint map $(p,v)\mapsto(p,\exp_p v)$. A one-point inverse function theorem alone does not prove convexity. The uniqueness is among segments in $U$, not among ones that leave and reenter it. On a circle, for instance, longer geodesics can wind before reaching the same endpoint.

On such a $U$ let $\gamma$ go from $p$ to $q$ in parameter interval $[0,1]$. The \emph{world function} is the smooth two-point function
\begin{equation}
 \sigma_g(q,p)=\tfrac12 g_p(v,v)=\tfrac12\int_0^1 g(\dot\gamma,\dot\gamma)\,\mathrm dt .\label{f:world-function}
\end{equation}
Vary the endpoints and denote the variation field by $W$. Torsion freeness lets the two parameter derivatives interchange in the derivative of the energy. Metric compatibility then makes its integrand $g(\nabla_tW,\dot\gamma)=\partial_t g(W,\dot\gamma)-g(W,\nabla_t\dot\gamma)$. The integral is just the endpoint term, so
\[
 \operatorname{grad}_q\sigma_g=\dot\gamma(1),\qquad
 \operatorname{grad}_p\sigma_g=-v,\qquad
 g_q^{-1}(\mathrm d_q\sigma_g,\mathrm d_q\sigma_g)=2\sigma_g .
\]
In Lorentzian signature a timelike interval has negative world function. Null-separated points have zero world function, but its endpoint differential need not vanish there.

The mixed derivatives measure the response of initial velocity to moving the final point. In independent coordinates at $q$ and $p$, set
\begin{equation}
 \Delta(q,p)=\frac{\big|\det(-\partial_{q^a}\partial_{p^b}\sigma_g)\big|}
 {\sqrt{|\det g(q)|}\,\sqrt{|\det g(p)|}} .\label{f:van-vleck}
\end{equation}
This is the Van Vleck--Morette determinant on the convex neighborhood; the absolute values make the formula independent of orientation and signature. Indeed $-\partial_{p^b}\sigma_g=g_{bc}(p)v^c$, so the mixed matrix has entries $g_{bc}(p)\,\partial v^c/\partial q^a$. Taking determinants shows that $\Delta=J^{-1}$, where $(\exp_p)^*\mathrm d\mathrm{vol}_g=J\,\mathrm d\mathrm{vol}_{g_p}$: the right-hand density apart from $J$ is constant on $T_pM$. In particular $\Delta(p,p)=1$. In exponential coordinates at $p$, $V=\operatorname{grad}_q\sigma_g$ is the radial field $v^i\partial_{v^i}$; its divergence in the density $J(v)\,\mathrm d\mathrm{vol}_{g_p}$ is $n+V(\log J)$ for $n=\dim M$. Consequently $U_0=\sqrt{\Delta}$ obeys the transport law
\[
 2V(U_0)+(\Box_g\sigma_g-n)U_0=0,\qquad U_0(p,p)=1,
\]
where $\Box_g\sigma_g=\operatorname{div}V$ differentiates in $q$. Conversely any smooth solution with this normalization equals $\sqrt{\Delta}$: its ratio to $\sqrt{\Delta}$ is constant along each radial geodesic and tends to one at its origin. This local formula needs neither completeness nor a choice of geodesic at a distant pair of points.

\begin{exercise}[Spreading from the north pole]
On the unit round sphere, use geodesic polar coordinates about the north pole to compute $\Delta$ for $0<r<\pi$ along the short radial branch. What fails at the antipode?

\emph{Solution.} The tangent-plane density in polar coordinates is $r\,\mathrm dr\,\mathrm d\phi$ while the sphere density is $\sin r\,\mathrm dr\,\mathrm d\phi$. Meridians are unit-speed geodesics, so $J=\sin r/r$ and $\Delta=r/\sin r$ with limit one at zero. At $r=\pi$ all radial directions have the same endpoint; the exponential has singular differential, and this inverse-Jacobian construction no longer applies.
\end{exercise}
% !TeX root = ../main.tex
\chapter{Lorentzian causality}\label{f:causality}

The Levi--Civita connection preserves squared speed. For a positive-definite metric this controls length; for a Lorentzian metric it also distinguishes directions that can carry a signal from directions that cannot. The distinction starts in one tangent space. Whether a signal can return to its starting event, or whether there is a surface that every signal must cross, is a question about whole paths.

\section{Choosing a future cone}

The inertia of a real nondegenerate symmetric form is its number of negative and positive squares in a diagonal basis. These numbers are independent of that basis by Sylvester's law of inertia; for a smooth nondegenerate metric they are locally constant, since the eigenvalues of its symmetric matrix in a local frame vary continuously and none can cross zero. A Lorentzian metric in dimension $n\geq2$ has one negative and $n-1$ positive squares. At a point, choose a basis in which $g(v,v)=-(v^0)^2+|\mathbf v|^2$. A nonzero vector is timelike, null, or spacelike according as this quantity is negative, zero, or positive. A causal vector is timelike or null. Timelike vectors have two components, $v^0>|\mathbf v|$ and $v^0<-|\mathbf v|$. These are convex cones: addition preserves the strict inequalities by the triangle inequality. Each closed cone also includes its null boundary and zero.

A \emph{time orientation} continuously selects one of the timelike components as future; nonzero vectors in its closure are the future causal vectors. Such a choice is equivalent to a smooth timelike vector field $D$. Indeed a field selects the component containing $D_x$ at each $x$. Conversely choose local smooth timelike fields in the selected components and sum them with a partition of unity. At each point the terms with positive weights all lie in the same open convex cone, so the sum is timelike, not zero. In a basis with $D$ along the positive time axis, $g(D,v)<0$ for every nonzero future causal $v$ and has the opposite sign for a past causal $v$. This supplies a test that does not depend on the basis. Parallel transport from Chapter~\ref{f:connections} along a path in a time-oriented manifold preserves the future timelike cone: it preserves the norm, and the sign of the pairing with $D$ cannot change without a zero. By continuity it preserves the future null directions too. A causal path, however, need not have parallel velocity.

\begin{exercise}[A time direction that does not close]
The isometry $(t,x)\mapsto(-t,x+1)$ of the Minkowski plane generates a free, properly discontinuous action. Is its smooth quotient time-orientable? Compare with Example~\ref{f:circle-line}.

\emph{Solution.} A time orientation would pull back to an invariant choice on the connected covering plane. There it must choose $v^0>0$ everywhere or $v^0<0$ everywhere for its timelike directions; the generator exchanges the choices. The subgroup generated by the square acts by $(t,x)\mapsto(t,x+2)$ and preserves them. It is the choice of time cone that twists, just as the choice of real unit frame twisted on the circle.
\end{exercise}

\section{Reachability by directed paths}

A future timelike path may be taken piecewise $C^1$ with future timelike velocities on each regular piece. For causal paths we allow locally Lipschitz curves with nonzero future causal velocity almost everywhere, up to increasing reparametrization; constant waiting intervals can be suppressed. This class does not require a velocity at a corner. Write $p\ll q$ if a future timelike path runs from $p$ to $q$, and $p\leq q$ if $p=q$ or such a causal path exists. Set
\[
 I^+(p)=\{q:p\ll q\},\qquad J^+(p)=\{q:p\leq q\};
\]
these are the chronological and causal futures; reversing the direction gives $I^-(p),J^-(p)$. For a set, take the union of the corresponding futures of its points. Concatenation gives transitivity. A chronological spacetime has no closed timelike path; a causal spacetime has no nonconstant closed causal path, so its relation $\leq$ is a partial order.

In a convex normal neighborhood $U$ of $p$, write $v=\exp_p^{-1}(q)$. Future timelike $v$ gives a future timelike radial geodesic to $q$, and future null $v$ a null one. The local causal comparison theorem (the Lorentzian Gauss lemma and its consequence for curves in $U$, \cite[Ch.~5]{ON}) gives the converses: for $q\ne p$, $q\in I^+(p;U)$ exactly for future timelike $v$, and $q\in J^+(p;U)$ exactly for future causal $v$. Thus inside this neighborhood the sign of $\sigma_g(q,p)$ from \eqref{f:world-function} tests causal character, but not future versus past. The notation $;U$ matters: a path permitted to leave $U$ need not satisfy this endpoint test.

The chronological future is open. On the last smooth part of a timelike path ending at $q$, pick a point $r$ sufficiently near $q$ that this part stays in a convex normal neighborhood of $q$ containing $r$. Then $\exp_r^{-1}(q)$ is future timelike by local comparison. Initial velocities to nearby $q'$ remain in this open cone, giving a new final segment to append to the path ending at $r$. In the Minkowski plane $I^+(p)$ is given by $\Delta t>|\Delta x|$ and $J^+(p)$ by $\Delta t\geq|\Delta x|$: integrate the velocity inequalities for necessity and use a straight segment for sufficiency. Thus the causal future need not be open. Nor must it be closed. Remove $(0,0)$ from this plane and let $p=(-1,-1)$ and $q=(1,1)$, in $(t,x)$ order. Any future causal path in these coordinates satisfies $|\dot x|\leq\dot t$. Reaching $q$ would saturate the integrated bound, forcing $\dot x=\dot t$ almost everywhere and hence $x=t$ throughout the path. The missing origin forbids this. Yet $(1+\epsilon,1)$ for $\epsilon>0$ can be reached by a straight timelike segment avoiding the origin. In this spacetime a limit of points in $I^+(p)$ need not even belong to $J^+(p)$.

\section{Cauchy surfaces and global hyperbolicity}

An inextendible causal curve has neither a future nor a past extension in the spacetime, even after reparametrization. An infinite parameter interval is not the test: a curve approaching an ordinary point in finite coordinate time could be made to take infinite parameter time simply by relabelling it. A smooth spacelike hypersurface is \emph{Cauchy} if every inextendible causal curve meets it exactly once. A \emph{time function} is a continuous function increasing strictly along future causal paths. A \emph{temporal function} $T$ is smooth with everywhere past timelike gradient; $\mathrm dT(v)=g(\operatorname{grad}T,v)>0$ for future causal $v\ne0$, so it is a time function. A temporal function onto $\mathbb R$ is \emph{Cauchy temporal} when all its levels are Cauchy hypersurfaces; equivalently, its values along every inextendible causal curve exhaust $\mathbb R$.

Call a neighborhood causally convex if every causal segment with both endpoints in it stays in it. \emph{Strong causality} means there are arbitrarily small such neighborhoods at every point. Unlike geodesic convexity, this tests all directed paths in the spacetime, not just geodesics inside a chosen open set. A smooth time-oriented spacetime is \emph{globally hyperbolic} if it is strongly causal and every causal diamond $J^+(p)\cap J^-(q)$ is compact. Compactness prevents a sequence of possible intermediate events from escaping the spacetime while both endpoints stay fixed.

\begin{theorem}[Cauchy characterization and smooth splitting]\label{f:cauchy-splitting}
Let $M$ be a connected smooth time-oriented Lorentzian manifold without boundary. Global hyperbolicity, existence of a smooth spacelike Cauchy hypersurface, and existence of a Cauchy temporal function onto $\mathbb R$ are equivalent. In this case one can choose coordinates on a product $M\simeq\mathbb R\times\Sigma$ in which
\[
 g=-\beta(T,y)\,\mathrm dT^2+h_T(y),\qquad \beta>0 ,
\]
with $h_T$ a smooth family of Riemannian metrics and each slice a spacelike Cauchy hypersurface.
\end{theorem}

The equivalence and the existence of a \emph{smooth} Cauchy temporal function are global theorems, not consequences of diagonalizing the metric \cite{BGP,BS}. Once $T$ is available, the reason for the product formula is instructive. Flow along $D=\operatorname{grad}T/g(\operatorname{grad}T,\operatorname{grad}T)$; then $D(T)=1$. A maximal flow line cannot stop at a finite value of $T$: if it had an endpoint, the smooth ODE would extend it, and without an endpoint it would be future- or past-inextendible missing later or earlier Cauchy levels. Hence every line runs through $T=0$ and all other levels. The flow and its inverse (flow back by time $T$) give the product diffeomorphism. Its time vector is normal to the slices, which removes the mixed metric terms; the theorem does not say that $\beta=1$ or that $h_T$ is independent of $T$.

\section{Hypersurfaces: normals, screens and transversals}

Let a smooth hypersurface $F$ be given locally by $r=0$ with $\nu=\mathrm dr\ne0$. Write $k=\nu^\sharp$ for the vector characterized by $g(k,X)=\nu(X)$. Then $T_pF=k^\perp$, and the radical of the restricted metric, the subspace orthogonal to all of $T_pF$ inside $T_pF$, is $T_pF\cap\mathbb R k$. If $g^{-1}(\nu,\nu)<0$, the removed direction is timelike and $g|_{T_pF}$ is positive definite: the hypersurface is spacelike there. A positive conormal norm instead gives a timelike hypersurface, with one negative direction left in $T_pF$. In either case $k$ is transverse and may be normalized. The conclusions are unchanged on multiplying $r$ by a nowhere-zero smooth factor. At a corner this calculation classifies each incident smooth face separately.

\begin{proposition}[The null screen]\label{f:null-screen}
At a null point of $F$, $k$ is tangent, the radical is $\mathbb R k$, and $T_pF/\mathbb R k$ has a canonical positive-definite metric. A screen, meaning a spacelike complement to this line inside $T_pF$, realizes that quotient but need not be unique. A transverse vector lies outside $T_pF$, not in a screen.
\end{proposition}
\begin{proof}
Because $\nu(k)=0$, the radical computation gives the line $\mathbb R k$. Choose $w$ with $g(k,w)=-1$, which is possible by nondegeneracy. Put $l=w+\tfrac12 g(w,w)k$. Then $g(l,l)=0$ and $g(k,l)=-1$, so $l$ is transverse and the plane of $k,l$ has one positive and one negative square. Its orthogonal complement $S$ is positive definite. We have $T_pF=\mathbb R k\oplus S$; changing a representative by a multiple of $k$ does not alter its pairing with another tangent vector. Thus the quotient metric is independent of the choice of $l$.
\end{proof}

On a face that is null everywhere, the radical lines are tangent to null geodesics up to a change of parameter. Indeed $k=\operatorname{grad}r$, so symmetry of the Hessian gives $g(\nabla_k k,Y)=g(\nabla_Y k,k)=\tfrac12 Yg(k,k)$; the Hessian is symmetric because the Levi--Civita connection has zero torsion. For $Y$ tangent to $F$ the last derivative vanishes there. Hence $\nabla_k k$ is a multiple of $k$, and an integrating factor along each integral curve rescales its tangent to a parallel one. This does not turn $k$ into a transverse normal. For example, with $g=-\mathrm du\otimes\mathrm dv-\mathrm dv\otimes\mathrm du+\sum_A(\mathrm dy^A)^2$, on the wall $u=0$ of $u\geq0$ one has $k=-\partial_v$. The field $-\partial_u$ is an outward null transversal, while the $\partial_{y^A}$ span a screen. The normal line $TM|_F/TF$ from Chapter~\ref{f:submanifolds} is a different quotient from the screen quotient $TF/\mathbb R k$: the transversal has nonzero class in the former, whereas $k$ has zero class there. The outward-first orientation uses a transversal even when the metric normal is tangent.

\begin{exercise}[A face of changing type]
In the Minkowski plane classify the graph $t=x^2$ using its conormal. Is there an interval of null generators on it?

\emph{Solution.} The conormal $\mathrm dt-2x\,\mathrm dx$ has norm squared $4x^2-1$. The graph is spacelike for $|x|<1/2$, timelike for $|x|>1/2$, and null just at $x=\pm1/2$. The conclusion about null geodesic generators requires a null hypersurface on an open part of the face, not merely a null tangent at an isolated point.
\end{exercise}

\section{A cylinder with a time shift}

Consider the Minkowski plane with $g=-\mathrm dt^2+\mathrm dx^2$ and future $\partial_t$. Fix $L>0$ and divide by the translations generated by $(t,x)\mapsto(t+\tau,x+L)$. The action is free and properly discontinuous since the spatial translates move by integer multiples of $L$. Its differential is the identity, so both metric and time orientation descend. Write $b=\tau/L$. The pair $(T,[x])=(t-bx,x\bmod L)$ gives global product coordinates on the quotient: to invert, choose a representative $x$ and put $t=T+bx$; another choice changes the lift by a deck transformation. Thus the manifold is a cylinder for every $\tau$.

\begin{theorem}[Causal classification of the shifted cylinder]\label{f:shift-cylinder}
If $|\tau|>L$, every point lies on a closed timelike curve. If $|\tau|=L$, the cylinder is chronological but has closed null curves. If $|\tau|<L$, the function $T=t-(\tau/L)x$ is Cauchy temporal and the cylinder is globally hyperbolic.
\end{theorem}
\begin{proof}
The displacement $V=(\tau,L)$ has squared norm $L^2-\tau^2$. If it is timelike or null, straight segments of displacement $V$ (or $-V$ for future orientation) project to closed causal curves of that type. In the null case a future timelike loop would lift to a path with displacement $mV$ for an integer $m$. The strict inequality $|\dot x|<\dot t$ on timelike pieces gives $|\Delta x|<\Delta t$, impossible both for a null displacement and for $m=0$.

Suppose now $|b|<1$. The covector $\mathrm dT=\mathrm dt-b\,\mathrm dx$ has squared norm $b^2-1<0$ and is positive on future causal velocities. Along a future causal lift,
\[
 \dot T\geq(1-|b|)\dot t,\qquad |\dot x|\leq\dot t .
\]
Integrating between any two parameter values bounds the distance on the spatial circle by the increase of $T$ divided by $1-|b|$. Thus with $T$ as parameter the spatial path is Lipschitz into a complete circle. At a finite supremum of $T$ it would converge, giving an endpoint in the global coordinates $(T,[x])$. One could extend the path by increasing $T$ at fixed $[x]$, a timelike continuation. A finite infimum is excluded by the same past-endpoint argument. Therefore an inextendible path takes every value of $T$, exactly once by strict increase. Theorem~\ref{f:cauchy-splitting} applies to this smooth boundaryless quotient.
\end{proof}

\begin{figure}[htbp]\centering
\begin{tikzpicture}[x=1.13cm,y=1.02cm,>=Stealth,font=\small]
 \fill[lightpine] (0,0)--(-2.45,2.45)--(2.45,2.45)--cycle;
 \draw[pine] (-2.45,2.45)--(0,0)--(2.45,2.45);
 \draw[->,slate] (-2.6,0)--(3.5,0) node[right] {$x$};
 \draw[->,slate] (0,-.23)--(0,2.7) node[above] {$t$};
 \draw[dashed,slate] (1.4,0) node[below] {$L$}--(1.4,2.38);
 \draw[thick,->,brick] (0,0)--(1.4,.5) node[right=3pt] {spacelike};
 \draw[thick,->,ink] (0,0)--(1.4,1.4) node[right=3pt] {null};
 \draw[thick,->,pine] (0,0)--(1.4,2.18) node[above left=1pt] {timelike};
\end{tikzpicture}
\caption{For fixed $L$ the endpoint of the identification vector moves along the dashed line (here $\tau\geq0$). The light cone tests the straight segment. Ruling out other closed timelike paths at the null threshold requires the integrated cone inequality, not just the drawing.}\label{f:shift-cone-figure}
\end{figure}

The endpoint argument depends on which points and curves belong to the space. For instance in the open half-plane $x>0$, the future null ray $(t,-t)$ for $t<0$ stops at the omitted wall. The diamond from $(-2,1)$ to $(2,1)$ contains $(0,1/j)$, $j=1,2,\ldots$, but this sequence has no limit point there. In the closed half-plane the ray has an endpoint and can continue along the timelike wall if such curves are admitted. Results for boundaryless spacetimes do not automatically extend to regions with faces or corners; the permissible curves at the boundary and the meaning of inextendibility are part of a boundary version of the problem.

\begin{exercise}[Arrival after winding]
On the globally hyperbolic cylinder, determine the earliest $s>0$ for a causal signal from $[(0,0)]$ to $[(s,0)]$ that winds at least once. Can a timelike signal attain the same first arrival?

\emph{Solution.} A lift ending after $m\ne0$ windings has displacement $(s+m\tau,mL)$. The necessary inequality $s+m\tau\geq |m|L$ is sufficient by the straight segment. Minimizing $|m|L-m\tau$ over nonzero integers gives $L-|\tau|$. Equality is null; timelike paths require a strict inequality. The comparison measures a global travel time, not a different cone at the point of emission.
\end{exercise}

\begin{exercise}[Which spatial ends are harmless?]
Let $(\Sigma,h)$ be a connected Riemannian manifold without boundary. For the product metric $-\mathrm dt^2+h$ on $\mathbb R\times\Sigma$, show that $t$ is Cauchy temporal if and only if $\Sigma$ is complete for its Riemannian distance (the infimum of path lengths).

\emph{Hint.} Future causal paths satisfy $|\dot y|_h\leq\dot t$. In the complete case repeat the finite-endpoint argument of Theorem~\ref{f:shift-cylinder}. In the incomplete case choose a nonconvergent Cauchy sequence and a subsequence with successive distances less than $2^{-j-1}$. Join successive points by paths of lengths less than $2^{-j}$. Following them with unit spatial speed takes finite total time, but no spatial endpoint can exist. Extend into the past at fixed initial position; this inextendible causal path misses all sufficiently late slices.
\end{exercise}
% !TeX root = ../main.tex
\chapter{Jets and the calculus of variations}\label{f:variation}
The differential of a field is often written $\partial_\mu u^A$. A change of frame can make these numbers nonzero even for a field with constant old components. Yet a Lagrangian is supposed to depend on a field and its first derivatives independently of who described them. Its argument cannot be a list of numbers taken out of their chart. We start by finding an intrinsic first derivative for a section of any bundle; a connection is not required for this step.

\section{The bundle of possible derivatives}\label{f:jet}
Let $\pi:E\to M$ be a smooth fiber bundle and $u$ a section near $p$. At $e=u(p)$ the derivative $\ell=du_p:T_pM\to T_eE$ satisfies $d\pi_e\circ\ell=1_{T_pM}$, by differentiating $\pi\circ u=1_M$. Conversely, in bundle coordinates $(x^\mu,z^A)$ near $e$ every linear map with this property has the form $v\mapsto(v,w^A_\mu v^\mu)$. A section with a coordinate affine-linear fiber component realizes any such $w$ locally. Thus the value $e$ and the splitting $\ell$ specify its \emph{first jet} at $p$. The difference of two splittings with the same value lands in the vertical space $V_eE=\ker d\pi_e$. Consequently first jets over $e$ are an affine space modeled on $\operatorname{Hom}(T_pM,V_eE)$: their differences add linearly, even when no splitting has been distinguished as zero. A connection chooses a horizontal splitting and thereby converts the first jet into an ordinary vertical-valued covariant derivative.

For arbitrary $k$ identify two local sections at the same point if their values and the derivatives of their coordinate functions of total order at most $k$ agree there in a common bundle chart. The class is $j^k_p u$ and all such classes give $J^kE$, with projections to $M$ and to lower jets. Multi-indices record how often each coordinate is differentiated; the arrays have symmetric mixed indices since ordinary partial derivatives commute. The chain rule through order $k$ uses only the old derivatives through this order. Applying also the inverse coordinate change proves equivalence independent of coordinates and gives smooth transitions on these arrays. In a new base coordinate $y$ and new fiber coordinate $Z=Z(x,z)$ the first array is explicitly
\[
 w'^a_i=\left(\frac{\partial Z^a}{\partial x^\mu}+\frac{\partial Z^a}{\partial z^A}w^A_\mu\right)\frac{\partial x^\mu}{\partial y^i} .
\]
The inhomogeneous term is why these numbers are not generally the components of a tensor. At a fixed point every compatible finite symmetric array occurs: take the corresponding polynomial in local coordinates on a sufficiently small neighborhood. A smooth \emph{section} of the jet bundle need not arise by differentiating a single field throughout a neighborhood. Such a section is called holonomic when it does. For example on the line, the proposed value zero everywhere and proposed slope one everywhere give perfectly good individual first jets but not a holonomic section.

\begin{example}Take $y=e^x$, $Z=e^x z$ and a section $z=h(x)$. In the new variables it is $y\,h(\log y)$. At $x=0$ put $(h,h',h'')=(a,b,c)$; at $y=1$ the new array is $(a,a+b,b+c)$. In fact the first derivative is $h+h'$ and the second is $(h'+h'')/y$. Agreement of arrays between \emph{two sections} is invariant, not the old numerical value of an entry. Nor does agreement to every order at one point imply agreement of germs. The smooth extension of $e^{-1/x^2}$ by zero from $x>0$ has the same derivatives there at $x=0$ as the zero function, but differs on every neighborhood of zero.\end{example}

\section{When do local representatives give one field?}
In a vector bundle it is meaningful to add field values in one fiber; in a general bundle it need not be. In either case smoothness can be tested in a common local product and hence by ordinary calculus on the component functions. Across a signed chart $(s,y)$, assume the given sections are smooth up to their closed halves. Perform any prescribed fiber transformation on an actual collar before differentiating. Let $u_-$ and $u_+$ now represent components in the common frame, written using their inward variables $r_-=-s$, $r_+=s$.
\begin{proposition}[Smooth pasting of sections]\label{f:jet-pasting}
The two sections are restrictions of one smooth section across the open face if and only if at every point of the face and for every $j,\alpha$,
\[
 (-1)^j\partial_{r_-}^{j}\partial_y^\alpha u_-(0,y)=\partial_{r_+}^{j}\partial_y^\alpha u_+(0,y).
\]
At a finite product intersection use every mixed derivative in common signed coordinates, for all incident sectors. At a remaining exterior wall, where smoothness means local extendibility, the same test on a choice of extensions across that exterior wall is necessary and sufficient.
\end{proposition}
\begin{proof}
Necessity is restriction of a genuine section or its local extension. Conversely, in signed coordinates each proposed derivative, including order zero, pastes to a continuous function $v_I$ on the box by finite closed pasting. Split a short coordinate segment at every wall it crosses. Fundamental calculus on the individual segments gives
\[
 v_I(x+he_j)-v_I(x)=\int_0^h v_{I+e_j}(x+te_j)\,dt .
\]
A line in a wall uses the common smooth restriction there. The quotient by $h$ tends to $v_{I+e_j}(x)$ on either side. It follows by induction that the pasted value actually has these derivatives. For an exterior wall do the argument on the given compatible extensions and restrict back. The identical necessity and sufficiency statements with total order bounded by $k$ hold for $C^k$ fields. In nonlinear fibers take a chart at the common value and shrink the base box by continuity. The construction is unique pointwise and agrees under change of bundle chart.
\end{proof}

Notice that the theorem starts with \emph{functions on half-neighborhoods}, not formal series assigned to a subset. For a one-form, the transition must change its coordinate covectors as well as its coefficient functions. For a connection its potentials transform affinely as in Chapter~\ref{f:connections}. For a gauge transformation the same argument applies to the group-valued representative in its bundle chart. Gluing fields only up to transformations requires gluing the transformations too. In categorical language, a groupoid is a collection of objects and composable invertible arrows with identities. A functor respects these operations. A restriction functor is an equivalence if every matching pair is isomorphic to a restriction and if, between any two global objects, the map on arrows is bijective (essential surjectivity and full faithfulness). Merely taking a set of orbits discards the stabilizers and does not prove the assertion about arrows.

\section{The first variation and its boundary term}\label{f:first-variation}
A first-order Lagrangian assigns a top density $L(j_p^1u)$ at $p$ to every first jet. In oriented coordinates write it as $\mathcal L(x,u,v)\,d^nx$ with $v^A_\mu=\partial_\mu u^A$ when evaluated on a field. For a family $u_\lambda$ over the fixed base, the variation $Y=\partial_\lambda u_\lambda|_0$ belongs pointwise to the vertical tangent space. The momentum conjugate to a derivative in these coordinates is $p_A^\mu=\partial\mathcal L/\partial v^A_\mu$. Differentiate the action and use the ordinary product rule \emph{before} any field equation:
\begin{align*}
 \partial_\lambda \mathcal L(j^1u_\lambda)|_0
 &=\mathcal L_{u^A}Y^A+p_A^\mu\partial_\mu Y^A\\
 &=\big(\mathcal L_{u^A}-\partial_\mu p_A^\mu(j^1u)\big)Y^A
       +\partial_\mu(p_A^\mu Y^A).
\end{align*}
Repeated indices are summed; the derivative on $p(j^1u)$ is the total derivative along the actual field. The Euler expression, valued in the vertical dual tensored with densities along $u$, is $\mathcal E_A=\mathcal L_{u^A}-\partial_\mu p_A^\mu$. One way to see why it is intrinsic is to compare the first variations under two charts for every compactly supported vertical variation: the integral is intrinsic, the divergences integrate to zero, and an everywhere nonnegative bump supported where a component of the difference has fixed nonzero sign would contradict equality. The same test proves the fundamental lemma used in deriving $\mathcal E=0$ in the interior. Labels before and after a coordinate change need not be the same array.

For the real scalar put $L=-\tfrac12(g^{ab}\partial_a u\,\partial_b u+m^2u^2)dV_g$ with a fixed Lorentzian metric. Metric compatibility and the volume form give
\[
 \delta L(Y)=(\Box_g u-m^2u)Y\,dV_g+d\Theta(Y),
 \qquad \Theta(Y)=-Y\,\iota_{\nabla u}dV_g ,
\]
where for oriented Stokes integrals $dV_g$ denotes the positive volume form and $\Box_g=\nabla_a\nabla^a$. The interior Euler operator has the displayed sign. In a flat time coordinate the contribution of a final time slice is $+\int u_tY\,d^{n-1}x$; on the right spatial end of a strip it is $-\int u_xY\,dt$ with the time measure positive. These signs come from one divergence, not from separate force conventions.

\section{Boundary data, the Hessian and propagation of matched jets}
Consider a slab with oriented boundary $\Sigma_f-\Sigma_i+\Gamma$. A physical wall may carry an action $\int_\Gamma\ell$. With a chosen wall condition, suppose the pullback of $\Theta+\delta\ell$ to the wall on allowed variations is $d_\Gamma C$, where $C$ is a spacetime $(n-2)$-form linear in a field variation. Then Stokes and $\partial\Gamma=-\partial\Sigma_f+\partial\Sigma_i$ give
\[
 \delta S(Y)=\int_M\mathcal E(Y)+\theta_{\Sigma_f}(Y)-\theta_{\Sigma_i}(Y),\qquad
 \theta_\Sigma=\int_\Sigma\Theta-\int_{\partial\Sigma} C .
\]
These identities fix the boundary convention used here \cite{HW}. On a smooth family of fields the symbol $\delta$ means exterior differentiation in the family parameters; $d$ means exterior differentiation in the base, with $\delta d=d\delta$ in these formulas. Thus $\omega_\Sigma=\delta\theta_\Sigma$ is closed. Its kernel and the existence of a quotient on which it becomes nondegenerate will be separate issues in Chapter~\ref{f:symplectic}. For commuting trial variations $u,v$ at a solution, subtract the two second variations. The linearization $J=D\mathcal E$ maps perturbations to dual densities. When both variations implement the stated wall condition one obtains
\begin{equation}
 \langle Jv,u\rangle-\langle Ju,v\rangle
   =\omega_{\Sigma_f}(u,v)-\omega_{\Sigma_i}(u,v) .\label{f:green-variation}
\end{equation}
Here $\langle f,u\rangle=\int f(u)$; on compactly supported interior trials this says $J$ is formally self-adjoint. It is not yet an equality of inverses on any chosen boundary domain. On two homogeneous solutions the slice forms agree, assuming the integrals exist.

An artificial face instead has two contributions. At fixed common trace $q=u|_F$, the boundary variations vanish. When its value is released on both pieces with no face action, the two outward momenta must sum to zero, not vanish individually. Write both in one signed coordinate, with $s\le0$ and $s\ge0$ on the two sides. The remaining term is $(p^s_{-,A}-p^s_{+,A})\delta q^A$ per unit common tangential coordinate measure, so its vanishing for all $\delta q$ gives equality of the \emph{cooriented} momenta. A prescribed face density adds its own Euler variation; the jump then measures its load instead of necessarily vanishing. An independent face field, if present, has an additional equation from its own variation.

\subsection{The normal Legendre map}\label{f:normal-legendre}
Can this finite amount of information enforce the infinite-order matching of Proposition~\ref{f:jet-pasting}? Suppose the two actual half-fields solve the same regular first-order variational problem smoothly up to an ordinary face. At given $s,y,q$ and tangential derivatives of $q$, the map $w^B=\partial_s u^B\mapsto p^s_A$ is the \emph{normal Legendre map}. If both realized velocities belong to one smooth inverse branch with invertible $A_{AB}^{ss}=\partial p^s_A/\partial w^B$, equality of $p^s$ gives equality of $w$. With $A_{AB}^{\mu\nu}=\partial p_A^\mu/\partial v^B_\nu$, the equation now reads
\begin{align*}
 A_{AB}^{ss}\,\partial_s^2u^B={}&\mathcal L_{u^A}-(\partial_{x^\mu}p_A^\mu)_{u,v}
 -\frac{\partial p_A^\mu}{\partial u^B}\partial_\mu u^B\\[-1mm]
 &-\sum_{(\mu,\nu)\ne(s,s)} A_{AB}^{\mu\nu}\partial_\mu\partial_\nu u^B .
\end{align*}
Every term on the right at the face is determined by the common jets of normal order zero or one and their tangential derivatives. Inverting $A^{ss}$ gives a common second normal derivative. After solving for $\partial_s^2u$, differentiate the resulting equation $k$ times in $s$. Its right side needs at most normal order $k+1$, with tangential derivatives; its left side gives order $k+2$. Tangential derivatives are taken along matching functions on the \emph{whole} face. Induction and the pasting result prove smoothness. This uses an invertible coefficient in the equation, not an existence theorem for half-solutions from arbitrary timelike data. The other half of the hypothesis matters: for $\mathcal L=\tfrac12u_t^2-(u_x^2-1)^2/4$ the static piecewise field $u=|x|$ has equal cooriented momenta zero and $A^{xx}=-2$ on each side, but the slopes lie on different branches and do not match.

\section*{Exercises}
\begin{exercise}On the trivial line over the $(x,y)$-plane you are given first-jet data $(q,p_x,p_y)$. If only the two functions $p_x,p_y$ are prescribed on a rectangle, find a necessary condition for a field with these partials and show it is sufficient. What if $q$ is prescribed too?\end{exercise}
\noindent\emph{Solution.} Commuting mixed derivatives gives $\partial_xp_y=\partial_yp_x$. Integrate $p_x$ along a horizontal segment on the bottom and then $p_y$ up at the final $x$. Differentiation using the equality gives both prescribed partials; an additive constant is free. A prescribed $q$ must itself have these partials and the correct base value. On the punctured plane the closed angular form in Chapter~\ref{f:forms} shows why a rectangular domain mattered.
\begin{exercise}Take the potential energy $\frac12\int c(x)(u')^2dx$, where $c$ is constant and positive on each side of zero, and include a spring energy $\tfrac12 kq^2$ at the shared trace $q=u(0)$. Derive the interface equation. Does it imply equality of slopes?\end{exercise}
\noindent\emph{Solution.} With fixed outer endpoints the coefficient of $\delta q$ is $c_-u'_-(0)-c_+u'_+(0)+kq$. Its vanishing is the equation. Even if $k=0$ it is the weighted slopes that agree. This inhomogeneous material is not an artificial subdivision of one smooth homogeneous density.
\begin{exercise}For a scalar with Lagrangian coefficient $\partial_u\phi\,\partial_v\phi$, test the cut $u=0$. Can equal value and cooriented momentum fail to determine the first transverse derivative even on two solutions?\end{exercise}
\noindent\emph{Answer.} $p^u=\partial_v\phi$ follows from the common trace along the face. On $u\le0$ take zero and on $u\ge0$ take $\phi=u$. Both solve $\partial_u\partial_v\phi=0$. The normal Hessian vanishes and the transverse derivatives differ.
\begin{exercise}For a matrix connection on a product collar, $F_{si}=\partial_s A_i-\partial_i A_s+[A_s,A_i]$. Which normal derivative is absent? Can one choose a frame in which matching $F_{si}$ directly matches the normal derivative of $A_i$?\end{exercise}
\noindent\emph{Solution.} No term contains $\partial_sA_s$. With $A'=hAh^{-1}-(dh)h^{-1}$, solve $\partial_s h=hA_s$, $h(0,y)=1$ locally. Then $A'_s=0$ throughout the collar and $F'_{si}=\partial_sA'_i$. This is a frame choice fixed at the face on each side, not the discovery of a new equation for $A_s$. The remaining Yang--Mills constraint still has to hold on the solutions being matched.

\part{Analysis of fields and their boundary values}
% !TeX root = ../main.tex
\chapter{Test functions, Hilbert spaces and distributions}\label{f:distributions}

Consider a nonnegative pulse of total integral one confined to a ball of radius $\varepsilon$ about zero. As the ball shrinks, the pulse has a simple limit in one experiment and no limit at all in another. A smooth detector with profile $\varphi$ records an average that tends to $\varphi(0)$. A detector that measures the squared amplitude of the pulse records a quantity that grows without bound. To speak of the limit of fields we must say what it means for two fields to be close. We shall build two answers: a distance based on mean square, and convergence judged by the readings of smooth tests. Their relation will let us differentiate fields across a surface at which they cease to be smooth.

\section{Completing a space of fields}\label{f:distributions-completion}

A norm on a real or complex vector space assigns a number $\|u\|\geq0$ to each vector, zero only for the zero vector; it satisfies $\|au\|=|a|\|u\|$ and the triangle inequality. It gives a distance $\|u-v\|$. A sequence is Cauchy if its terms become arbitrarily close to one another, whether or not there is a vector to which they converge. A normed space in which every Cauchy sequence converges is a \emph{Banach space}. The completion of a normed space $X$ can be constructed without guessing its missing vectors. Take Cauchy sequences in $X$, identifying $(u_j)$ and $(v_j)$ when $\|u_j-v_j\|\to0$. The norm of a class is $\lim_j\|u_j\|$; the reverse triangle inequality makes this limit well defined. Constant sequences embed $X$, and the terms of any representative approximate its class. To see completeness, approximate the $j$th term of a Cauchy sequence of classes by an element of $X$ to within $2^{-j}$. Those elements form a Cauchy sequence in $X$, and its class is the required limit.

Integration gives important examples of norms. Fix a positive smooth density $d\mu$ on a manifold, or Lebesgue measure $dx$ on $\mathbb R^n$. In the Lebesgue integral, a nonnegative simple function $\sum_i a_i{\bf1}_{A_i}$ on disjoint measurable sets has integral $\sum_i a_i\mu(A_i)$. The integral of a nonnegative measurable function is the supremum of the integrals of simple functions below it. Integrate signed or complex functions by taking their parts when the absolute value is integrable. In particular, changing a function on a null set changes none of these integrals. For $1\leq p<\infty$ put
\[
  \|u\|_p=\left(\int |u|^p\,d\mu\right)^{1/p};
\]
$\|u\|_\infty$ is the smallest almost-everywhere upper bound on $|u|$. The space $L^p$ consists of functions with finite norm, \emph{identified up to equality almost everywhere}. We use H\"older's inequality $\int |uv|\,d\mu\leq\|u\|_p\|v\|_{p'}$ for $1/p+1/p'=1$, and Minkowski's inequality $\|u+v\|_p\leq\|u\|_p+\|v\|_p$ for $p\geq1$ \cite[Appendix~B, integral inequalities]{Evans}. Without the identification the supposed norm would vanish on some nonzero functions.

An inner product $(u,v)$ here is conjugate-linear in its first entry, linear in its second, conjugate symmetric, and positive definite. It induces the norm $\sqrt{(u,u)}$ and satisfies $|(u,v)|\leq\|u\|\,\|v\|$ (minimize $\|u-av\|^2$ over $a$). A complete inner product space is a \emph{Hilbert space}. In $L^2$ we use $(u,v)=\int\overline u v\,d\mu$. Here is why a Cauchy sequence has a limit there. Choose a subsequence $u_{j_k}$ with $\|u_{j_{k+1}}-u_{j_k}\|_2\leq 2^{-k}$. The partial sums of $\sum_k |u_{j_{k+1}}-u_{j_k}|$ have $L^2$ norms bounded by $\sum_k 2^{-k}$. By monotone convergence (the integrals of an increasing sequence of nonnegative measurable functions converge to the integral of its limit \cite[Ch.~I, Thm.~I.10]{RS}), applied to the squares, the full sum belongs to $L^2$. Thus $u_{j_k}$ converges almost everywhere to a function $u$, and the same bound on the tail gives $\|u-u_{j_k}\|_2\leq\sum_{l\geq k}2^{-l}$. The original Cauchy sequence has this limit too. In Euclidean space $C_c^\infty$ is dense in $L^p$ for $1\leq p<\infty$: one cuts off at large distance and smooths by convolution with a compactly supported approximate identity \cite[Appendix~C, mollifiers]{Evans}. Consequently completing smooth tests in mean square really does give $L^2$, not a smaller space of smooth representatives.

The \emph{continuous dual} $X'$ consists of linear functionals $\ell:X\to\mathbb C$ (or $\mathbb R$) satisfying $|\ell(u)|\leq C\|u\|$. The least such $C$ is the dual norm. Continuity at zero implies this bound by scaling a ball about zero, and the bound implies continuity everywhere. More generally, a bounded linear map into a complete normed space extends uniquely from a dense subspace: apply it to any approximating sequence and take the limit. We will use this principle to define a boundary value.

\begin{proposition}[Hilbert space representation]\label{f:hilbert-representation}
For every $\ell$ in the continuous dual of a Hilbert space there is a unique $h$ with $\ell(u)=(h,u)$ and $\|\ell\|=\|h\|$.
\end{proposition}
\begin{proof}
For $\ell\ne0$ take the closed affine set $A=\{v:\ell(v)=1\}$ and let $d=\inf_{v\in A}\|v\|$. A minimizing sequence is Cauchy: its midpoints are in $A$, and the parallelogram identity gives $\|v-w\|^2\leq 2\|v\|^2+2\|w\|^2-4d^2$. Its limit $z\in A$ thus minimizes the norm there. For $k\in\ker\ell$ every $z+ak$ belongs to $A$; minimizing its squared norm at $a=0$ gives $(z,k)=0$ (vary both real and imaginary parts). Since $u-\ell(u)z\in\ker\ell$, we have $(z,u)=\ell(u)\|z\|^2$. Take $h=z/\|z\|^2$. Cauchy--Schwarz gives the norm bound, with equality at $h/\|h\|$. The zero functional has $h=0$, and testing the difference of two representing vectors against itself proves uniqueness.
\end{proof}

\section{Compactly supported tests and their dual}\label{f:test-topology}

Point evaluation is not an $L^2$ functional. Even on smooth functions it is unbounded in the $L^2$ norm: a bump $\varphi(x/\varepsilon)$ equal to one at zero has norm $\varepsilon^{n/2}\|\varphi\|_2$. To keep point evaluation, and also evaluation of derivatives, give the tests a finer topology. For an open set $U\subset\mathbb R^n$ write $\mathcal D(U)=C_c^\infty(U)$. For a compact $K\subset U$ let $\mathcal D_K$ consist of the tests supported in $K$, and put
\[
 p_{K,m}(\varphi)=\max_{|\alpha|\leq m}\ \sup_{x\in K}|\partial^\alpha\varphi(x)|,\qquad m=0,1,2,\ldots .
\]
Here $\alpha$ is a list of nonnegative integers and $|\alpha|$ is their sum. A seminorm satisfies the norm inequalities but may vanish on a nonzero vector. A family of seminorms defines neighborhoods of zero by requiring finitely many seminorms to be smaller than prescribed positive bounds; the resulting topology is called locally convex. On $\mathcal D_K$ convergence means convergence in every $p_{K,m}$. On the union $\mathcal D(U)$ one takes all seminorms whose restriction to each $\mathcal D_K$ is continuous. This is the finest locally convex topology making the inclusions continuous, the \emph{test-function topology}. Its convergent sequences have a common compact support bound and converge there with all derivatives \cite[Ch.~II, \S2.1]{HormI}. A bump translated off to infinity therefore does not converge to zero as a test, even though it vanishes eventually on each fixed compact set.

\begin{definition}\label{f:order-bound}
A distribution on $U$ is a continuous linear functional $T$ on $\mathcal D(U)$. Equivalently, for every compact $K\subset U$ there are $m\geq0$ and $C<\infty$ such that
\begin{equation}
 |T(\varphi)|\leq C p_{K,m}(\varphi)\qquad (\varphi\in\mathcal D_K).\label{f:distribution-estimate}
\end{equation}
The space of distributions is denoted by $\mathcal D'(U)$.
\end{definition}
Indeed a linear functional on the union is continuous exactly when its restrictions to the $\mathcal D_K$ are continuous, by the defining locally convex topology. On one $K$, a neighborhood at zero imposes bounds on finitely many seminorms. They are increasing in $m$, so scaling gives \eqref{f:distribution-estimate} for the largest index needed. Conversely that estimate proves continuity there. A common choice of $m$ for all compact sets bounds the \emph{order} of $T$ on $U$; local bounds need not give one common order. A locally integrable $f$ defines an order-zero distribution $T_f(\varphi)=\int f\varphi\,dx$, since $|T_f(\varphi)|\leq p_{K,0}(\varphi)\int_K |f|\,dx$. The functional $\delta_a(\varphi)=\varphi(a)$ also has order zero, but it cannot equal any $T_f$: integration of $|f|$ over shrinking balls tends to zero, whereas a bounded bump in the ball can have value one at $a$.

The support of $T$ is the complement of the union of all open sets on which it annihilates every supported test. It annihilates tests supported in the union as well: use a partition of unity on the compact support of such a test to write it as a finite sum of tests in the open sets. Thus the support of $\delta_a$ is exactly $\{a\}$. Convergence $T_j\to T$ in the distributional sense means $T_j(\varphi)\to T(\varphi)$ for every test. It makes no assertion about pointwise convergence or convergence in a norm.

On a manifold there is a distinction hidden by $dx$. A distributional \emph{density} acts on compactly supported smooth functions; a distributional \emph{scalar} acts on compactly supported smooth densities. A locally integrable field $u$ with values in a bundle $V$ acts on a test dual density $a\in C_c^\infty(M;V^*\otimes |\Lambda^{\rm top}T^*M|)$ by $\int_M a(u)$. This defines the test space for distributional sections without choosing a volume or a fibre metric. In a chart and frame we apply the compact estimates to the components. A fixed positive density identifies the scalar and density versions, but the integral already knew which object supplied the measure. These distributional pairings are bilinear; the associated Hilbert pairing, when needed, conjugates the test.

\begin{exercise}[Two measurements of one pulse]
Let $\rho\in C_c^\infty(\mathbb R^n)$ be nonnegative with integral one, and set $\rho_\varepsilon(x)=\varepsilon^{-n}\rho((x-a)/\varepsilon)$. Determine its distributional limit, its $L^1$ norm and its $L^2$ norm.
\end{exercise}
\noindent\emph{Solution.} Substitution gives $\int\rho(y)\varphi(a+\varepsilon y)\,dy\to\varphi(a)$ by uniform continuity on the compact integration region. Thus the limit is $\delta_a$, although $\|\rho_\varepsilon\|_1=1$ and $\|\rho_\varepsilon\|_2=\varepsilon^{-n/2}\|\rho\|_2$. No $L^2$ limit is possible.

\section{Differentiating across a surface}\label{f:jump-derivative}

Integration by parts against compactly supported tests suggests the definition $(\partial_j T)(\varphi)=-T(\partial_j\varphi)$. Differentiation of a test raises the index in \eqref{f:distribution-estimate} by at most one, so this defines a distribution. When $T=T_f$ and $\partial_j T=T_h$ with $h$ locally integrable, $h$ is called the \emph{weak derivative} of $f$. Not all distributional derivatives are represented by functions. If $b$ is smooth, $(bT)(\varphi)=T(b\varphi)$ makes sense and the usual product rule with $b$ follows by differentiating $b\varphi$. This construction does not define a product of two arbitrary distributions.

For instance $\delta'_a(\varphi)=-\varphi'(a)$ has support $\{a\}$ and order exactly one. In fact $\partial_x^j\delta_a$ needs order $j$: to rule out a bound of order $m<j$, use $\varepsilon^m\psi((x-a)/\varepsilon)$ with $\psi\in C_c^\infty$ and $\psi^{(j)}(0)\ne0$. Its derivatives through order $m$ stay bounded as $\varepsilon\downarrow0$, but its $j$th derivative at $a$ grows like $\varepsilon^{m-j}$. Thus $\sum_{j\geq1}\partial_x^j\delta_j$ on the line gives a distribution of locally finite but globally unbounded order: only finitely many terms meet any compact set.

Let $r$ be a transverse coordinate and $y$ the coordinates along $r=0$. Write $u$ in the two half-spaces as $u_-$ and $u_+$, smooth up to the common plane, and set $[u]=u_+(0,y)-u_-(0,y)$. Integrating $-u\partial_r\varphi$ on the two halves gives the two interior derivatives and the boundary integral $\int [u](y)\varphi(0,y)\,dy$. Thus
\begin{equation}
 \partial_r u=(\partial_r u)_{\rm pieces}+[u]\delta(r),\qquad
 \partial_r^2 u=(\partial_r^2 u)_{\rm pieces}+[\partial_r u]\delta(r)+[u]\delta'(r).
 \label{f:jump-formula}
\end{equation}
A continuous field produces no sheet in its first derivative, but a jump in its slope is a sheet in its second derivative. A tangential derivative creates no such term: its integration by parts does not cross the interface. In one dimension the first formula gives the derivative of the step function $\Theta$ as $\delta_0$, independently of which value was assigned to $\Theta(0)$.

\begin{exercise}[A singular derivative with no jump law]
On the line show that $\partial_x\log|x|$ is the principal value of $1/x$, that is, the distribution $\varphi\mapsto\lim_{\varepsilon\downarrow0}\int_{|x|>\varepsilon}\varphi(x)/x\,dx$.
\end{exercise}
\noindent\emph{Solution.} The logarithm is locally integrable. Integrating $-\log|x|\,\varphi'(x)$ away from $(-\varepsilon,\varepsilon)$ gives the truncated integral plus the boundary term $\log\varepsilon\,[\varphi(\varepsilon)-\varphi(-\varepsilon)]$. That term tends to zero. The truncated integral has a limit: near zero its paired integrand is $(\varphi(x)-\varphi(-x))/x$, bounded by $2\sup|\varphi'|$. Separate integrals of $\varphi(x)/x$ on the two sides need not exist.

\section{Fourier transformation on tests and on fields}\label{f:fourier-convention}

Fourier transformation does not keep compact support. Instead we use the Schwartz space $\mathcal S(\mathbb R^n)$: its elements are smooth and every $\sup_x|x^\alpha\partial^\beta\varphi(x)|$ is finite. These seminorms specify convergence there. Its continuous linear dual $\mathcal S'$ consists of \emph{tempered distributions}; continuity means a bound by finitely many of the Schwartz seminorms. Cutting a Schwartz test off on successively larger balls approximates it in all those seminorms, by rapid decay and the product rule. Thus a tempered distribution is determined by its compact tests. Every polynomially bounded, locally integrable function (with the growth bound outside a ball) gives one, as does $\delta_a$. Local integrability alone is not enough to pair with all rapidly decreasing tests.

We use the following Fourier inversion and Plancherel theorem as an input \cite[Ch.~VII, \S7.1]{HormI}. For $\varphi\in\mathcal S$ the transform is a continuous bijection of $\mathcal S$ onto itself, with transform and inverse
\begin{equation}
 \widehat\varphi(\xi)=\int_{\mathbb R^n} e^{-ix\cdot\xi}\varphi(x)\,dx,\qquad
 \varphi(x)=(2\pi)^{-n}\int_{\mathbb R^n}e^{ix\cdot\xi}\widehat\varphi(\xi)\,d\xi .\label{f:fourier-pair}
\end{equation}
With this normalization $\varphi\mapsto(2\pi)^{-n/2}\widehat\varphi$ extends to a \emph{unitary} map on $L^2$, meaning a surjection preserving the inner product. In particular,
\begin{equation}
 (u,v)_{L^2}=(2\pi)^{-n}\int\overline{\widehat u(\xi)}\widehat v(\xi)\,d\xi \qquad (u,v\in L^2).\label{f:plancherel}
\end{equation}
The $L^2$ transform is obtained by completion; it need not be an absolutely convergent integral. For $T\in\mathcal S'$ define $\widehat T(\varphi)=T(\widehat\varphi)$. This transpose, with no conjugation or inverse, gives the ordinary transform for integrable fields by interchanging the two integrals. In particular $\widehat{\delta_0}=1$ and $\widehat 1=(2\pi)^n\delta_0$. Differentiating the kernel on a test gives $\widehat{\partial_j T}=i\xi_j\widehat T$. In this representation the question whether a weak derivative is square integrable becomes a question about a weighted integral of $|\widehat u|^2$.

\begin{exercise}[Recovering a fundamental solution]
Find the tempered solution $K$ of $(1-\partial_x^2)K=\delta_0$ on the line and verify its normalization without a transform table.
\end{exercise}
\noindent\emph{Solution.} The transformed equation is $(1+\xi^2)\widehat K=1$. Division is legitimate on Schwartz tests because multiplication by $(1+\xi^2)^{-1}$ preserves them. Thus there is at most one solution. Try $K=c e^{-|x|}$: away from zero $K''=K$, and its slope jumps by $-2c$. Equation~\eqref{f:jump-formula} gives $(1-\partial_x^2)K=2c\delta_0$. Taking $c=1/2$ proves existence as well. Although the forcing is confined to one point, this inverse produces a field on the entire line.
% !TeX root = ../main.tex
\chapter{Sobolev spaces and the trace problem}\label{f:traces}

Suppose a sequence of smooth fields converges in the bulk. Must its restrictions to a dividing surface converge? For mean-square convergence the answer is no: a layer of unit height can become arbitrarily thin without changing its height at the surface. The thinner the layer, however, the larger its normal frequencies. We shall find how much control of these frequencies makes restriction continuous. The answer is different for the value of the field and for its flux through the surface.

\section{Measuring derivatives in mean square}\label{f:sobolev-scale}

Write $\langle\xi\rangle=(1+|\xi|^2)^{1/2}$, using the Euclidean length of the coordinate frequency, even when the coordinates describe a Lorentzian spacetime. A Lorentzian quadratic form changes sign and cannot supply a norm. With the convention \eqref{f:fourier-pair}, the \emph{Sobolev space} $H^s(\mathbb R^n)$, $s\in\mathbb R$, consists of the tempered distributions whose Fourier transforms are functions with finite norm
\begin{equation}
 \|u\|_{H^s}^2=(2\pi)^{-n}\int_{\mathbb R^n}\langle\xi\rangle^{2s}|\widehat u(\xi)|^2\,d\xi .\label{f:sobolev-norm}
\end{equation}
The inner product $(u,v)_{H^s}$ uses $\overline{\widehat u}\,\widehat v$ in place of $|\widehat u|^2$. Thus $H^0=L^2$ by Plancherel, while larger $s$ imposes a stronger condition at high frequencies. Multiplication by $(2\pi)^{-n/2}\langle\xi\rangle^s$ after Fourier transformation takes $H^s$ isometrically \emph{onto} $L^2$: dividing an $L^2$ function by this weight gives a tempered distribution by Cauchy--Schwarz against rapidly decreasing tests. Hence $H^s$ is complete. Schwartz functions are dense since their transforms, multiplied by the weight, still range over all Schwartz functions, a dense subspace of $L^2$.

For an integer $m\geq0$ this definition says exactly that $u$ and all its distributional derivatives up to order $m$ belong to $L^2$. Indeed Plancherel and $\widehat{\partial^\alpha u}=(i\xi)^\alpha\widehat u$ give
\[
 \sum_{|\alpha|\leq m}\|\partial^\alpha u\|_2^2
  =(2\pi)^{-n}\int\!\left(\sum_{|\alpha|\leq m}|\xi^\alpha|^2\right)|\widehat u|^2\,d\xi .
\]
Expansion of $(1+\sum_j\xi_j^2)^m$ gives the same monomials with positive coefficients, so the two norms bound one another up to constants. In particular $H^1$ measures weak first derivatives; it does not require a differentiable representative. For any real $s$, $|\xi_j|\leq\langle\xi\rangle$ shows that $\partial_j:H^s\to H^{s-1}$ is bounded.

The Hilbert inner product on $H^s$ is not the pairing of a load with a field. For a field $q\in H^s$ and a load $F\in H^{-s}$ we use the pairing that equals $(q,F)_{L^2}$ when both entries are square integrable:
\begin{equation}
 \langle F,q\rangle=(2\pi)^{-n}\int\widehat F(\xi)\overline{\widehat q(\xi)}\,d\xi .\label{f:sobolev-duality}
\end{equation}
Insert weights $\langle\xi\rangle^{-s}$ and $\langle\xi\rangle^s$ in the two factors to get $|\langle F,q\rangle|\leq\|F\|_{H^{-s}}\|q\|_{H^s}$. Conversely take any bounded conjugate-linear functional of $q$. Proposition~\ref{f:hilbert-representation}, applied after complex conjugation, writes it as $(q,v)_{H^s}$ for some $v\in H^s$. Taking $\widehat F=\langle\xi\rangle^{2s}\widehat v$ in \eqref{f:sobolev-duality} represents this functional and $\|F\|_{H^{-s}}=\|v\|_{H^s}$. This is the identification of $H^{-s}$ with the \emph{anti-dual} of $H^s$; for real fields the distinction from the linear dual disappears. In the bilinear distribution notation the same functional is $q\mapsto F(\overline q)$ on smooth tests.

\begin{exercise}[A point load]
In what dimensions and for what $s$ does a point source belong to $H^{-s}(\mathbb R^n)$? Compare this with the continuity of point evaluation.
\end{exercise}
\noindent\emph{Solution.} The transform of $\delta_0$ is one. Its squared norm is proportional to $\int_0^\infty (1+r^2)^{-s}r^{n-1}\,dr$, finite exactly for $s>n/2$, with a logarithmic divergence at equality. Duality therefore gives exactly this condition for bounded point evaluation on smooth $H^s$ fields. Above the threshold $\widehat u\in L^1$ by the same weighted Cauchy--Schwarz inequality; Fourier inversion gives a continuous representative.

\section{Restriction to a plane}\label{f:plane-restriction}

Put $x=(y,r)\in\mathbb R^{n-1}\times\mathbb R$, $n\geq2$, with frequency $(k,\lambda)$, and write $a=\langle k\rangle$. For $u\in\mathcal S$, restriction $\gamma u=u(\cdot,0)$ removes one Fourier inversion:
\begin{equation}
  \widehat{\gamma u}(k)=\frac1{2\pi}\int_{\mathbb R}\widehat u(k,\lambda)\,d\lambda .\label{f:trace-fourier}
\end{equation}
Even arbitrarily high normal frequencies contribute at $r=0$ with no oscillatory cancellation. The following estimate controls their sum.

\begin{theorem}[Trace on a hyperplane]\label{f:plane-trace}
For $s>1/2$, restriction extends uniquely to a bounded surjection $\gamma:H^s(\mathbb R^n)\to H^{s-1/2}(\mathbb R^{n-1})$. It has a bounded linear right inverse. For $s\leq1/2$ there is no continuous restriction from $H^s$ even into distributions on the plane agreeing with restriction of Schwartz functions.
\end{theorem}
\begin{proof}
Let $I_s=\int_{\mathbb R}(1+t^2)^{-s}\,dt$, finite for $s>1/2$. Apply Cauchy--Schwarz in \eqref{f:trace-fourier}, with the reciprocal weights in the two factors. Substitution $\lambda=at$ gives
\begin{align*}
  (2\pi)^2 |\widehat{\gamma u}(k)|^2
  &\leq \underbrace{\int_{\mathbb R}(a^2+\lambda^2)^{-s}\,d\lambda}_{a^{1-2s}I_s}
         \int_{\mathbb R}(a^2+\lambda^2)^s |\widehat u(k,\lambda)|^2\,d\lambda .
\end{align*}
Multiply by $a^{2s-1}$ before integrating in $k$. With the two different dimensions in \eqref{f:sobolev-norm} the result is $\|\gamma u\|_{H^{s-1/2}}^2\leq (I_s/2\pi)\|u\|_{H^s}^2$. Density gives the unique extension. For a general $u\in H^s$, the very same inequality makes the normal integral in \eqref{f:trace-fourier} absolutely convergent for almost every $k$ and identifies it with this extended trace.

To recover a prescribed $q\in H^{s-1/2}$, make the two factors in Cauchy--Schwarz proportional as functions of $\lambda$. The normal integral must be $2\pi\widehat q(k)$, which fixes the coefficient:
\begin{equation}
 \widehat{R_s q}(k,\lambda)=\frac{2\pi}{I_s}\,a^{2s-1}(a^2+\lambda^2)^{-s}\widehat q(k).\label{f:trace-lift}
\end{equation}
Integrating in $\lambda$ gives $\gamma R_s q=q$. Integrating its squared modulus with the bulk weight gives $\|R_s q\|_{H^s}^2=(2\pi/I_s)\|q\|_{H^{s-1/2}}^2$, as the same normal integral appears once again. In particular the half-derivative loss cannot be reduced. For example, shift a fixed nonzero Fourier bump $\widehat q$ to frequency $Ne_1$ to obtain $\widehat q_N$. For $\delta>0$, the ratio of the $H^{s-1/2+\delta}$ and $H^{s-1/2}$ norms of $q_N$ grows like $N^\delta$. The fields $R_s q_N$ therefore violate any bound into the stronger space.

At the endpoint we must rule out an estimate, not merely observe that our proof fails. Take a nonzero $h\in C_c^\infty(\mathbb R^{n-1})$ in frequency space. Let $0\leq\chi\leq1$ be smooth, one on $[-1,1]$ and zero outside $[-2,2]$. For $L>2$ set
\[
 b_L(\lambda)=\frac{\chi(\lambda/L)}{\sqrt{1+\lambda^2}},\quad
 B_L=\int b_L(\lambda)\,d\lambda,\qquad
 \widehat u_L(k,\lambda)=2\pi h(k)b_L(\lambda)/B_L .
\]
Then $u_L$ is Schwartz and its trace has transform $h$, independently of $L$. But $B_L$ grows like $\log L$ (integrate $1/\sqrt{1+\lambda^2}$ up to $L$). On the fixed support of $h$, $(a^2+\lambda^2)^{1/2}\leq C(1+\lambda^2)^{1/2}$, so $\|u_L\|_{H^{1/2}}^2\leq C' B_L/B_L^2\to0$, using $\chi^2\leq\chi$. The nonzero fixed traces could not converge even weakly to zero. The same sequence tends to zero in $H^s$ for every smaller $s$.
\end{proof}

The trace is thus a limit of restrictions, not evaluation of an arbitrary representative on a null set. For example, in $n=3$ multiply $|x|^{-1/4}$ by a smooth bump equal to one near zero. The product belongs to $H^1$: its squared gradient is of order $|x|^{-5/2}$, integrable with the radial measure $r^2dr$. Integrating by parts outside a ball of radius $\varepsilon$ leaves a boundary term of order $\varepsilon^{2-1/4}$, so there is no hidden first-derivative sheet. This field has an $H^{1/2}$ plane trace through zero but no continuous representative at the origin. Its value at that single point is not trace data.

\begin{exercise}[What the right inverse does at $s=1$]
Find the tangential Fourier transform $\widehat{R_1q}^{\,y}(k,r)$ at distance $|r|$ from the plane. Is this extension smooth across the plane for every smooth $q$?
\end{exercise}
\noindent\emph{Solution.} Here $I_1=\pi$. Direct integration on $r>0$ and $r<0$ gives $\int e^{-ir\lambda}e^{-a|r|}\,dr=2a/(a^2+\lambda^2)$. Hence \eqref{f:trace-lift} gives $\widehat{R_1q}^{\,y}(k,r)=e^{-a|r|}\widehat q(k)$. Its normal derivative jumps by $-2a\widehat q(k)$; a lift bounded in $H^1$ need not be a smooth continuation. Nor does this formula solve a time evolution problem.

\section{One-sided fields and smooth boundaries}\label{f:localized-trace}

Often the field is initially defined only on one side of the proposed cut. For an open $\Omega\subset\mathbb R^n$ define $H^s(\Omega)$ as the restrictions of whole-space $H^s$ distributions, with norm the infimum of $\|U\|_{H^s(\mathbb R^n)}$ over extensions $U$. Distributional restriction to an \emph{open} set only shrinks the test space: extend a compact interior test by zero. It is not restriction to the boundary. The infimum is a quotient norm, since vanishing on interior tests defines a closed subspace of $H^s$.

Why is the trace on a half-space independent of the whole-space extension? For $s>1/2$ equip every parallel slice with the trace just constructed. Translation is continuous in $H^s$ (on transforms it multiplies by a phase tending to one); therefore $r\mapsto\gamma_r U$ is continuous in $H^{s-1/2}$. Integrating $\gamma_r U$ against a compact smooth test in $r$ recovers the distribution $U$ tested in that variable. This holds first for Schwartz fields and then by the uniform slice estimate. If $U$ vanishes in $r>0$, its slices paired with any tangential test vanish there as distributions of $r$, hence everywhere there by continuity. Letting $r\downarrow0$ gives $\gamma_0 U=0$. Apply this to the difference of two extensions.

To pass to a curved face we use the following localization and extension facts \cite[Ch.~4, \S\S1--4]{Taylor}. For fixed real $s$, a smooth compactly supported multiplier is bounded on $H^s(\mathbb R^n)$; pullback by a smooth coordinate diffeomorphism, localized by cutoffs on relatively compact chart patches, is bounded in both coordinate descriptions. For a half-space, or a bounded domain with smooth boundary, and $s\geq0$, there is a bounded linear extension from $H^s(\Omega)$ into whole-space $H^s$ which restricts to the identity on $\Omega$. At nonnegative integral $s=m$ the restriction norms are equivalent on such domains to the norms formed from all interior weak derivatives through order $m$. In particular the extension applies to a field given originally by its weak derivatives on the domain. One should not extend a general field by \emph{zero}: a smooth nonzero boundary value then produces the jump in \eqref{f:jump-formula}.

Here is the use of those facts. Flatten a smooth boundary in a chart, multiply the field by a cutoff supported within that chart, and extend the resulting half-space field. The flat trace estimate bounds its trace on the smaller patch where the cutoff is one. Independence of the extension and agreement on smooth fields identify the traces on overlaps by density. On a bounded smooth $\Omega$, finitely many such patches define $H^{s-1/2}(\partial\Omega)$ using tangential coordinates and a partition of unity. They give a bounded trace into this boundary space. To lift a given boundary field, split it with the partition of unity, use \eqref{f:trace-lift} on each flattened patch, multiply by a second cutoff equal to one on the boundary support of that piece, pull back and add. This gives a bounded right inverse on the smooth boundary. The same argument is local in a smooth manifold or a bundle frame. It does not assign independent values on meeting faces of a corner. Applied to a first derivative it gives an ordinary derivative trace in $H^{s-3/2}$ for $s>3/2$; an $H^1$ bound alone does not give that trace.

\section{Flux without a pointwise normal component}\label{f:normal-flux}

Let $\Omega$ be a bounded smooth Euclidean domain. For a smooth vector field $p$ the integration by parts identity for divergence involves $p\cdot\nu$, where $\nu$ is the outward unit normal. For an $L^2$ field this expression has no intrinsic boundary value. The identity itself admits an extension on the graph domain
\[
 H(\operatorname{div};\Omega)=\{p\in L^2(\Omega;\mathbb C^n):\operatorname{div}p\in L^2(\Omega)\},
\]
where the divergence is distributional and the squared graph norm is $\|p\|_2^2+\|\operatorname{div}p\|_2^2$. For the following argument we need the zero-trace characterization: on a bounded smooth domain, $\ker(\gamma:H^1(\Omega)\to H^{1/2}(\partial\Omega))$ is $H^1_0(\Omega)$, the closure in $H^1$ of compactly supported interior tests \cite[\S5.5, trace-zero theorem]{Evans}.

\begin{proposition}[Weak outward flux]\label{f:weak-flux}
There is a bounded normal-trace map $H(\operatorname{div};\Omega)\to H^{-1/2}(\partial\Omega)$ agreeing with $p\cdot\nu$ on smooth fields, with surface measure understood in the pairing.
\end{proposition}
\begin{proof}
For $q\in H^{1/2}(\partial\Omega)$ choose $v\in H^1(\Omega)$ with trace $q$. Define
\begin{equation}
 \langle\gamma_\nu p,q\rangle=\int_\Omega\big((\operatorname{div}p)\overline v+p\cdot\nabla\overline v\big)\,dx .\label{f:flux-pairing}
\end{equation}
Two choices differ by an $H^1_0$ field. The expression on their difference vanishes for compact interior tests by definition of divergence, hence for their $H^1$ closure by Cauchy--Schwarz. Taking the bounded lift from the preceding section bounds the absolute value by $C\|p\|_{H(\operatorname{div})}\|q\|_{H^{1/2}}$. It therefore defines an element of the boundary anti-dual, denoted $H^{-1/2}$ using \eqref{f:sobolev-duality} in charts and the surface density. For smooth $p,v$ the divergence theorem gives exactly $\int_{\partial\Omega}(p\cdot\nu)\overline q\,dS$.
\end{proof}

In a coordinate patch one can use the components of a vector density and its coordinate divergence in the same integral. The normal flux is then a boundary density; no Lorentzian unit normal is needed to define it. For example $p=A\nabla u$ belongs to $L^2$ when $A$ is bounded and $u\in H^1$. The additional requirement $\operatorname{div}(A\nabla u)\in L^2$ is an equation-domain condition, not part of the definition of $H^1$. This is how a differential equation can supply a flux which the field's Sobolev index alone did not provide.

\begin{exercise}[Testing the joined flux]
A smooth interface separates two domains. Let $p_\pm$ be in the divergence graph domain on their sides and define $p$ by joining them off the interface. On tests supported away from the exterior boundary, compute $\operatorname{div}p$ and identify the matching condition for an $L^2$ divergence across the cut.
\end{exercise}
\noindent\emph{Solution.} Apply \eqref{f:flux-pairing} on each side with the same test $v$. The action of $\operatorname{div}p$ on $\overline v$ is the sum of the two bulk pairings with $\operatorname{div}p_\pm$, \emph{minus} $\langle\gamma_{\nu_-}p_-+\gamma_{\nu_+}p_+,\gamma v\rangle$. The interface term vanishes when the outward fluxes sum to zero. Conversely if the global divergence belongs to $L^2$, it equals the two bulk divergences almost everywhere, since the interface has measure zero. Their difference on tests must therefore vanish. Every smooth compact face test can be lifted in a collar, so the sum of the outward fluxes is zero there as a distribution.
% !TeX root = ../main.tex
\chapter{Hyperbolic operators and Green operators}\label{f:waves}

A trace theorem tells us which boundary values an arbitrary field can have. It does not tell us which boundary values a wave can have if there was no wave yesterday. For the massless wave of speed one on $x>0$, an outgoing disturbance has the form $q(t-x)$: its value at the end is $q(t)$ and its spatial derivative there is $-q'(t)$. These are two measurements of one history. To understand this distinction we first construct the evolution of a forced wave, then use that evolution to determine what can be prescribed at a cut.

\section{The principal symbol and the choice of initial surface}\label{f:wave-symbol}

A linear differential operator of order $j$ has the local form $P=\sum_{|\alpha|\leq j}a_\alpha(x)\partial^\alpha$, with smooth matrix coefficients for a system. In the Fourier convention \eqref{f:fourier-pair}, its \emph{principal symbol} is
\[
 \sigma_j(P)(x,\xi)=\sum_{|\alpha|=j}a_\alpha(x)(i\xi)^\alpha .
\]
One can see why this is intrinsically a function of a covector: the coefficient of $L^j$ in $e^{-iL\phi}P(e^{iL\phi}u)$ at $x$ is $\sigma_j(P)(x,d\phi_x)u(x)$. A nonzero covector at which the symbol is not invertible is \emph{characteristic} (for a system, take equal input and output ranks). A hypersurface $r=0$, $dr\ne0$, is noncharacteristic if its conormals are noncharacteristic. In adapted coordinates this says that one can solve the equation algebraically for the highest pure normal derivative. It does not say that solving for it gives a stable evolution.

Fix signature $(-,+,\ldots,+)$, and define the scalar operator, for a constant $m\geq0$, by
\begin{equation}
 P=\Box_g-m^2,\qquad \Box_g u=\nabla^\mu\nabla_\mu u
   =|\det g|^{-1/2}\partial_\mu\big(|\det g|^{1/2}g^{\mu\nu}\partial_\nu u\big).
 \label{f:wave-operator}
\end{equation}
This is the Euler operator of the action $-\tfrac12\int(g^{\mu\nu}\partial_\mu u\partial_\nu u+m^2u^2)\,dV_g$ after removing the volume density. Its principal symbol is $-g^{\mu\nu}\xi_\mu\xi_\nu$. A second-order system with this symbol times the identity is called \emph{normally hyperbolic} in our sign convention. Its characteristics are the null covectors. A spacelike slice is noncharacteristic, but so is a timelike wall. In Minkowski space there are inextendible timelike curves which never meet the wall. A Cauchy surface, by contrast, meets every inextendible timelike curve exactly once; the wave theorem below requires a smooth spacelike Cauchy surface. On a null face one loses even the algebraic determination of the second normal derivative. For example $4\partial_u\partial_v w=0$ has solutions $a(u)+b(v)$: data at $u=0$ do not determine $a(u)-a(0)$, and adding the tangential derivative $\partial_vw(0,v)$ gives no new information.

\begin{exercise}[Evolution in a spatial direction]
For the equation $w_{tt}-w_{xx}-w_{yy}=0$ with periodic $y$, show that continuation of two traces away from the timelike plane $x=0$ can amplify a small error in the data by an arbitrarily large factor.
\end{exercise}
\noindent\emph{Solution.} Fix $a>0$ and take $w_N=e^{-Na/2}\cos(Ny)\cosh(Nx)$. It is a time-independent solution. At $x=0$ the value tends to zero with all its tangential derivatives and the normal derivative is identically zero. At $x=a$ the amplitude grows like $e^{Na/2}/2$. Noncharacteristicity did let us compute $w_{xx}$; it gave no estimate in the proposed direction of evolution.

\section{Energy, finite speed and existence in flat space}\label{f:wave-energy}

For the estimates it is convenient to work with the \emph{positive evolution operator} $D=\partial_t^2-\Delta_x+m^2=-P$. Thus $Du=h$ means $Pu=-h$; every inverse of $P$ below is the \emph{negative} of the corresponding inverse of $D$. For a smooth real solution in $d$ spatial dimensions multiply $Du=h$ by $u_t$. The product rule gives a local identity, not just a conserved total:
\begin{equation}
 e=\tfrac12(u_t^2+|\nabla_x u|^2+m^2u^2),\qquad
 \partial_t e-\operatorname{div}_x(u_t\nabla_xu)=h u_t .\label{f:wave-balance}
\end{equation}
With sufficient spatial decay to discard the boundary at infinity, put $\mathcal E(t)=\int e(t,x)\,dx$. Then $\mathcal E'\leq\|h\|_2\sqrt{2\mathcal E}$. Integrating the inequality for $\sqrt{2\mathcal E+\varepsilon}$ and letting $\varepsilon\downarrow0$ yields, for $t\geq t_0$,
\begin{equation}
 \sqrt{2\mathcal E(t)}\leq \sqrt{2\mathcal E(t_0)}+\int_{t_0}^t\|h(\tau,\cdot)\|_2\,d\tau .\label{f:energy-estimate}
\end{equation}
Apply this also to differences of solutions. It gives uniqueness with fixed Cauchy data and continuous dependence in the energy norm. If $m=0$, energy controls derivatives; the initial value fixes the remaining constant, and $\|u(t)\|_2\leq\|u(t_0)\|_2+\int_{t_0}^t\|u_t\|_2\,d\tau$ supplies the missing norm on finite time intervals.

For a local conclusion, integrate \eqref{f:wave-balance} over the shrinking ball $|x-x_*|<T-t$. Differentiating this integral in time contributes $-e$ on the moving sphere, so its boundary integrand is
\[
 u_t\partial_\nu u-e=-\tfrac12(u_t-\partial_\nu u)^2-\tfrac12|\nabla_{\rm tan}u|^2-\tfrac12m^2u^2\leq0 .
\]
If the source is absent in this backward cone and the data vanish on its initial ball, the energy in each shrinking ball remains zero. The solution is zero there, also in the massless case by the initial value, and continuity gives $u(T,x_*)=0$. This is finite speed one. On a static internal cut the two boundary power terms are $u_t$ times the two outward spatial derivatives. If the field and the cooriented derivative match, these powers cancel on the part of the cut \emph{inside the cone}. Vanishing power alone would not imply matching: two reflecting Dirichlet walls also have zero power.

Neither estimate produces a solution. For that purpose take smooth compactly supported displacement $a$, velocity $b$ at $t_0$, and smooth compact spacetime force $h$. Fourier transform only in $x$ and set $\omega_k=(|k|^2+m^2)^{1/2}$, $s_k(v)=\sin(\omega_k v)/\omega_k$, interpreted as $v$ at $\omega_k=0$. The scalar oscillator at each $k$ has solution
\begin{align}
 \widehat u(t,k)={}&\cos(\omega_k(t-t_0))\widehat a(k)+s_k(t-t_0)\widehat b(k)\notag\\
 &+\int_{t_0}^t s_k(t-\tau)\widehat h(\tau,k)\,d\tau .\label{f:wave-existence}
\end{align}
Indeed $s_k(0)=0$, $s_k'(0)=1$ and $s_k''+\omega_k^2s_k=0$: the second time derivative of the integral produces exactly $\widehat h(t,k)$. On a bounded time interval the transforms of the data and force decay faster than any power of $|k|$, uniformly in the integration time. The kernels and their time derivatives grow at most polynomially in $|k|$ there. Thus the inverse transform and all the derivatives needed for the equation and initial conditions converge absolutely. This constructs the smooth solution and the cone argument proves its support bound. Choosing zero data at a time before the source gives the retarded solution. By approximation, \eqref{f:energy-estimate} also constructs limits in $C_t H^1_x\cap C^1_t L^2_x$ for $a\in H^1$, $b\in L^2$ and $h\in L^1_tL^2_x$ on finite intervals; their equation holds on compact tests by passage to the limit. Here continuity in time, in the indicated spatial norms, is a consequence of evolution. It is not the trace of a general spacetime $L^2$ field.

\section{One-sided inverses and adjoint domains}\label{f:green-domains}

Let $M$ be smooth, time oriented, globally hyperbolic and without boundary, $V\to M$ a smooth finite-rank vector bundle, and $P$ a smooth normally hyperbolic operator on it. We use the global Cauchy and Green theorems of B\"ar, Ginoux and Pf\"affle \cite[Thm.~3.2.11 and Cor.~3.4.3]{BGP}: compactly supported smooth data $(u,\nabla_nu)$ on a smooth spacelike Cauchy hypersurface and compactly supported smooth forcing give a unique smooth solution supported in $J^+(K)\cup J^-(K)$, where $K$ is the union of the three supports. Here $n$ is the future unit normal and one may use the connection belonging to the operator (for a scalar, the normal derivative). On compact smooth spacetime sources there are unique retarded and advanced Green operators, mapping $C_c^\infty(M;V)$ to $C^\infty(M;V)$, characterized by
\begin{equation}
 \begin{gathered}
  PG^{R/A} f=f,\qquad G^{R/A}P v=v \quad(f,v\in C_c^\infty(M;V)),\\
  \operatorname{supp}G^R f\subset J^+(\operatorname{supp}f),\qquad
  \operatorname{supp}G^A f\subset J^-(\operatorname{supp}f).
 \end{gathered}\label{f:green-identities}
\end{equation}
In that reference the leading differential expression is $-g^{\mu\nu}\partial_\mu\partial_\nu$, and the future-supported map is called $G_+$ (``advanced''). For our scalar $P$ one applies their theorem to $-P$ and negates its Green operators; our word \emph{retarded} specifies future support of the response. The flat-space proof above illustrates the theorem, not its proof for curved metrics.

The output on a compact source need not be compact. To apply a retarded inverse a second time, we need a class of supports that the first inverse preserves. Call a closed set $A$ \emph{past compact} if $A\cap J^-(C)$ is compact for every compact $C\subset M$; reverse the sign for \emph{future compact}. Write $C^\infty_{\rm pc}(M;V)$ and $C^\infty_{\rm fc}(M;V)$ for the corresponding spaces of smooth sections. We use the geometric fact that $J^+(K)\cap J^-(C)$ is compact for compact $K,C$ in a globally hyperbolic spacetime \cite[Lemma~A.5.7]{BGP}. In particular a field supported in $J^+(K)$ for compact $K$ is past compact. Differential operators preserve either support condition because they do not enlarge supports.

\begin{proposition}[One-sided extension]\label{f:one-sided-extension}
Under the same boundaryless, globally hyperbolic hypotheses, $P$ is invertible on $C^\infty_{\rm pc}(M;V)$, with inverse extending $G^R$ and satisfying $\operatorname{supp}G^R f\subset J^+(\operatorname{supp}f)$. On $C^\infty_{\rm fc}(M;V)$ its inverse extends $G^A$ with the past support bound.
\end{proposition}
\begin{proof}
Let $A=\operatorname{supp}f$ be past compact and $O$ a relatively compact open neighborhood of an observation point. Choose $\chi\in C_c^\infty(M)$ equal to one near the compact set $A\cap J^-(\overline O)$. On $O$ define the answer as $G^R(\chi f)|_O$, using only the compact-source inverse. For any two such cutoffs the difference of the cut-off sources has support disjoint from $J^-(\overline O)$, so its retarded response vanishes on $O$ by \eqref{f:green-identities}. Shrinking $O$ gives the same construction on overlaps, so the local answers fit together to a smooth response $u$. It is linear: one cutoff can be used for the union of two past-compact supports. On each $O$ we have $Pu=\chi f=f$ and $\operatorname{supp}u\cap O\subset J^+(A)$. To see that the output is past compact, note that
\[
 J^+(A)\cap J^-(C)=J^+\big(A\cap J^-(C)\big)\cap J^-(C)
\]
is a compact diamond for each compact $C$. In particular $u$ is in the domain for another application.

For the other inverse identity take a past-compact \emph{field} $v$ and choose $\chi$ as above for its support. It also cuts off $Pv$ correctly, since $\operatorname{supp}Pv\subset\operatorname{supp}v$. The compact fields $P(\chi v)$ and $\chi Pv$ differ only away from $J^-(\overline O)$: wherever $v$ can contribute in that past, $\chi$ is locally one. Hence on $O$ the extended $G^R Pv$ equals $G^R P(\chi v)=\chi v=v$. This proves both inverse identities and uniqueness on the past-compact space. Reversing the causal directions proves the advanced statement.
\end{proof}

Without a support condition $G^RPv-v$, even if defined, can contain a free solution with nonzero past data. On a finite slab one needs $Pv$ in the chosen source space, $v$ satisfying the wall condition, and zero Cauchy data at the past cap for the retarded identity or the future cap for the advanced one. Compact support away from both caps suffices for the cap conditions; mere compactness up to the caps does not. Smooth sources can be taken to vanish near the relevant initial cap, avoiding compatibility conditions at its corners. Likewise, a physical timelike boundary requires a separate, well-posed initial--boundary problem. At a wall the one-sided extension needs causal uniqueness and propagation for that boundary problem, with compositions acting in its wall domain. Even $\chi v$ in the cutoff proof may fail a normal-derivative boundary condition unless $\chi$ is chosen compatibly. The tempered inverse of $1-\partial_x^2$ in Chapter~\ref{f:distributions} illustrates a different choice altogether: its kernel $e^{-|x|}/2$ is nonzero on both sides of the source, with no propagation direction selected.

For the scalar, represent source densities using $dV_g$, and write $(u,v)_M=\int\overline u v\,dV_g$ when the supports have compact intersection. The \emph{formal adjoint} $P^\dagger$ is defined by $(Pu,v)_M=(u,P^\dagger v)_M$ on compact interior tests; a fibre inner product gives the analogous notation for bundles. For \eqref{f:wave-operator}, the difference $\overline{Pu}\,v-\overline u\,Pv$ is the divergence of $(\nabla\overline u)v-\overline u\nabla v$, so $P^\dagger=P$. Put $u=G_P^R f$ and $v=G_{P^\dagger}^A h$. Their supports overlap only in $J^+(\operatorname{supp}f)\cap J^-(\operatorname{supp}h)$, compact by global hyperbolicity. A cutoff equal to one near this intersection permits integration by parts with no end term, giving
\[
 (G_P^R f,h)_M=(u,P^\dagger v)_M=(Pu,v)_M=(f,G_{P^\dagger}^A h)_M .
\]
At a wall one needs \emph{adjoint boundary conditions} to cancel the remaining boundary integral; formal symmetry of the interior expression does not do that. With matching adjoint domains and formally self-adjoint $P$, on the common source space the \emph{causal propagator}
\begin{equation}
 E=G^R-G^A\qquad\text{satisfies}\qquad PEf=0,\quad EPv=0,\quad(Ef,h)_M=-(f,Eh)_M .\label{f:causal-propagator}
\end{equation}
In the boundaryless problem its support on $f$ lies in $J^+(K)\cup J^-(K)$, $K=\operatorname{supp}f$. For a spatial oscillator mode of the scalar $P$, the kernels are $-\Theta(t)s_k(t)$ and $\Theta(-t)s_k(t)$: their slopes both jump by $-1$ at zero, so applying $-(\partial_t^2+\omega_k^2)$ gives $\delta_0$. Their difference is $-s_k(t)$, a homogeneous solution with the sign fixed by the source equation.

\begin{exercise}[The adjoint of an absorbing wall]
On a static spatial wall take $\partial_\nu u+u_t=0$, with $\nu$ outward. Compute the power leaving the bulk and the adjoint wall condition. Assume the variations vanish near the time caps.
\end{exercise}
\noindent\emph{Solution.} The contribution to the derivative of bulk energy is $\int u_t\partial_\nu u=-\int u_t^2$ for a real field. On the wall the adjoint boundary integral for $P$ is $\int ((\partial_\nu\overline u)v-\overline u\partial_\nu v)$. Substituting $-\overline u_t$ and integrating it by parts in time leaves $\int\overline u(v_t-\partial_\nu v)$. Thus the adjoint condition is $\partial_\nu v-v_t=0$, with the opposite sign. Using the forward absorbing law for both fields would leave a boundary term.

\section{Determining a retarded interface history}\label{f:interface-history}

Cut a linear problem in two. Keep its physical outer boundary conditions and choose, at the artificial face on each piece, an auxiliary problem that is solvable with one prescribed trace. Denote the pair of retarded particular solutions with frozen zero trace by $G_0f$; a pair of homogeneous solutions with adjustable boundary histories $\eta$ and zero past data is $H\eta$. We require $G_0f+H\eta$ to parametrize the regional responses in the declared class. The map $H$ here solves an equation: the Sobolev lift \eqref{f:trace-lift} was not required to do so. Let $\mathcal B$ measure transmission error. For an ordinary scalar joining two media this is the difference of the values and the sum of the outward fluxes; if a common value has already been imposed only the flux row remains. Put
\[
 Nf=\mathcal BG_0f,\qquad M\eta=\mathcal B H\eta .
\]
The adjustable history must satisfy $M\eta=-Nf$. If a prescribed sheet load calls for nonzero mismatch $b$, the right hand side is $b-Nf$. In particular the number of measurements is no guarantee that arbitrary measurements can be prescribed.

\begin{proposition}[Regional reconstruction]\label{f:causal-sewing}
Fix source, history and mismatch spaces for the retarded problems above. Suppose zero mismatch assembles the regional fields into a global solution of the stated regularity, without a sheet source. Suppose $\operatorname{ran}N\subseteq\operatorname{ran}M$ and there is a continuous linear $L$ on $\operatorname{ran}N$ with $MLy=y$ there. Assume the regional maps are continuous and the zero-past homogeneous global solution is unique. With $K=\operatorname{supp}f$ and $\Gamma$ the cut, require $\operatorname{supp}G_0 f\subset J^+(K)$, $\operatorname{supp}LNf\subset\Gamma\cap J^+(K)$, and $\operatorname{supp}H\eta\subset J^+(\operatorname{supp}\eta)$, using the propagation cones of the assembled problem. Then the global retarded response on the chosen sources is
\begin{equation}
 G^R_{\rm joined} f=G_0f-HLNf .\label{f:sewn-inverse}
\end{equation}
\end{proposition}
\begin{proof}
The correction satisfies the homogeneous interior equations and the same homogeneous outer conditions, so it changes no bulk forcing. The mismatch becomes $Nf-MLNf=0$. A signal from $LNf$ reaches only points in $J^+(\operatorname{supp}LNf)\subset J^+(K)$ by transitivity of causal reachability. The assembled field therefore has the required support, and uniqueness identifies it. Two allowed choices of $L$ differ on a given mismatch by an element of $\ker M$, whose lifted difference would be a homogeneous joined solution with zero past, hence zero. This argument needs a right inverse on the generated errors, not a two-sided inverse on an arbitrary set of boundary measurements.
\end{proof}

Consider two strings of positive mass densities $\rho_\pm$ and tensions $T_\pm$, joined at $x=0$. Use the positive equation $\rho u_{tt}-\partial_x(Tu_x)=h$. Give the two ends a common displacement. By \eqref{f:jump-formula} the sheet term in the positive equation is $(T_-u_{-,x}-T_+u_{+,x})\delta(x)$; the Euler operator of opposite sign has the opposite sheet term. Continuity of $Tu_x$ therefore joins piecewise smooth solutions without a singular load. Put $c_\pm=\sqrt{T_\pm/\rho_\pm}$ and $Z_\pm=T_\pm/c_\pm$; the second constant is called the impedance. Start with smooth compact forces supported off the cut after $t=0$, zero initial data, and no incoming wave from infinity. The frozen Dirichlet solutions exist on each half-line: extend the force oddly across the end, solve on the full line by \eqref{f:wave-existence} after rescaling the speed and force, and restrict the odd solution. For a common smooth history $q$ vanishing for $t\leq0$, the outgoing homogeneous fields are
\[
 H_-q(t,x)=q(t+x/c_-),\quad x<0;\qquad H_+q(t,x)=q(t-x/c_+),\quad x>0 .
\]
Both outward weighted derivatives $T\partial_\nu$ now have the same sign at the cut. Their sum and its retarded inverse are
\begin{equation}
 Mq=(Z_-+Z_+)q',\qquad (Ly)(t)=\frac{1}{Z_-+Z_+}\int_{-\infty}^t y(\tau)\,d\tau .\label{f:retarded-history-lift}
\end{equation}
These maps act continuously on smooth histories zero for $t\leq0$, with convergence of every derivative on compact time intervals. In particular $M$ is onto there. If one also required $q$ compact in time, any $y$ of nonzero integral would be out of range. Even on a single half-line, prescribing an outward force $g$ as well as displacement $q$ in the homogeneous experiment requires $g=Zq'$: the range of the pair is a graph, not a product. For a sourced problem one must first subtract the frozen particular force. This is the overdetermination that a trace theorem alone cannot see.

For a force confined to $y<0$, we can also compute the actual frozen error $Nh$. Let $v$ be the frozen left solution and $w=v_t-c_-v_x$. Then $(\partial_t+c_-\partial_x)w=h/\rho_-$. Integrate along the ray $y=-c_-(t-\tau)$ ending at the face at time $t$. There $v_t=0$ by the frozen condition, so
\[
 Nh(t)=T_-v_x(t,0)=-\frac{T_-}{\rho_-c_-}\int_{-\infty}^t h(\tau,-c_-(t-\tau))\,d\tau
     =-\int_{-\infty}^0 h(t+y/c_-,y)\,dy .
\]
The right-hand field is $H_+(-LNh)$. At a detector $x>0$ this gives
\begin{equation}
 u_+(t,x)=\frac{1}{Z_-+Z_+}\int_{-\infty}^0\!dy\int_{-\infty}^{\,t-x/c_+ - |y|/c_-} h(\tau,y)\,d\tau .\label{f:string-transmission}
\end{equation}
The integration limit adds the travel times in the two media. The same half-line support bounds apply on the left, including to the correction radiated back from the cut. A homogeneous difference with zero Dirichlet trace on a half-line has zero boundary power, hence vanishes for zero past data. This gives uniqueness for each regional parametrization. For the joined problem use energy densities $(\rho u_t^2+T u_x^2)/2$ on the two strings: the internal terms are the common $u_t$ times the sum of outward weighted derivatives, hence cancel. The hypotheses of the reconstruction therefore hold in this piecewise smooth class. In a medium with a genuine jump in coefficients, a reflected wave can be part of the joined solution.

\begin{exercise}[Measuring the impedance jump]
Send towards the cut a right-moving pulse $i(t-x/c_-)$ on the left with $i\in C_c^\infty((0,\infty))$. At time zero this specifies compact data strictly to the left. Find the reflected and transmitted displacements with no incoming wave on the right.
\end{exercise}
\noindent\emph{Solution.} If $r(t+x/c_-)$ is reflected and $q(t-x/c_+)$ transmitted, continuity gives $r=q-i$ at the face. The force sum is $Z_-(-i'+r')+Z_+q'=0$. With zero past interface history this gives
\[
 q=\frac{2Z_-}{Z_-+Z_+}i,\qquad r=\frac{Z_--Z_+}{Z_-+Z_+}i .
\]
In particular equal impedances remove the reflection even if the two speeds differ; the arrival time on the far side still measures its speed.

\section{Eliminating one of two cuts}\label{f:delayed-schur}

If two cuts carry independent admissible histories $q,r$, the two transmission equations can have the form $Aq+Br=b$, $Cq+Dr=d$. The blocks are now operators on histories. If $A$ has a causal inverse on the required domain, the first row gives $q=A^{-1}(b-Br)$ \emph{while $r$ is still free}. Substitution gives the Schur complement equation $(D-CA^{-1}B)r=d-CA^{-1}b$. This is equivalent to both rows if the operations are defined and the eliminated history is restored. To solve uniquely by this route the remaining equation must also have a unique causal solution on the data in question. A time delay, for example, has a time advance as its formal inverse; formal invertibility alone would give the wrong support.

One can see the delay in an exact calculation. Take unit tension, density and speed on the whole line and insert cuts at $0$ and $\ell>0$. Use the Laplace transform $\widetilde q(z)=\int_0^\infty e^{-zt}q(t)\,dt$, initially on smooth compact driving histories vanishing near zero, $\Re z>0$. Between the cuts the homogeneous lift with endpoint values $\widetilde q,\widetilde r$ obeys $U''=z^2 U$, hence
\[
 U(x)=\frac{\sinh(z(\ell-x))\widetilde q+\sinh(zx)\widetilde r}{\sinh(z\ell)} .
\]
The exterior half-lines contribute $z\widetilde q$ and $z\widetilde r$ to the two outward sums. Adding $-U'(0)$ and $U'(\ell)$ gives equal diagonal blocks $A=D=z(1+\coth(z\ell))$ and off-diagonal blocks $B=C=-z/\sinh(z\ell)$. Writing the hyperbolic functions in exponentials gives
\[
 A^{-1}=\frac{1-e^{-2z\ell}}{2z},\qquad A^{-1}B=-e^{-z\ell},\qquad D-CA^{-1}B=2z .
\]
Thus the remaining equation is $2z\widetilde r=\widetilde d+e^{-z\ell}\widetilde b$, not $2z\widetilde r=\widetilde d$. In time, if $b,d$ denote the required corrections to the two outward sums, the result is
\[
 r(t)=\tfrac12\int_{-\infty}^t \big(d(\tau)+b(\tau-\ell)\big)\,d\tau .
\]
Indeed $e^{-z\ell}$ translates a zero-past input to the right by $\ell$, and $1/z$ integrates it from the past, by substitution and integration by parts in the defining transform. The first pivot also has a causal inverse, half the integral over $[t-2\ell,t]$. These time formulas extend to smooth zero-past data on each finite interval; the direct blocks expand in delayed derivatives using $(1-e^{-2z\ell})^{-1}=\sum_{j\geq0}e^{-2jz\ell}$, with only finitely many arrivals on that interval. Eliminating a cut removed an unknown from the equation. It did not shorten the time a wave takes to cross the region.
% !TeX root = ../main.tex
\chapter{Complexes and the inverse of a gauge equation}\label{f:gauge-complex}
The scalar wave operator acquired its retarded inverse by imposing zero incoming data. In electromagnetism the same prescription does not distinguish potentials differing by the derivative of a gauge parameter of compact support. This is not a technical failure of an estimate: the operator is required by symmetry to kill a whole family of fields. What can replace an inverse? We first answer a question of linear algebra, for which neither a metric nor causality is needed.

\section{Cohomology and homotopy}
A graded vector space is a family of spaces $C^k$ indexed here by integers. A cochain complex is such a family with maps $D_k:C^k\to C^{k+1}$ satisfying $D_{k+1}D_k=0$. Write simply $D$ when the degree of the argument is known. A cocycle is a vector in $\ker D_k$; a coboundary is a vector in $\operatorname{im}D_{k-1}$. The zero composition is exactly what makes the quotient $H^k(C)=\ker D_k/\operatorname{im}D_{k-1}$ well defined. It records cocycles which cannot be removed by changes coming from one degree below. For instance $d^2=0$ makes differential forms a cochain complex. The angular one-form of Chapter~\ref{f:forms} represents a nonzero class, as its period detects. There is no requirement that a complex be exact: exactness at degree $k$ is the \emph{additional conclusion} $H^k=0$.

A cochain map $T:C\to C'$ consists of degree-preserving maps commuting with the differentials, so it maps closed vectors to closed vectors and boundaries to boundaries. If two cochain maps satisfy $T-S=D' h+hD$ with maps $h:C^k\to C'^{k-1}$, they are called chain homotopic. On a closed vector $v$ the difference is the boundary $D'hv$, so the induced maps on cohomology coincide. In particular if
\begin{equation}
 Dh+hD=1,\label{f:contraction}
\end{equation}
then every closed vector is exact: $v=D(hv)$ for $Dv=0$. The map $h$ is a contracting homotopy. With $1-\Pi$ on the right, where $\Pi^2=\Pi$ commutes with $D$, the same reasoning says every class has a representative in the retained part $\Pi C$. It does not authorize one to forget that part. A two-term complex $C^0\xrightarrow{P}C^1$ with no other nonzero terms has \eqref{f:contraction} exactly when $h_1$ is a two-sided inverse of $P$. Moving to a longer complex lets neighboring maps absorb the parts a single inverse cannot handle.

\section{What symmetry implies about an equation}
Let $K$ send infinitesimal proper gauge parameters to field variations at a solution, $J$ denote the linearized Euler operator, and $K^*$ the adjoint relative to the field--source and parameter--parameter-density pairings. Gauge-related fields are solutions together, so $JK=0$. Invariance of the action for compactly supported local gauge parameters gives the Noether identity $K^*\mathcal E=0$ (integrate the derivatives of the parameter by parts and use its arbitrariness). Linearizing at $\mathcal E=0$ gives $K^*J=0$. If a boundary is present, these statements require that the transformations preserve the boundary conditions and that the adjoint pairing leave no unwanted surface term. Nonzero endpoint or surface charges are not being declared zero by this argument.

For an irreducible linearized theory the resulting portion of the deformation complex is
\begin{equation}
 U\ \xrightarrow[{-1\ \to\ 0}]{K}\ V\
 \xrightarrow[{0\ \to\ 1}]{J}\ V^\vee\
 \xrightarrow[{1\ \to\ 2}]{K^*}\ U^\vee .\label{f:deformation-complex}
\end{equation}
The symbols now denote section spaces on stated support and boundary domains, with dual \emph{densities} for sources, not algebraic duals containing arbitrary functionals. At degree zero $\ker J/\operatorname{im}K$ is the space of linearized solutions up to proper gauge; at degree one a closed source $f$ obeys the compatibility condition $K^*f=0$. Each differential must map its claimed domain into the next. For example, $K\lambda=d\lambda$ in Maxwell theory. If the tangential potential is fixed on a wall, $j^*d\lambda=d(j^*\lambda)$ must vanish there; a parameter with zero trace suffices. Permission to change a frame in the bulk does not settle the wall domain.

This grading is particularly natural in the Batalin--Vilkovisky description \cite{CMR,Khavkine}. An ordinary gauge parameter is encoded by an odd variable called a ghost (an odd variable anticommutes with other odd variables). Fields and ghosts also have dual variables called antifields; in the usual function grading a field of degree $k$ has antifield of degree $-1-k$. An odd differential $Q$ has $Q^2=0$ and acts by the graded product rule $Q(ab)=(Qa)b+(-1)^{|a|}a\,Qb$. At a zero of $Q$ its derivative is a square-zero map on the tangent complex. The degrees indicated on \eqref{f:deformation-complex} are those of \emph{deformations}, with parameters one degree below field tangents, not degrees being assigned to ordinary commuting sections as if they were ghost coordinate functions. For a simple model the gauge part of the differential on functions sends an abelian potential $A$ to $dc$ and its ghost $c$ to zero. Thus it squares to zero on $A$ and leaves $dA$ invariant. Equations of motion require additional components, represented by antifields in the BV formulation. This describes what the linear complex measures without assuming existence of a full quantum BV theory with boundary. In a differential graded description of Ward identities, an exact local violation can be removed by changing a local functional; a violation in nonzero cohomology is a different assertion.

\section{Gauge fixing constructs a homotopy}
Choose a gauge condition $C:V\to U$ and a symmetric invertible weight $B:U\to U^\vee$. Form $P=J+C^*BC$, $R=CK$, and $R^*=K^*C^*$. Assume $J$ is formally symmetric on the paired domains. A source for the field can be solved using $P$, but only if the gauge condition and parameter equation can be solved as well. Fix either a retarded or an advanced support problem. Suppose the indicated operators admit two-sided inverses $G_P,G_R,G_{R^*}$ on compatible one-sided domains, with the required uniqueness and boundary data, and all the following compositions are defined there. These hypotheses involve cap and wall conditions as well as algebra.

\begin{proposition}\label{f:green-homotopy}
On that complex a contracting homotopy has three possibly nonzero components, on degrees $0,1,2$ respectively:
\[
 h_0=G_R C,\qquad h_1=G_P,\qquad h_2=C^*G_{R^*} .
\]
\end{proposition}
\begin{proof}
From $JK=0$ and $K^*J=0$ obtain $PK=C^*BR$ and $K^*P=R^*BC$. With the inverse, support and domain assumptions these identities imply $KG_R=G_PC^*B$ and $BCG_P=G_{R^*}K^*$. On $V$ one thus has $Kh_0+h_1J=G_P(C^*BC+J)=1$. On $V^\vee$ one has $Jh_1+h_2K^*=(J+C^*BC)G_P=1$. On the two remaining spaces $h_0K=G_RR=1$ and $K^*h_2=R^*G_{R^*}=1$. These are all four degrees of \eqref{f:contraction}.
\end{proof}

\begin{example}[Changing the gauge condition without changing the response class]\label{f:gauge-matrix}
Use Euclidean pairings on $U=\mathbb R$, $V=\mathbb R^2$. Set $K\lambda=(\lambda,0)$, $J=\operatorname{diag}(0,q)$, $q\ne0$. A compatible source is $(0,s)$; all its preimages under $J$ have second coordinate $s/q$. Choose $C(x,y)=a x+d y$ with $a\ne0$, and $B=b\ne0$. Then
\[
 P=\begin{pmatrix}ba^2&bad\\ bad&q+bd^2\end{pmatrix},\qquad
 \det P=ba^2q,\qquad P^{-1}(0,s)^{\mathsf T}
   =\big(-ds/(a q),\,s/q\big)^{\mathsf T}.
\]
Changing $d$ moves the first coordinate by a vector in $\operatorname{im}K$ and does not change the value seen by another closed source. Conversely a test source $(1,0)$ would read a spurious dependence on $d$. It was never closed. At $a=0$ the gauge condition fails to slice the orbit: $CK=0$ and the proposed inversion fails. At $q=0$ there is a different problem, an actual residual physical direction.
\end{example}

\section{A differential complex on spacetime}\label{f:hodge}
On oriented pseudo-Riemannian spacetime the metric gives alternating covectors a nondegenerate pairing, generally indefinite. The Hodge star $*:\Lambda^kT^*M\to\Lambda^{n-k}T^*M$ is fixed by $\alpha\wedge *\beta=\langle\alpha,\beta\rangle_g\,\epsilon_g$, where the $k$-form pairing includes the coordinate factor $1/k!$. To see existence, use an oriented orthonormal coframe: the star sends a wedge of distinct coframe covectors to the complementary wedge with the sign from the metric times the permutation sign. This description also gives $*^2=(-1)^{k(n-k)+1}$ in Lorentz signature with one minus sign. Write $d^*$ for the formal adjoint of $d$ under the integrated metric pairing on compactly supported forms without boundary. On one-forms, $d^*A=-\nabla^a A_a$ by Stokes. It is neither pointwise complex conjugation nor the field-space differential $\delta$. Adjointness and $d^2=0$ imply $(d^*)^2=0$.

Take $K=d$ on scalar parameters, $J=-d^*d$ on one-form potentials, $C=d^*$, and $B=-1$ after metric identification of dual densities. The Maxwell equation $JA=f$ requires $d^*f=0$, the conservation of its source. The gauge-fixed field and parameter operators are $P=-W_1$ and $R=W_0$ with $W_k=dd^*+d^*d$ on $k$-forms. Up to overall sign these have metric wave principal part; on globally hyperbolic spacetimes without boundary apply the wave theorem of Chapter~\ref{f:waves} and the one-sided extension on p.~\pageref{f:one-sided-extension}. For compact input supported in $K$, the retarded outputs and all their derivatives have support in $J^+(K)$, and the advanced ones in $J^-(K)$. These are the larger spaces used when a second inverse is applied. Write $S_k$ for the same-side causal inverse of $W_k$. Then $W_1 d=dW_0$ since $d^2=0$, and uniqueness with the same one-sided data gives $S_1d=dS_0$. Consequently the three homotopy blocks are $S_0d^*$, $-S_1$, and $dS_0$. On one-forms, for instance, $dS_0d^*+S_1d^*d=S_1W_1=1$. With a wall, $d$ might take a scalar in its allowed domain to a one-form outside the field domain. The algebraic identity $W_1d=dW_0$ would still hold, but the inverse intertwining conclusion would not follow from it.

Why use two one-sided homotopies? On their common compact-source input one has $D(h^R-h^A)+(h^R-h^A)D=0$. A closed source is therefore sent to a homogeneous field; changing the source by an exact one changes the field by an exact one (with the sign fixed by this identity). This is the replacement for saying that the causal propagator descends from test sources to physical solutions. If two gauge choices admit homotopies on the \emph{same} retarded complex, their two retarded responses on a closed source differ by a closed field $w$. Then $w=D\widehat h w$ is a proper gauge change there, and similarly on the advanced complex. Pairing with another closed source does not distinguish the two choices. This proof requires precisely the domain where the homotopy on $w$ and the adjoint pairing exist; a residual projector or a surface charge changes the answer.

\section*{Exercises}
\begin{exercise}In Example~\ref{f:gauge-matrix} set $q=0$. Find all the cohomology groups of the four-term complex and explain why a contraction on the whole complex is impossible.\end{exercise}
\noindent\emph{Solution.} $K$ is still injective and $K^*$ onto. At degree zero $\mathbb R^2/\mathbb R(1,0)$ is one-dimensional; at degree one $\ker K^*=\mathbb R(0,1)$ is also one-dimensional since $J=0$. Endpoint cohomologies vanish. Equation~\eqref{f:contraction} would make both middle groups vanish.
\begin{exercise}In a finite-dimensional complex with a positive inner product let $\Delta=DD^*+D^*D$. Prove that $\ker\Delta=\ker D\cap\ker D^*$. If $G$ is the inverse of $\Delta$ on the orthogonal complement of this kernel and zero on the kernel, construct a homotopy to orthogonal projection onto the kernel.\end{exercise}
\noindent\emph{Hint.} Pair $\Delta v$ with $v$. Both $D$ and $D^*$ commute with $\Delta$ on the full graded space and preserve its eigenspaces. Hence $h=D^*G$ satisfies $Dh+hD=\Delta G=1-\Pi$. No causal support condition or indefinite inner product is present in this exercise.
\begin{exercise}\label{f:gauge-boundary-ex}On a spatial region with boundary take a scalar charge density constraint $\partial_iE^i=0$. Pair a variation $\partial_i\delta E^i$ with $\lambda$ and integrate by parts. Why might a boundary value of $\lambda$ be physically measurable although the bulk transformation is a gradient?\end{exercise}
\noindent\emph{Check.} $\int\lambda\,\partial_i\delta E^i=-\int (\partial_i\lambda)\delta E^i+\int_{\partial}\lambda\,\delta E^n$. If the surface flux is variable and the parameter is allowed to be nonzero there, its pairing need not vanish. Setting the proper parameter to zero there, or imposing a flux condition, changes this test. Chapter~\ref{f:symplectic} gives its interpretation as a charge rather than a gauge degeneracy.

\part{Classical and quantum observables}
% !TeX root = ../main.tex
\chapter{Symplectic reduction and classical observables}\label{f:symplectic}
Suppose two angular readings determine a position only through their sum. Opposite changes in the two angles leave the position unchanged. Do they change anything measurable? If the conjugate momenta can vary independently, the relative rotation has a measurable generator. If the momenta must agree, the same rotation pairs to zero with every allowed displacement. To distinguish the two cases we start with a linear question: given an observable $H$, can we find a vector $X$ such that pairing any displacement $Y$ with $X$ gives the change $\delta H(Y)$? For the wave equation we shall solve this inverse problem without phase-space coordinates, using compact sources.

\section{What an alternating form can detect}
Let $V$ be a finite-dimensional real vector space. An alternating bilinear form $b$ satisfies $b(x,x)=0$; applying this to $x+y$ gives $b(x,y)=-b(y,x)$. We put the unknown vector in the \emph{second} slot: $b^\flat(x)=b(\,\cdot\,,x)\in V^*$. There are three different spaces to keep track of. For a subspace $W\subset V$ its \emph{annihilator} is $W^\circ=\{\lambda\in V^*: \lambda|_W=0\}$. Its orthogonal with respect to $b$ is $W^\perp=(b^\flat)^{-1}(W^\circ)\subset V$. The \emph{radical} of the whole form is $R=V^\perp$. Thus an annihilator consists of covectors, whereas a radical consists of vectors. A nondegenerate alternating form, one for which $R=0$, is called symplectic.

\begin{lemma}\label{f:alternating-quotient}
The equation $b(\,\cdot\,,x)=\lambda$ has a solution exactly for $\lambda\in R^\circ$; its solutions then form a translate of $R$. The form induced on $V/R$ is symplectic. If $V$ itself is symplectic, $(W^\perp)^\perp=W$ and the radical of the restriction to $W$ is $W\cap W^\perp$.
\end{lemma}
\begin{proof}
Skewness puts the image of $b^\flat$ in $R^\circ$. Both spaces have dimension $\dim V-\dim R$, by rank--nullity for the first and extension of a basis of $R$ for the second. The kernel is $R$, giving the claimed ambiguity. Adding an element of $R$ to either entry leaves $b$ unchanged; a class pairing to zero with all classes has a representative in $R$. In the nondegenerate case $b^\flat$ is an isomorphism, so $\dim W^\perp=\dim V-\dim W$. The inclusion $W\subset (W^\perp)^\perp$ is then equality by dimension. A vector in $W$ pairs to zero with all of $W$ precisely when it also lies in $W^\perp$.
\end{proof}

In a nonzero symplectic space choose $e\ne0$, then $f$ with $b(f,e)=1$. The plane they span is already nondegenerate. It and its orthogonal have zero intersection and complementary dimensions; the orthogonal is nondegenerate as well, since a vector there annihilating the orthogonal annihilates all of $V$. Repetition decomposes $V$ into such planes, proving in particular that $\dim V$ is even. Imposing a condition that restricts displacements to $W$ can destroy this pairing even when it was perfect on $V$. The appropriate linear reduction is then $W/(W\cap W^\perp)$, not $V/W$.

\section{Hamiltonian fields and the Poisson bracket}
On a smooth manifold $M$ a two-form $\omega$ gives an alternating pairing on each tangent space, smoothly in its point of attachment. It is \emph{symplectic} if it is nondegenerate at every point and closed. We write $\delta$ for the exterior derivative on phase space, as distinct from differentiation on spacetime. A closed two-form allowed to have a kernel is \emph{presymplectic}. For instance on the cotangent bundle $\pi:T^*Q\to Q$ a point $(q,p)$ already contains a covector. Applying it to the projection of a displacement defines the canonical one-form $\theta(Y)=p(\dd\pi\,Y)$. In local coordinates, with repeated coordinate indices summed,
\[
 \theta=p_i\,\delta q^i,\qquad \omega=\delta\theta=\delta p_i\wedge\delta q^i .
\]
Closedness here follows from $\delta^2=0$; nondegeneracy follows by testing independently with displacements of $p$ and of $q$. A general symplectic manifold need not have a globally defined potential $\theta$.

\begin{definition}\label{f:hamiltonian-slots}
The Hamiltonian vector field of a smooth real function $H$ and the bracket of two functions are specified by
\begin{equation}
 \delta H(Y)=\omega(Y,X_H),\qquad \{F,G\}=\delta F(X_G).
 \label{f:poisson-convention}
\end{equation}
On a symplectic manifold the field exists uniquely: the inverse of the smooth, pointwise invertible matrix of $\omega^\flat$ depends smoothly on its entries.
\end{definition}
In the cotangent coordinates, $\omega(Y,X_H)=Yp_i\,X_H q^i-Yq^i\,X_H p_i$. Reading off the coefficients gives Hamilton's equations $X_H q^i=\partial H/\partial p_i$, $X_H p_i=-\partial H/\partial q^i$ \cite{Arnold}. In particular $\{q^i,p_j\}=\delta^i_j$: the second entry generates the motion of the first. On a constant-rank presymplectic manifold a local Hamiltonian vector exists precisely when $\delta H$ annihilates the radical everywhere on the patch. Indeed the lemma gives pointwise solvability, and solving on a local smooth complement to the kernel gives a smooth choice. Two choices differ by a null field. A function varying along a null direction cannot have such a generator.

\begin{proposition}\label{f:poisson-jacobi}
The bracket on a symplectic manifold is a Poisson bracket: it is bilinear, antisymmetric, a derivation in each argument for multiplication of functions, and satisfies the Jacobi identity.
\end{proposition}
\begin{proof}
Bilinearity follows from the defining linear equations, and $\{F,G\}=\omega(X_G,X_F)$ gives antisymmetry. Differentiating a product gives the derivation rule. For Jacobi, $\iota_{X_F}\omega=-\delta F$, so Cartan's formula and closedness imply $\mathcal L_{X_F}\omega=0$. Therefore
\[
 \iota_{[X_F,X_G]}\omega=\mathcal L_{X_F}(\iota_{X_G}\omega)
      =-\delta(X_F G)=\delta\{F,G\}.
\]
Thus $[X_F,X_G]=-X_{\{F,G\}}$. Apply both sides to $H$: the result is $\{\{H,G\},F\}-\{\{H,F\},G\}=-\{H,\{F,G\}\}$, which is Jacobi after using antisymmetry. It was closedness, not merely an invertible skew matrix at each point, that supplied the differential identity.
\end{proof}

More generally a \emph{Poisson manifold} has a bracket on its smooth functions with the four properties of the proposition; no inverse two-form is required. For example, on $(q,p,s)\in\mathbb R^3$ take $\{F,G\}=F_qG_p-F_pG_q$. The symplectic planes $s=\text{constant}$ carry all the motions, and $s$ is a \emph{Casimir}, a function whose bracket with every function vanishes. This is a bracket defined on \emph{all} functions. Simply declaring the same three-dimensional space presymplectic with form $\delta p\wedge\delta q$ poses a different inverse problem for $s$ (Exercise~\ref{f:ex-null-coordinate}).

\section{Momentum constraints and the reduced space}
For a closed two-form of constant rank, the null spaces fit together into a distribution of tangent planes. If $Y,Z$ are null fields, Cartan's formula gives $\mathcal L_Y\omega=0$ and $\iota_{[Y,Z]}\omega=\mathcal L_Y\iota_Z\omega-\iota_Z\mathcal L_Y\omega=0$. The Frobenius theorem for a smooth constant-rank distribution closed under brackets \cite{LeeSM} therefore integrates these planes to leaves. A function with a Hamiltonian vector is constant on every connected null leaf. This local construction does not ensure that the set of leaves is a smooth manifold. A group action gives an important case in which the global quotient can be controlled.

Let a Lie group $G$ act smoothly, preserving $\omega$, on a symplectic manifold $M$. For $\xi$ in its Lie algebra $\mathfrak g$, write $\xi_M(x)$ for the velocity of $\exp(t\xi)x$ at $t=0$. Preservation of $\omega$ implies $-\iota_{\xi_M}\omega$ is closed, but it need not be exact. A \emph{momentum map} $J:M\to\mathfrak g^*$ supplies actual functions $J_\xi(x)=\langle J(x),\xi\rangle$ with $\delta J_\xi=-\iota_{\xi_M}\omega$. We assume equivariance: $J(gx)(\xi)=J(x)(\operatorname{Ad}_{g^{-1}}\xi)$, where $\operatorname{Ad}$ is obtained by differentiating conjugation in $G$. In particular $G$ preserves $C=J^{-1}(0)$.

\begin{theorem}[Regular reduction at zero]\label{f:regular-reduction}
Suppose $0$ is a regular value of $J$ and $G$ acts freely and properly on $C$. The orbit space $C/G$ has a unique symplectic form $\omega_{\rm red}$ such that $\pi_C^*\omega_{\rm red}=i^*\omega$, with $i:C\hookrightarrow M$ the inclusion and $\pi_C$ the orbit projection.
\end{theorem}
\begin{proof}
Regular means $\dd J_x$ is surjective at each $x\in C$, so $C$ is a submanifold with tangent space $\ker\dd J_x$. Write $O_x=\{\xi_M(x):\xi\in\mathfrak g\}$ for the orbit tangent. Since $\delta J_\xi(Y)=\omega(Y,\xi_M)$ for every $\xi$, we have $T_xC=O_x^\perp$. Equivariance gives $O_x\subset T_xC$. Lemma~\ref{f:alternating-quotient} now identifies the kernel of $i^*\omega$ at $x$ with $T_xC\cap O_x=O_x$.

We use the smooth quotient theorem for free proper actions \cite{Marsden}: freeness means no nonidentity element fixes a point, and properness means $(g,x)\mapsto(gx,x)$ is proper (compact sets have compact inverse images). The projection to the smooth quotient is a submersion whose differential has kernel $O_x$. The constraint and the orbit directions each cost $\dim G$ dimensions, giving $\dim(C/G)=\dim M-2\dim G$ when $C$ is nonempty. To pair two quotient tangents, lift them to $T_xC$ and pair the lifts. Changing a lift by an orbit tangent changes nothing; changing $x$ on its orbit changes nothing by invariance. Local smooth sections of the submersion show that the resulting form is smooth and closed. It is nondegenerate because a lift annihilating all lifts lies in $O_x$, and it is unique because every tangent pair has lifts.
\end{proof}
An invariant smooth $F$ on $M$ gives an observable $\bar F$ on $C/G$: $X_F$ is tangent to $C$ since $\delta J_\xi(X_F)=-\xi_M F=0$. For $Y\in TC$ the pullback identity gives
\[
 \omega_{\rm red}(\dd\pi_C\,Y,\dd\pi_C\,X_F)=\omega(Y,X_F)
                  =\delta\bar F(\dd\pi_C\,Y).
\]
Thus $\dd\pi_C\,X_F=X_{\bar F}$, and the bracket of invariant extensions can be computed before restriction, independently of the extensions. At a nonzero momentum value only the subgroup stabilizing that value preserves the level; replacing zero by an arbitrary value in $C/G$ would not in general define a quotient action.

An action symmetry and a null direction of a restricted form are not synonyms. In action language an off-shell symmetry can change the action by an endpoint difference $\alpha_X|_i^f$. If $\mathcal L_X\theta=\delta\alpha_X$ on the allowed domain, the Noether charge $H_X=\theta(X)-\alpha_X$ satisfies $\delta H_X=-\iota_X\omega$ by Cartan's formula \cite{HW}. A varying charge detects a physical motion; it does not identify two descriptions of the same point. Boundary conditions determine which variations can test this charge.

\begin{example}[Two angular coordinates, one invariant]\label{f:two-rotor}
On $T^*(S^1\times S^1)$ use angles $\alpha_1,\alpha_2$ modulo $2\pi$ and momenta $e_1,e_2$. The circle acting by $(\alpha_1+t,\alpha_2-t)$ is generated by $J=e_1-e_2$. Its zero level is regular, and the circle action is free and proper. On this level $e_1=e_2=e$; the orbits are exactly the fibers of $(\alpha_1,\alpha_2,e)\mapsto(w,e)$ with $w=\exp(i(\alpha_1+\alpha_2))$. Indeed any two decompositions of the same total angle differ by $(t,-t)$. Writing $w=e^{i\alpha}$ locally, the restricted potential is $e\,\delta\alpha$ and $\{w,e\}=iw$ by complex linearity. Its value is independent of the local lift of the angle. Rotating just the first circle is still a symmetry on the level, but it changes $w$ and its charge is $e$, not the vanishing $J$. Only the relative rotation has been divided out. Geometrically the two angles can record open transports along consecutive arcs of a circle; $w$ is the transport around the closed path.
\end{example}

\section{From compact sources to Hamiltonian solutions}
For fields there is no finite-dimensional rank argument identifying every covector with a vector. We construct the inverse on a stated domain. Take a time-oriented globally hyperbolic spacetime without boundary, of signature $(-,+,\ldots,+)$, with metric volume $\dd V$. The real scalar action on a slab is
\[
 S[\phi]=-\tfrac12\int\big(\nabla_\mu\phi\nabla^\mu\phi+m^2\phi^2\big)\,\dd V,
 \qquad P=\Box_g-m^2,\quad \Box_g=\nabla_\mu\nabla^\mu .
\]
Varying the integrand in a direction $Y$ gives $(P\phi)Y-\nabla_\mu(Y\nabla^\mu\phi)$; the linearization of the Euler term is therefore $P$ with the displayed sign. It is the action Hessian relative to $\dd V$. With $n$ the future unit normal, the final cap potential of Chapter~\ref{f:variation} is $\theta_\Sigma(Y)=\int_\Sigma(\nabla_n\phi)Y\,\dd\Sigma$. Thus the form on two variations is
\[
 \omega_\Sigma(u,v)=\int_\Sigma\big[(\nabla_n u)v-(\nabla_n v)u\big]\,\dd\Sigma .
\]
The divergence of $u\nabla^\mu v-v\nabla^\mu u$ is $uPv-vPu$. With the future cap contributing minus the contraction with $n_\mu$ in this signature, the slab identity is
\begin{equation}
 \int_{\Sigma_-}^{\Sigma_+}(uPv-vPu)\,\dd V
       =\omega_{\Sigma_+}(u,v)-\omega_{\Sigma_-}(u,v).
 \label{f:observable-green}
\end{equation}
We use it when the products have compact spatial support in the slab. For homogeneous solutions it proves slice independence.

Let $\mathcal D=C_c^\infty(M;\mathbb R)$ and $\langle f,u\rangle=\int_M fu\,\dd V$: $f$ represents a source density using this fixed measure. By the Green and Cauchy theorems for normally hyperbolic operators (Chapter~\ref{f:waves}; \cite{BGP}), there are unique retarded and advanced inverses $G^R,G^A$ on $\mathcal D$, supported in $J^+(\supp f)$ and $J^-(\supp f)$ respectively on input $f$. They invert $P$ on both sides on tests. Put $E=G^R-G^A$ and let $\mathcal S_{\rm sc}$ consist of smooth real homogeneous solutions supported in $J(K)=J^+(K)\cup J^-(K)$ for some compact $K$ (spacelike-compact solutions).

\begin{theorem}[Compact source presentation]\label{f:source-exact}
The sequence of real vector spaces
\[
 0\longrightarrow\mathcal D\xrightarrow{P}\mathcal D
       \xrightarrow{E}\mathcal S_{\rm sc}\longrightarrow0
\]
is exact: $P$ is injective on tests, $\ker E=P\mathcal D$, and $E$ reaches every solution in the indicated space.
\end{theorem}
\begin{proof}
A compactly supported homogeneous solution has zero data on an early Cauchy surface, so is zero. The inverse identities imply $PE=0$ and $EP=0$ on tests. If $Ef=0$, the common value $k=G^R f=G^A f$ satisfies $Pk=f$. Its support is contained in the compact causal diamond $J^+(\supp f)\cap J^-(\supp f)$, so $k$ is a test, not just a smooth preimage.

Given $u\in\mathcal S_{\rm sc}$, choose a smooth function $\chi$ of Cauchy time, zero before one level and one after a later level. Since $Pu=0$, the source
\[
 f=P(\chi u)=[P,\chi]u=(\Box_g\chi)u+2\nabla_\mu\chi\nabla^\mu u
\]
has support in $J(K)$ within that intervening closed slab. This set is compact: the future part is bounded between $K$ and the compact set $J^+(K)\cap\Sigma_+$, and the past part has the analogous earlier bound. Causal diamonds of compact sets are compact. Both $\chi u$ and $(\chi-1)u$ solve $Pv=f$; the first has zero early data, the second zero late data. Uniqueness therefore gives
\[
 G^R f=\chi u,\qquad G^A f=(\chi-1)u,\qquad Ef=u.
\]
Changing the switch in this time coordinate changes the representative by $P((\chi-\widetilde\chi)u)$ with compactly supported argument. A time cutoff alone would not make the source compact for an arbitrary spatially extended solution.
\end{proof}

The smeared observable is $F_f(\phi)=\langle f,\phi\rangle$ on solutions. Formal self-adjointness gives $F_{Pk}=0$ for $k\in\mathcal D$. To find its generator, take $Pu=0$, $v=G^R f$ in \eqref{f:observable-green} with caps before and after the source. The early retarded data vanish; on the late cap $v=Ef$ with its normal derivative. It follows that $\delta F_f(u)=\omega(u,Ef)$ on every slice for $u\in\mathcal S_{\rm sc}$. Taking instead $u=G^R f$, $v=G^A h$ makes both cap pairings vanish, yielding $\langle h,G^R f\rangle=\langle f,G^A h\rangle$. In particular $E$ is skew for the source pairing.

\begin{corollary}\label{f:peierls-labels}
The labels of these linear observables are precisely $V=\mathcal D/P\mathcal D$. Their nondegenerate Poisson pairing is
\begin{equation}
 \sigma([f],[h])=\{F_f,F_h\}_{\rm P}=\langle f,Eh\rangle
          =\omega(Eh,Ef)=-\omega(Ef,Eh).
 \label{f:label-sign}
\end{equation}
\end{corollary}
\begin{proof}
If $F_f$ vanishes on all spacelike-compact solutions, $0=\langle f,Eh\rangle=-\langle h,Ef\rangle$ for all tests $h$, so $Ef=0$. Exactness gives $f\in P\mathcal D$. The same test proves that the pairing has no radical on the quotient. The bracket follows from \eqref{f:poisson-convention} with $X_{F_h}=Eh$. Its minus sign in the last expression exchanges two arguments; it does not reverse the causal propagator.
\end{proof}
Since $E$ is onto, \eqref{f:label-sign} also proves nondegeneracy of $\omega$ on $\mathcal S_{\rm sc}$. It does not identify that space with its entire algebraic dual: the covectors inverted here are represented by compact smooth sources. Adding $-\varepsilon F_h$ to the action gives $Pv=h$ for its first-order response. Thus the Hamiltonian perturbation really is the retarded response minus the advanced one: this is the Peierls construction. The bracket extends by the product rule to finite polynomials in smearings. It makes no assertion about a product of distributions at the same point. With a timelike wall the variational domain, wall action and possible corner in $\theta$ must be specified: both causal problems have to exist and obey the adjoint Green identity on that domain. For gauge equations one must also restrict sources to annihilate proper gauge directions and justify an inverse modulo those directions \cite{Khavkine}. Neither extension follows from the boundaryless scalar sequence.

\section*{Exercises}
\begin{exercise}[Basic: a missing conjugate]\label{f:ex-null-coordinate}
For $\omega=\delta p\wedge\delta q$ on $\mathbb R^3$ with coordinates $(q,p,s)$, find all functions admitting a Hamiltonian vector field and all such fields for $H=p$. Compare with the Poisson bracket that ignores $s$.
\end{exercise}
\noindent\emph{Solution.} The radical is spanned by $\partial_s$, so $H_s=0$ is necessary and sufficient. For $H=p$ all choices are $\partial_q+c(q,p,s)\partial_s$. In particular $s$ admits no field for the two-form, whereas in the Poisson description it is a Casimir with zero Hamiltonian field.

\begin{exercise}[Intermediate: where regularity fails]
Rotate $\mathbb R^2$ about its origin, acting on the momentum vector by the same rotation. Show that $J=xp_y-yp_x$ generates the action. Identify the reduced form on $J=0$ where $r=(x^2+y^2)^{1/2}>0$. Why does the proof not describe the orbit of $(x,y,p_x,p_y)=0$?
\end{exercise}
\noindent\emph{Solution.} Hamilton's equations give $X_J=(-y,x,-p_y,p_x)$ in this ordering. In polar coordinates $p_x\delta x+p_y\delta y=p_r\delta r+J\delta\alpha$, where $p_r=(xp_x+yp_y)/r$. Dividing out the angle on the zero level leaves $\delta p_r\wedge\delta r$. At the origin $\dd J=0$ and every rotation fixes the point: both regularity and freeness fail there.

\begin{exercise}[Advanced: the sign measured by a switched source]
For $P=-(\partial_t^2+\nu^2)$, $\nu>0$, choose a switch $\chi$ as in Theorem~\ref{f:source-exact}. Construct tests $f_c,f_s$ for the solutions $\cos(\nu t),\sin(\nu t)$ and compute $\sigma([f_c],[f_s])$ by integration without using a propagator kernel.
\end{exercise}
\noindent\emph{Solution.} Write $c=\cos(\nu t)$ and $s=\sin(\nu t)$. The tests are $f_c=-\chi''c+2\nu\chi's$ and $f_s=-\chi''s-2\nu\chi'c$. The cutoff argument gives $Ef_s=s$. Integrating the $-\chi''cs$ term by parts leaves $\int f_c s\,\dd t=\nu\int\chi'(c^2-s^2+2s^2)\,\dd t=\nu$. On any time slice $\omega(c,s)=\dot c s-\dot s c=-\nu$, confirming which pairing belongs to the labels.
% !TeX root = ../main.tex
\chapter{Algebras, representations and states}\label{f:quantum}
The bracket in Chapter~\ref{f:symplectic} gives a number for each pair of linear observable labels. It does not give the expectation of either observable, much less the expectation of their product. To see what the bracket does determine, start with a real vector space $V$ and an alternating form $\sigma$. For the scalar, $v=[f]$ and $\sigma([f],[h])=\langle f,Eh\rangle$, not $\omega(Ef,Eh)$. We shall use this one input to construct two quantum algebras. A state is an extra choice; a Hilbert space can then be constructed rather than postulated. We display $\hbar>0$ in the defining relations and specify when natural units $\hbar=1$ are used.

\section{An algebra by generators and relations}
A complex unital algebra is a complex vector space with an associative bilinear multiplication and an identity $1$. An involution $*$ is conjugate-linear and satisfies $(AB)^*=B^*A^*$, $(A^*)^*=A$ and $1^*=1$. It models taking an adjoint; it is defined even before there are operators to take adjoints of. The elements $A=A^*$ are called Hermitian. Imposing a relation $R=0$ means also imposing $ARB=0$ for arbitrary $A,B$. Accordingly one divides by a \emph{two-sided ideal}, a linear subspace stable under multiplication on both sides. For an involution on the quotient one also needs $I^*=I$. These conditions make $[A][B]=[AB]$ and $[A]^*=[A^*]$ independent of representatives. Any unital $*$-homomorphism killing $I$ factors uniquely through the quotient, by sending $[A]$ to the image of $A$.

\begin{definition}[Polynomial canonical commutation relations]\label{f:polynomial-ccr}
Complexify $V$ to $V_{\mathbb C}$. Using the \hyperref[f:tensor-product]{algebraic tensor product}, form the tensor algebra $T(V_{\mathbb C})=\bigoplus_{n\ge0}V_{\mathbb C}^{\otimes n}$, with degree zero $\mathbb C$, product concatenation, and star conjugating entries and reversing their order. The polynomial CCR algebra is its quotient by the two-sided $*$-ideal generated, for $v,w\in V$, by
\[
 v\otimes w-w\otimes v-i\hbar\sigma(v,w)1 .
\]
The adjoint of each displayed relation is its negative. The classes $\Phi(v)$ of degree-one tensors are real-linear in $v$, are Hermitian for real $v$, and satisfy $[\Phi(v),\Phi(w)]=i\hbar\sigma(v,w)1$, where $[A,B]=AB-BA$.
\end{definition}
These are finite words and their finite sums, not convergent power series. The algebra is \emph{universal}: given Hermitian elements $T_v$ of another unital $*$-algebra with this linearity and commutator, multiply the $T_v$ in the order of each word and extend linearly. The generating relations map to zero, hence so does their ideal, and the map on the quotient is the unique homomorphism with $\Phi(v)\mapsto T_v$. The relations are consistent: on the vector space of commuting polynomials $\operatorname{Sym}(V_{\mathbb C})$ let $M_v$ multiply by $v$ and let the derivation $D_v$ send $w$ to $\sigma(v,w)$. Then $[D_v,M_w]=\sigma(v,w)1$; hence $M_v+(i\hbar/2)D_v$ obey the desired commutators and the unit acts nontrivially. This last check only uses an associative algebra, not a Hilbert-space inner product.

If $r\in\operatorname{rad}\sigma$, $\Phi(r)$ commutes with every word. Conversely a central linear generator has $\sigma(r,v)=0$ for every $v$. One may retain these central quantities or prescribe their values in a sector, for example $\Phi(r)=\lambda(r)1$ for a specified real-linear $\lambda$. Imposing every classical condition as an ideal of an unreduced quantum algebra would be different: in a single pair with $[Q,P]=i\hbar 1$, the two-sided ideal generated by $Q$ already contains $1$. Its quotient describes no canonical pair at all.

\section{Bounded Weyl generators and the meaning of regularity}
A \emph{representation} assigns to each algebra element $A$ an operator $\pi(A)$, linearly, with $\pi(1)$ the identity and $\pi(AB)=\pi(A)\pi(B)$. A Hilbert-space representation of the polynomial algebra generally needs \emph{unbounded} operators. Specify a common dense domain $\mathcal D_\pi$ carried into itself by all the operators, so every stated product acts there. To represent the involution we require $\langle\xi,\pi(A)\eta\rangle=\langle\pi(A^*)\xi,\eta\rangle$ on this domain. A Hermitian symbol acts symmetrically, not necessarily self-adjointly after closure. The Weyl construction avoids these domain questions in the algebra multiplication itself. A $C^*$-algebra is a complete normed $*$-algebra with submultiplicative norm and $\|A^*A\|=\|A\|^2$; bounded operators on a Hilbert space are the model \cite{BR}.

\begin{definition}[Weyl relations]\label{f:weyl-relations}
Use one unitary symbol $W(v)$ for each $v\in V$, with
\begin{equation}
 W(0)=1,\qquad W(v)^*=W(-v),\qquad
 W(v)W(w)=e^{-i\hbar\sigma(v,w)/2}W(v+w).
 \label{f:weyl-phase}
\end{equation}
The Weyl algebra $\mathcal W(V,\sigma)$ is the universal unital $C^*$-algebra with these relations. The exponent convention is $W(tv)=\exp(it\Phi(v))$ when a field operator exists; there is no further factor of $1/\hbar$ in this convention.
\end{definition}
To construct it, first take finite formal sums of the symbols. The two parenthesizations of three factors have the same phase because $\sigma(u,v)+\sigma(u+v,w)=\sigma(v,w)+\sigma(u,v+w)$. Reversing the product conjugates its phase, as required by $*$. Take the supremum of the operator norm over Hilbert-space realizations by unitaries: the norm of $\sum c_j W(v_j)$ is at most $\sum |c_j|$. It separates elements. Indeed on $\ell^2(V)$ with orthonormal vectors $e_w$ indexed by the labels, put
\[
 U(v)e_w=e^{-i\hbar\sigma(v,w)/2}e_{v+w}.
\]
These operators permute a basis up to phases, obey \eqref{f:weyl-phase}, and a finite sum with distinct labels applied to $e_0$ records every coefficient. Completion in the supremum norm gives the $C^*$-algebra. Each unitary realization extends by continuity, uniquely since the finite span is dense. Also $W(v)W(tw)=e^{-i\hbar t\sigma(v,w)}W(tw)W(v)$. Commutation for all real $t$ and all $w$ is equivalent to $v\in\operatorname{rad}\sigma$: in particular nondegeneracy excludes nonzero central Weyl labels, not the scalar unit.

For example, to remove the radical $R$ in its zero-charge sector, divide by the closed ideal generated by $W(r)-1$, $r\in R$. In this quotient $W(v+r)=W(v)W(r)$ has the same class as $W(v)$. These classes obey the Weyl relations of $(V/R,\bar\sigma)$; conversely sending $W(v)$ to the generator labeled by $v+R$ kills the ideal. The two maps are inverse on generators. Retaining a central label, fixing a character on it, and discarding the label without specifying a sector are different operations.

A representation is \emph{regular} if $t\mapsto\pi(W(tv))$ is strongly continuous for each $v$, that is, continuous on each Hilbert vector in norm. Stone's theorem then gives a self-adjoint generator $\widehat\Phi(v)$ for this one-parameter group. One may differentiate the Weyl relation in both parameters on common smooth vectors for a finite collection of directions, obtaining $[\widehat\Phi(v),\widehat\Phi(w)]=i\hbar\sigma(v,w)1$. Such vectors are dense: for the finite-dimensional span $L$, $(t,v)\mapsto e^{it}\pi(W(v))$ is a continuous unitary representation of $\mathbb R\times L$ with group law $(t,v)(s,w)=(t+s-\hbar\sigma(v,w)/2,v+w)$. Averaging a vector against smooth, compactly supported approximate identities on this group gives smooth approximants. Differentiating is not an operation in the abstract $C^*$-norm. The realization $U$ above is not even regular when $V\ne0$: for $v\ne0$ and $t\ne0$, $\|U(tv)e_0-e_0\|=\sqrt2$. Nor do symmetric polynomial generators satisfying a commutator on a domain, without self-adjointness and control of exponentiation, give the Weyl relations. The two universal constructions answer different questions.

\section{Positivity constructs a cyclic representation}
A \emph{state} on a unital $*$-algebra $\mathcal A$ is a complex-linear functional $\rho$ with $\rho(1)=1$ and $\rho(A^*A)\ge0$ for every $A$. Positivity is stronger than assigning real numbers to Hermitian elements: it constrains measurements of sums of observables. In particular $(A,B)_\rho=\rho(A^*B)$ is a positive semidefinite Hermitian form (conjugate-linear in the first entry). Polarizing its real nonnegative diagonal gives Hermitian symmetry; positivity on $A+zB$ for all complex $z$ gives $|(A,B)_\rho|^2\le(A,A)_\rho(B,B)_\rho$. Consequently a vector of zero length pairs to zero with every vector.

\begin{theorem}[Gelfand--Naimark--Segal construction]\label{f:gns}
Every state is the vector state of a cyclic representation. On the domains generated by the cyclic vectors this realization is unique up to a unitary taking one such vector to the other. For a general $*$-algebra these are common invariant dense domains; for a $C^*$-algebra the representation extends by bounded operators to the completed Hilbert space.
\end{theorem}
\begin{proof}
Set $N_\rho=\{A:(A,A)_\rho=0\}$. By the inequality this is the subspace orthogonal to everything. If $A\in N_\rho$ then $\rho((BA)^*BA)=(A,B^*BA)_\rho=0$ for every $B$, so it is a \emph{left} ideal. On the quotient vector space define
\[
 \langle[A],[B]\rangle=\rho(A^*B),\qquad
 \pi_\rho(B)[A]=[BA],\qquad \Omega_\rho=[1].
\]
The left ideal is exactly what makes left multiplication well defined here; we do not multiply two classes as elements of a quotient algebra. The inner product is positive definite, and the adjoint pairing holds since $(BA)^*C=A^*B^*C$. Complete to $\mathcal H_\rho$. The vectors $\pi_\rho(A)\Omega_\rho=[A]$ form the dense domain (this is cyclicity) and $\langle\Omega_\rho,\pi_\rho(A)\Omega_\rho\rangle=\rho(A)$. In the $C^*$ case, $\|B\|^2 1-B^*B$ has a positive square root, so positivity on its sandwich between $A^*$ and $A$ gives $\|[BA]\|^2\le\|B\|^2\|[A]\|^2$. Thus each operator extends to the completion. Finally, in any cyclic realization of the same state, $[A]\mapsto\pi(A)\Omega$ preserves these inner products and has dense range. Its unitary extension intertwines on the generated domain, and everywhere for bounded operators.
\end{proof}

A real-linear bijection $S:V_1\to V_2$ preserving the forms gives a $*$-isomorphism $\alpha_S$ by $W_1(v)\mapsto W_2(Sv)$; apply universality in both directions. The same argument applies to polynomial generators. A state $\rho_2$ pulls back to $\rho_2\circ\alpha_S$. For two \emph{specified} states, a unitary intertwiner taking one cyclic vector to the other exists precisely when $\rho_1=\rho_2\circ\alpha_S$, by the inner-product argument above. Another question is whether an automorphism is implemented in a \emph{fixed} representation by $\pi(\alpha_S(A))=U\pi(A)U^*$. Such a unitary need not fix the state vector; it need not exist at all (Exercise~\ref{f:ex-implementation}).

\section{Characteristic functions and mixed covariances}
A Weyl state is determined by its characteristic function $v\mapsto\rho(W(v))$, since finite sums are dense. For a real symmetric bilinear form $\mu$, consider the centered Gaussian candidate $\chi(v)=\exp[-\mu(v,v)/2]$. Which $\mu$ define states? The answer includes a restriction on the correlations between different labels, not just on each variance.

\begin{proposition}[Gaussian positivity]\label{f:gaussian-positivity}
The candidate is a state characteristic function if and only if $\Lambda=\mu+i\hbar\sigma/2$, extended sesquilinearly to $V_{\mathbb C}$, is positive semidefinite. The Gaussian representation is regular, with ordered two-point function $\langle\Omega_\rho,\widehat\Phi(v)\widehat\Phi(w)\Omega_\rho\rangle=\Lambda(v,w)$.
\end{proposition}
\begin{proof}
Positivity on finite Weyl sums is exactly positive semidefiniteness of the matrix $K_{jk}=\rho(W(v_j)^*W(v_k))$ for every finite list. For the candidate,
\[
 K_{jk}=e^{-\mu(v_j,v_j)/2}\,e^{\Lambda(v_j,v_k)}\,e^{-\mu(v_k,v_k)/2} .
\]
If $[\Lambda(v_j,v_k)]$ is a Gram matrix of vectors, its entrywise $n$th power is a Gram matrix of their $n$-fold tensor powers. Summing with weights $1/n!$, including the all-ones matrix at $n=0$, proves positivity of the entrywise exponential. The outside diagonal factors preserve it. The resulting state on finite sums extends: its cyclic representation has unitary Weyl generators, hence is bounded by the defining supremum norm. Conversely the Gram matrices of $(W(tv_j)-1)/t$ in the state are positive. Their entries tend to $\Lambda(v_j,v_k)$: expand $K_{jk}$ at labels $tv_j,tv_k$, subtract $\chi(-tv_j)$ and $\chi(tv_k)$, and add $1$ before dividing by $t^2$. Limits preserve positivity.

On any finite-dimensional label span, the kernel $\langle\pi(W(x))\Omega_\rho,\pi(W(y))\Omega_\rho\rangle$ is smooth in $(x,y)$ by the displayed formula. The vector orbit of $\Omega_\rho$ is smooth there as well: the inner products of two difference quotients are mixed difference quotients of this kernel, so those vector quotients are Cauchy in norm. The same argument applies to iterated differences. Weyl translates of $\Omega_\rho$ have the same property and span a dense subspace; boundedness of the unitaries gives strong continuity everywhere. The ordered product in the proposition is thus taken on a smooth vector, not in the abstract $C^*$-algebra. Applying $-\partial_s\partial_t$ to $e^{-i\hbar st\sigma(v,w)/2}\chi(sv+tw)$ at zero gives its expectation.
\end{proof}
In particular $\mu(v,v)\mu(w,w)\ge\mu(v,w)^2+\hbar^2\sigma(v,w)^2/4$. The commutator fixes the imaginary part of the ordered moment but not the real covariance $\mu(v,w)$ \cite{WaldQFT}.

\begin{corollary}[Ordered Wick pairings]\label{f:wick-pairings}
In the centered Gaussian state above, also called a centered \emph{quasifree} state, take real labels $v_1,\ldots,v_n$. Write $X_j=\widehat\Phi(v_j)$ and $\Lambda_{jk}=\Lambda(v_j,v_k)$. The cyclic-vector expectation in the specified operator order is
\[
 \langle X_1\cdots X_n\rangle_\rho=
 \begin{cases}
  0, & n\text{ odd},\\[2pt]
  \displaystyle\sum_{\mathcal P}\prod_{(j,k)\in\mathcal P}\Lambda_{jk}, & n\text{ even}.
 \end{cases}
\]
Here $\mathcal P$ runs over partitions into pairs, writing the earlier index first in each pair. No time ordering is understood.
\end{corollary}
\begin{proof}
Put $u=\sum_j t_jv_j$ with real parameters. Multiplying the Weyl operators in the written order, each new factor acquires a phase from every earlier factor. Hence their generating function is
\begin{align*}
 Z(t)&=\rho\big(W(t_1v_1)\cdots W(t_nv_n)\big)
 =e^{-(i\hbar/2)\sum_{j<k}t_jt_k\sigma(v_j,v_k)}\chi(u)\\
 &=\exp\!\left[-\tfrac12\sum_j t_j^2\mu(v_j,v_j)
                       -\sum_{j<k}t_jt_k\Lambda_{jk}\right].
\end{align*}
On the smooth cyclic vector of Proposition~\ref{f:gaussian-positivity} the ordered moment is $i^{-n}\partial_{t_1}\cdots\partial_{t_n}Z(0)$. The diagonal terms cannot contribute to the coefficient of $t_1\cdots t_n$. Every contributing term selects disjoint pairs covering all the indices, so $n=2m$. The $m!$ possible orders of the selected terms cancel the exponential's $1/m!$; their sign $(-1)^m$ cancels $i^{-2m}$. This proves the formula. Regularity alone gives field generators, not moments of every order in an arbitrary state vector; even when those moments exist, they need not obey this Gaussian rule.
\end{proof}

For $X=\widehat\Phi(v)$ and $Y=\widehat\Phi(w)$ the four-field example is
\[
 \langle X^2Y^2\rangle_\rho
 =\Lambda(v,v)\Lambda(w,w)+2\Lambda(v,w)^2 .
\]
Subtracting the mean from each square leaves two cross pairings, both with $X$ before $Y$: the last factor is a square, not an absolute square. More generally $N(z)=e^{\mu(z,z)/2}\pi(W(z))$ generates \emph{Wick powers} by $i^{-k}\partial_s^k N(sz)|_{s=0}$ on smooth vectors. Removing the contractions internal to each factor leaves the cross term in the identity
\[
 N(sv)N(tw)=e^{-st\Lambda(v,w)}N(sv+tw),
\]
obtained directly from the Weyl relation. In differentiating three times in each variable, two cross contractions choose two out of three slots on each side and a bijection between those slots: $\binom32^2\,2!=18$. The remaining derivatives leave one field from each side in a single Wick product. This is the multiplicity in the fish term \eqref{f:ext-fish}; time ordering and extension to coincident points are separate operations there. Here all scalar field labels are smeared tests modulo the equation of motion.

\begin{example}[An oscillator in position space]\label{f:oscillator-state}
For $v=(a,b)$, $w=(c,d)$ let $\sigma(v,w)=ad-bc$. On $L^2(\mathbb R,\dd x)$ take $Q\psi=x\psi$, $P\psi=-i\hbar\psi'$ on Schwartz functions, so that $\widehat\Phi(v)=aQ+bP$ there. Its Weyl group acts by
\[
 (W(a,b)\psi)(x)=e^{ia(x+\hbar b/2)}\psi(x+\hbar b).
\]
Composing two such translations gives the phase difference $\hbar(bc-ad)/2$ relative to $W(a+c,b+d)$, just as in \eqref{f:weyl-phase}. Take the normalized oscillator wave function $\psi_0(x)=(\pi l^2)^{-1/4}e^{-x^2/(2l^2)}$, $l>0$. In the expectation of $W(a,b)$ the product of Gaussians is centered at $x=-\hbar b/2$. Shifting to $y=x+\hbar b/2$ leaves the Fourier integral of $e^{-y^2/l^2}$, giving
\[
 \langle\psi_0,W(a,b)\psi_0\rangle
   =\exp\!\left[-\tfrac14(l^2 a^2+\hbar^2 b^2/l^2)\right].
\]
The two variances are $l^2/2$ and $\hbar^2/(2l^2)$; their product saturates the bound. Neither the width nor the vector was specified by the Weyl relations.
\end{example}

Two copies already distinguish symplectic and statistical independence. Set $\hbar=l=1$ here. Translate the first oscillator wave function by a random displacement $X$ of mean zero and Gaussian variance $s^2$, and the second by $\varepsilon X$ with $\varepsilon=+1$ or $-1$. Average these product vector states; convex averaging preserves positivity. Translation of a wave function to center $X$ multiplies its characteristic function by $e^{iaX}$. The joint characteristic functions are therefore
\[
 \chi_\varepsilon=\exp\!\left[-\tfrac14\sum_{j=1}^2(a_j^2+b_j^2)
                          -\tfrac{s^2}{2}(a_1+\varepsilon a_2)^2\right].
\]
Setting either pair of labels to zero gives identical marginals for the two signs. Yet $\rho(Q_1Q_2)=\varepsilon s^2$ while $[Q_1,Q_2]=0$. These correlations need not even be entanglement: the mixture was constructed from product vectors. For two general label subspaces, take finite lists from each. If $C$ is the trial real cross-covariance matrix, Gaussian positivity requires
\[
 \begin{pmatrix}
  \mu_1+i\hbar\sigma_{11}/2 & C+i\hbar\sigma_{12}/2\\[2pt]
  (C+i\hbar\sigma_{12}/2)^* & \mu_2+i\hbar\sigma_{22}/2
 \end{pmatrix}\ \ge 0
\]
for \emph{every} choice of lists. Here $*$ is conjugate transpose, $\mu_j$ are the two marginal blocks, and $\sigma_{jk}$ are fixed by the joint bracket. Any linear combination representing a zero label must pair to zero with both lists. A tensor product of marginal states does not in general solve this extension problem. For a spacetime field, a further question is whether the two-point function is a distribution with admissible singularities (Chapter~\ref{f:microlocal}).

\section{Which generators belong to the whole?}
There is a situation where a tensor product is justified at the algebraic level. If $V=V_1\oplus V_2$ and $\sigma(V_1,V_2)=0$, the Weyl generators of the summands commute and $W(v_1+v_2)=W(v_1)W(v_2)$. Conversely any two Weyl realizations with commuting ranges give a realization of this direct sum by that product. This is the universal property of the maximal $C^*$ tensor product (the completion allowing all commuting representations). It neither requires nor selects a product state. If the mixed bracket is nonzero, operators on separate tensor factors would commute and give the wrong phase. If labels are shared, treating the copies as independent misses their identification. Even when two subalgebras commute, this does not establish a factorization $\mathcal H=\mathcal H_1\otimes\mathcal H_2$ in which each local algebra acts as all bounded operators on its own Hilbert factor \cite{Haag}.

There is also a generation question even before a state is specified. For subspaces $V_1,V_2\subset V$ the Weyl relations generate the labels in $V_1+V_2$ regardless of their mixed bracket. No other single Weyl label is in the generated $C^*$-subalgebra: in the basis representation above, its elements send $e_0$ into the closed span of $e_v$ with $v\in V_1+V_2$, whereas $W(w)e_0=e_w$ is orthogonal to that span if $w$ is outside the sum. Checking the bracket of known generators does not test whether a proposed observable is generated.

For a nonlinear example return to the two arcs of Example~\ref{f:two-rotor}. In the algebra of configuration functions take the finite Laurent polynomials in the two open transports $z_j=e^{i\alpha_j}$. At the shared endpoints the relative frame rotation sends $(z_1,z_2)$ to $(tz_1,t^{-1}z_2)$, $|t|=1$. The monomial $z_1^k z_2^l$ has weight $k-l$; linear independence of the Laurent monomials makes the invariant subalgebra exactly the Laurent polynomials in $w=z_1z_2$. This $U(1)$ Wilson loop gives the multiplication operator $(U\psi)(\alpha)=e^{i\alpha}\psi(\alpha)$ on $L^2(S^1)$. With angular momentum $L=-i\hbar\partial_\alpha$, differentiating the product gives $[L,U]=\hbar U$ on smooth periodic functions, in agreement with $\{w,e\}=iw$. If one retains on each arc only functions already invariant under independent changes of its endpoint frames, each arc algebra contains just constants within this configuration class. Their products cannot distinguish $w=1$ from $w=-1$. The lost observable is recovered by multiplying the \emph{open} transports and then taking invariants, not by changing a state or a commutator.

\section{Local operator algebras and a separating collar}\label{f:local-factors}
A bounded-operator algebra in a representation has another natural closure besides norm closure. For a set $S\subset\mathcal B(\mathcal H)$ its \emph{commutant} $S'$ consists of all bounded operators commuting with every element of $S$. The double commutant theorem says that for a unital $*$-algebra $S$, $S''$ is its closure for weak operator convergence (convergence of every matrix coefficient). A unital $*$-algebra with $\mathfrak M=\mathfrak M''$ is a von Neumann algebra; it is a \emph{factor} when its center $\mathfrak M\cap\mathfrak M'$ consists only of scalars \cite{BR}.

The Hilbert tensor product is the completion of the algebraic tensor product for the inner product obtained by extending $\langle x\otimes y,x'\otimes y'\rangle=\langle x,x'\rangle\langle y,y'\rangle$ sesquilinearly. For example, take $\mathfrak M=M_2(\mathbb C)\otimes 1$ on $\mathbb C^2\otimes\mathbb C^n$. In a two-by-two array of $n$-dimensional blocks, an operator commuting with $E_{11}\otimes1$ is block diagonal. Commuting also with $E_{12}\otimes1$ forces the diagonal blocks to agree. Conversely $1\otimes T$ commutes with the whole first algebra, so $\mathfrak M'=1\otimes M_n(\mathbb C)$ and the center is scalar. For $n=2$ the Bell vector $(e_1\otimes e_1+e_2\otimes e_2)/\sqrt2$ gives probability $1/2$ for finding the first basis state on both sides, not the product $1/4$ of the marginal probabilities. The state is entangled; the tensor factors still exist.

Projections distinguish a failure of factorization that is independent of this correlation test. A projection $p=p^*=p^2$ is \emph{minimal} in $\mathfrak M$ if it is nonzero and has no nonzero proper subprojection in $\mathfrak M$. In the matrix example $E_{11}\otimes1$ is minimal in $\mathfrak M$ even though its range has dimension $n$. A factor with a minimal projection is type~I; in a suitable Hilbert tensor decomposition it acts as $\mathcal B(\mathcal K)\otimes1$. A projection is \emph{finite} if no proper subprojection $q<p$ is equivalent to it, where equivalence means $v^*v=p$, $vv^*=q$ for some partial isometry $v\in\mathfrak M$. A factor is type~III precisely when it has no nonzero finite projection; in particular it has no minimal projection. The absence of minimal projections alone does not characterize type~III.

Consider the vacuum representation $\pi_0$ of the massive real free scalar in four-dimensional Minkowski space. For a nonempty bounded double cone $O=I^+(p)\cap I^-(q)$, let $\mathfrak R(O)$ be the double commutant of the Weyl unitaries $\pi_0(W([f]))$ with $f$ supported in $O$. A structural theorem for this field says that $\mathfrak R(O)$ is a type~III factor \cite{Haag,Yngvason}. It cannot be the full bounded-operator algebra on a Hilbert factor associated to the sharp region: such an algebra contains minimal projections. This conclusion uses the local field and its representation, not just a pair of nonzero vacuum correlations. Restricting the vacuum to the region means evaluating $\langle\Omega,A\Omega\rangle$ for $A\in\mathfrak R(O)$; it requires no partial trace.

If the closure of a double cone $O_1$ lies inside a larger double cone $O_2$, the two boundaries leave a collar. The \emph{split property} of this field \cite{Haag,FewSplit} supplies a type~I factor $\mathfrak N$ with
\[
 \mathfrak R(O_1)\subset\mathfrak N\subset\mathfrak R(O_2).
\]
In a tensor decomposition for $\mathfrak N$, the inner algebra acts on the first factor, while $\mathfrak R(O_2)'\subset\mathfrak N'$ acts on the second. The latter commutant contains observables localized spacelike to $O_2$, by locality. It is the \emph{intermediate} algebra that is type~I, not either double-cone algebra. The split theorem does not supply a limiting type~I factor when the collar is removed.

\section*{Exercises}
\begin{exercise}[Basic: null vectors are not relations]
Take $\mathcal A=M_2(\mathbb C)$, $\rho(A)=\langle e_1,Ae_1\rangle$ for the first coordinate vector. Identify $N_\rho$, exhibit its failure to be two-sided, and identify the cyclic Hilbert space.
\end{exercise}
\noindent\emph{Solution.} $\rho(A^*A)=\|Ae_1\|^2$, so $N_\rho$ consists of matrices with zero first column. For matrix units, $E_{12}\in N_\rho$ but $E_{12}E_{21}=E_{11}\notin N_\rho$. The isometry $[A]\mapsto Ae_1$ is onto $\mathbb C^2$, and the representation is ordinary matrix multiplication on vectors.

\begin{exercise}[Intermediate: why no bounded polynomial pair?]
Suppose bounded operators on a nonzero Hilbert space satisfy $[A,B]=1$. Use $[A,B^n]$ to obtain a contradiction. What does this say about a nonzero scalar CCR?
\end{exercise}
\noindent\emph{Solution.} The product rule gives $[A,B^n]=nB^{n-1}$. No power of $B$ vanishes, by this identity and $B^0=1$. Taking norms gives $n\|B^{n-1}\|\le2\|A\|\|B\|\|B^{n-1}\|$, impossible for all $n$. Rescaling one operator reduces any nonzero scalar commutator to this one. The conclusion concerns the fields, not their Weyl unitaries.

\begin{exercise}[Advanced: an isomorphism without a unitary]\label{f:ex-implementation}
In units $\hbar=1$, take the algebraic direct sum of countably many standard real label planes $(a_j,b_j)$ (each vector has finite support). Use the Gaussian state with $\mu(v,v)=\tfrac12\sum_j(a_j^2+b_j^2)$. Show that $S:(a_j,b_j)\mapsto(ja_j,b_j/j)$, $j\ge1$, induces a Weyl automorphism not unitarily implementable in this state's representation.
\end{exercise}
\noindent\emph{Solution.} The map and its inverse preserve finite support and each plane's alternating form. Take $v_n$ with only $a_n=1/n$ nonzero. Then $\|\pi(W(v_n))\Omega-\Omega\|^2=2(1-e^{-1/(4n^2)})\to0$. Each fixed $W(w)$ commutes with $W(v_n)$ for all sufficiently large $n$ by finite support. Convergence to the identity thus holds on a dense span of $\pi(W(w))\Omega$, and strongly everywhere by the unit norm bound. Unitary conjugation would give the same strong limit for $\pi(W(Sv_n))$, but its squared displacement of $\Omega$ is $2(1-e^{-1/4})$ for every $n$. The algebraic isomorphism has not failed; it is the proposed operator on this Hilbert space that does not exist.
% !TeX root = ../main.tex
\chapter{Singular directions and two-point functions}\label{f:microlocal}

A delta function on a sheet is singular at every point of the sheet. Yet a
probe moving across the sheet can register a perfectly well-defined impulse,
whereas a probe confined to the sheet would be asked to sample that impulse
at zero distance for its entire history. Singular support alone does not
distinguish the two measurements. We need to know in which directions a
localized field carries arbitrarily high frequencies.

\section{A Fourier test with two localizations}
Let $u$ be a distribution on an open subset of $\mathbb R^n$. Localize it
first in position: take a smooth compactly supported cutoff $\chi$ equal to
one on a neighborhood of $x_0$. Then
\[
  \widehat{\chi u}(\xi)=\langle u,\chi(x)e^{-ix\cdot\xi}\rangle
\]
is a smooth function of at most polynomial growth. Indeed the test function
in this pairing has fixed compact support, and differentiating it a finite
number of times in $x$ costs at most a polynomial in $\xi$. If $u$ is
smooth on the support of $\chi$, integration by parts gives rapid decay in
all directions. Conversely, rapid decay in all directions allows Fourier
inversion after multiplication by every monomial $(i\xi)^\alpha$; hence
$\chi u$ is smooth. The second localization keeps only directions in an
open cone $\Gamma$: if $\xi\in\Gamma$, so is $a\xi$ for every $a>0$.
One ray by itself is not an open cone in dimension greater than one.

\begin{definition}\label{f:wf-test}
A pair $(x_0,\xi_0)$, $\xi_0\ne0$, is \emph{regular for $u$} if there
exist such a cutoff $\chi$ and an open cone $\Gamma$ containing $\xi_0$
for which, for every nonnegative integer $N$,
\begin{equation}
 \sup_{\xi\in\Gamma}(1+|\xi|)^N|\widehat{\chi u}(\xi)|<\infty .
 \label{f:wf-decay}
\end{equation}
The complementary pairs form the \emph{wavefront set} $\operatorname{WF}(u)$.
It is closed in the cotangent bundle with the zero section removed, and
is invariant under positive rescaling of the covector.
\end{definition}

The same cutoff and cone work for nearby pairs, which explains closedness.
A more useful fact is that one can shrink the position neighborhood even
when this changes the Fourier transform everywhere. If $\psi$ is supported
where $\chi=1$, the transform of $\psi u$ is the convolution of
$\widehat\psi$ with $\widehat{\chi u}$, divided by $(2\pi)^n$.
Take a smaller cone whose unit directions have closure inside $\Gamma$.
For $\xi$ in this cone, split the convolution integral over $\eta$ into
$\eta\in\Gamma$ and its complement. In the first part use
$1+|\xi|\le(1+|\xi-\eta|)(1+|\eta|)$: every power of the second factor
can be absorbed by \eqref{f:wf-decay}, and $\widehat\psi$ is rapidly
decreasing. In the second part angular separation gives
$|\xi-\eta|\ge c(|\xi|+|\eta|)$ for some $c>0$. There the decay of
$\widehat\psi(\xi-\eta)$ absorbs both the polynomial growth of
$\widehat{\chi u}(\eta)$ and any desired power of $|\xi|$.

\begin{proposition}
The projection of $\operatorname{WF}(u)$ onto the base is the ordinary
singular support of $u$.
\end{proposition}
\begin{proof}
A smooth neighborhood supplies a single cutoff good in all directions.
Conversely, if no direction over $x_0$ is singular, cover the unit sphere
by finitely many of the \emph{smaller} cones just described. Choose one
cutoff equal to one near $x_0$ and supported where all the corresponding
position cutoffs are one. Its transform is rapidly decreasing on each of
the finitely many cones, hence everywhere. Fourier inversion now gives a
smooth representative near $x_0$.
\end{proof}

\section{A sheet and a one-sided singularity}
In coordinates $(y,s)\in\mathbb R^{n-1}\times\mathbb R$, the distribution
$\delta(s)$ acts by integrating a test function over $s=0$. With dual
coordinates $(\eta,\tau)$ one has
\[
 \widehat{\chi\delta(s)}(\eta,\tau)
    =\int e^{-iy\cdot\eta}\chi(y,0)\,\mathrm{d}y .
\]
In a cone with $|\eta|\ge c|(\eta,\tau)|$, rapid decay follows by
integrating in $y$. In contrast, when $\chi$ equals one near a point of
the sheet its restriction to the sheet is not identically zero. Its
transform is nonzero at some fixed $\eta_*$; the sequence
$(\eta_*,\tau)$ tends in direction to either normal ray as
$\tau\to\pm\infty$, without any decay. Thus the wavefront set consists
precisely of $(y,0;0,\tau)$ with $\tau\ne0$. The two signs matter.

A distribution can have just one sign even at a point singularity.
Set $b_\pm(x)=(x\pm i0)^{-1}$, the limits in distributions of
$(x\pm i\epsilon)^{-1}$ as $\epsilon\downarrow0$. Their real part tends
to the principal value of $1/x$; their imaginary parts tend to
$\mp\pi\delta_0$, since $\epsilon/(x^2+\epsilon^2)$ has integral $\pi$
and concentrates at zero. To find the singular directions, use the actual
Fourier sign, not the sign of the imaginary part. For instance
\[
  (x+i\epsilon)^{-1}=-i\int_0^\infty e^{ipx-\epsilon p}\,\mathrm{d}p.
\]
The conjugate formula has $+i$ and $e^{-ipx}$. With forward phase
$e^{-ix\xi}$ and inverse factor $1/(2\pi)$ this gives
\begin{equation}
 \widehat b_+(\xi)=-2\pi i\,\mathbf 1_{\xi>0},\qquad
 \widehat b_-(\xi)= 2\pi i\,\mathbf 1_{\xi<0} .
 \label{f:wf-bv}
\end{equation}
These are transforms of the \emph{unlocalized} distributions. To perform
the defining test, multiply by $\chi$ and convolve. The answer for $b_+$
is $-i\int_{-\infty}^\xi\widehat\chi(v)\,\mathrm{d}v$. It decreases
rapidly toward $-\infty$ but tends to $-2\pi i\chi(0)$ toward $+\infty$.
The answer for $b_-$ is $i\int_\xi^\infty\widehat\chi(v)\,\mathrm{d}v$.
Consequently $\operatorname{WF}(b_\pm)=\{(0,\xi):\ \pm\xi>0\}$. Although
both singular supports are $\{0\}$, the singular frequencies are disjoint.

\begin{exercise}
Determine the two limiting high-frequency values of
$\widehat{\chi\,\mathrm{pv}(1/x)}$ when $\chi=1$ near zero. Do the same
for $\chi\delta_0$. Why is the singularity of neither distribution
one-sided?
\end{exercise}
\noindent\emph{Hint.} Take the half-sum of $b_+$ and $b_-$ in the first
case and $(b_--b_+)/(2\pi i)$ in the second. On the positive ray the
answers tend to $-i\pi$ and $1$, respectively; now reverse the ray.

\section{Restriction and multiplication as pullbacks}
For a smooth map $F:Y\to X$, the transpose of the differential sends a
covector $k$ at $F(y)$ to $dF_y^T k$ at $y$. A rapid phase on $X$ may
cease to oscillate after composition with $F$: its new differential can
be zero. Define the set of offending directions of the map by
\[
 N_F=\{(F(y),k):k\ne0,\ dF_y^Tk=0\}.
\]
For an embedding this is its nonzero conormal bundle: its covectors
annihilate the tangent space of the embedded submanifold.

\begin{theorem}[H\"ormander's pullback theorem, smooth manifolds \cite{HormI}]
\label{f:wf-pullback}
If $u$ is a distribution on $X$ and $N_F\cap\operatorname{WF}(u)$ is
empty, the distributional pullback $F^*u$ exists canonically, agrees
with composition on smooth functions, and satisfies
\[
 \operatorname{WF}(F^*u)\subset
 \{(y,dF_y^Tk):(F(y),k)\in\operatorname{WF}(u)\}.
\]
Here distributions are scalar, acting on test densities. For a
diffeomorphism the inclusion is equality, by applying the theorem to
the inverse.
\end{theorem}

The transformation law also explains why these directions are covectors
rather than vectors. In coordinates the oscillatory factor before Fourier
transformation on $Y$ is $e^{i F(y)\cdot k}$; its new Fourier phase has
differential $dF_y^Tk-\eta$. Directions separated from the possible
stationary ones can be estimated by integration by parts. In particular,
straightening a smooth sheet by a diffeomorphism turns the calculation for
$\delta(s)$ into the intrinsic statement: a nonvanishing smooth density
of delta mass on a sheet has exactly its nonzero conormals as wavefront
set. Smooth nonzero density factors do not change this conclusion.

For multiplication use the diagonal map $d:x\mapsto(x,x)$ and set
$uv=d^*(u\otimes v)$. Its transpose sends $(k,l)$ to $k+l$.
The tensor-product wavefront estimate allows a singular covector from
either factor alone (zero in the other slot), or one from each. Thus a
sufficient test is that no $(x,k)\in\operatorname{WF}(u)$ have
$(x,-k)\in\operatorname{WF}(v)$. In particular $b_+^2$ passes, whereas
$b_+b_-$ does not. The latter's common-width regularizations give
$1/(x^2+\epsilon^2)$, with pairing asymptotic to
$\pi\psi(0)/\epsilon$ on a test $\psi$ with $\psi(0)\ne0$. A failed
sufficient test is not a general impossibility theorem: the square of a
bounded step function, for example, is defined by ordinary almost
everywhere multiplication.

For a kernel $K(x,y)$ it is useful to record $(x,k;y,l)$ in its
\emph{primed} relation when $(x,k;y,-l)\in\operatorname{WF}(K)$. Then the
delta kernel of the identity relates a covector to itself. There is a
simple composition rule when every singular pair of both kernels has two
nonzero entries. For $K(x,y),L(y,z)$ assume also that the projection of
their joint support in $(x,y,z)$ onto $(x,z)$ is proper (compact sets
have compact inverse images there). Multiply their lifts to the triple
space and integrate in $y$ with the chosen smooth density. Cancellation in
the product test cannot occur: a singularity from $K$ has nonzero $x$
covector, one from $L$ has nonzero $z$ covector. Under integration only
singular directions with zero \emph{total} $y$ covector can remain. The
primed relation of the result is therefore contained in the composition
of the two primed relations \cite{HormI}. Dropping the support assumption
can instead leave a divergent integral even for smooth kernels.

\begin{exercise}
Pull the sheet $\delta(s)$ in the $(y,s)$ plane back by the curves
$v\mapsto(v,av)$ for $a\ne0$ and $v\mapsto(v,v^2)$. Where does the test
fail? For the second curve, try a smooth approximation
$\delta_\epsilon(s)=\epsilon^{-1}\rho(s/\epsilon)$ with $\rho\ge0$,
$\rho(0)>0$, compact support and unit integral.
\end{exercise}
\noindent\emph{Solution.} The transverse pullback is
$|a|^{-1}\delta(v)$. At the tangency the conormal $\mathrm{d}s$ pulls
back to zero. Setting $v=\sqrt\epsilon\,w$ in the regularized pairing
gives leading term $\epsilon^{-1/2}\psi(0)\int\rho(w^2)\,\mathrm{d}w$.
The transverse formula has no limiting analogue here.

\section{Which singularities can a wave transport?}
Now take $P=\Box_g-m^2-\xi R$ for a real scalar on a smooth Lorentzian
spacetime, with $\Box_g=\nabla^a\nabla_a$, $R$ the scalar curvature and
$\xi$ a constant curvature coupling. Its characteristic covectors are the
nonzero zeros of
$p(x,k)=g^{ab}k_a k_b/2$. With Fourier symbol $\partial_a\mapsto ik_a$,
$p$ is the principal symbol multiplied by $-1/2$; this rescaling only
reparametrizes its Hamilton curves. A hypersurface is
\emph{non-characteristic for $P$} if its nonzero conormals are not
characteristic. This is a property of an operator and a surface, not a
test for restricting an arbitrary distribution to the surface. For
example, $t=0$ in Minkowski space is non-characteristic for $P$, yet
$\delta(t)$ has precisely the conormals that obstruct its own trace.

\begin{theorem}[Microlocal elliptic regularity and propagation, scalar wave operator \cite{HormIII}]
\label{f:wf-propagation}
Let $u$ be a distribution in a smooth boundaryless open region, and
let $O$ be an open conic set in its nonzero cotangent bundle on which
$Pu$ is microlocally smooth. Then $\operatorname{WF}(u)\cap O$ lies in
$p=0$ and is a union of maximal segments in $O$ of the Hamilton curves
of $p$ there.
\end{theorem}

The hypothesis on $Pu$ says $O\cap\operatorname{WF}(Pu)=\varnothing$.
The Hamilton curves in the characteristic set are called
\emph{bicharacteristics}. The real principal type hypothesis of the
propagation theorem holds here: $\partial_k p=k^\sharp\ne0$ on the
nonzero null cone. Off that cone the leading symbol is invertible, giving
the elliptic regularity assertion. On the cone the Hamilton equations are
\[
 \dot x^a=g^{ab}k_b,\qquad
 \dot k_a=-\tfrac12(\partial_a g^{bc})k_b k_c
          =\tfrac12(\partial_a g_{bc})\dot x^b\dot x^c .
\]
The last expression equals $\Gamma^c{}_{ab}k_c\dot x^b$ for the
Levi--Civita connection. Thus $k$ is parallel along $x$ and $x$ is a
null geodesic. The theorem transports a singularity \emph{if it is
present}; it neither puts one on every ray nor chooses its frequency
sign. At a source singularity the hypothesis on $Pu$ must be rechecked.

In particular, if $Pu$ is smooth, every non-characteristic hypersurface
passes the pullback test for $u$; so do its spacetime derivatives, which
cannot enlarge the wavefront set. Both spacelike and timelike
hypersurfaces are non-characteristic for this operator. This does not
make a timelike surface a Cauchy surface. Conversely a perfectly smooth
solution has a restriction even to a characteristic surface. These are
different questions from well-posedness for data on that surface.

\begin{exercise}
In two-dimensional Minkowski space, take $u(t,x)=b_-(t-x)$. Check that
$\Box u=0$ distributionally, determine its wavefront set, and follow one
singular covector along a Hamilton curve.
\end{exercise}
\noindent\emph{Hint.} Use the coordinates $(t-x,x)$ to identify the
covectors as $a\,\mathrm{d}(t-x)$ with $a<0$ over $t=x$. The two second
derivatives in $-\partial_t^2+\partial_x^2$ cancel. The metric dual of
that covector points along the line, not across it.

\section{The frequency selected by a Hadamard state}
Return to a smooth time-oriented globally hyperbolic spacetime without
boundary. For real test functions $f,h$, let $W(f,h)=\Lambda([f],[h])$
be the ordered moment of Proposition~\ref{f:gaussian-positivity}, and
extend $W$ bilinearly to complex tests. Assume continuity on spacetime
tests so the Schwartz kernel theorem gives a distribution $W$ on the
product spacetime \cite{HormI}. This $W$ is not a Weyl generator.
With $\hbar=1$ the free equation and CCR read
$P_x W=P_{x'} W=0$ and $W-W^T=iE$, with $E=G_R-G_A$ and
$W^T(x,x')=W(x',x)$. Hermiticity and positivity mean that $W(\bar f,h)$
is a Hermitian form with $W(\bar f,f)\ge0$, using the spacetime volume
density. Which of the two frequency
orientations is its short-distance singularity allowed to carry?

With signature $(-,+,\ldots,+)$ call a causal covector $k$ future if
$k^\sharp$ is future, equivalently $k(v)<0$ on future timelike vectors
$v$ for nonzero causal $k$. A flat positive-energy phase
$e^{-i\omega t+i\mathbf p\cdot\mathbf x}$ has $k_t=-\omega$, not
$+\omega$, for our Fourier transform. Individual plane waves are smooth;
it is the superposition up to unbounded momentum that produces the
singularity. In a massive field the large-momentum directions still
approach the null cone, in agreement with the principal-symbol test.

\begin{theorem}[Microlocal characterization of Hadamard states \cite{Rad,SV}]
\label{f:wf-hadamard}
For a quasifree Klein--Gordon state with distributional two-point function
on a smooth time-oriented globally hyperbolic spacetime without boundary,
the Hadamard condition in four dimensions \cite{Rad}, and more generally
in spacetime dimension at least three \cite{SV}, is equivalently
\begin{equation}
 \operatorname{WF}(W)=C^+:=
 \{(x,k;x',-k'):(x,k)\sim(x',k'),\ k^\sharp\text{ future null}\}.
 \label{f:wf-spectrum}
\end{equation}
The relation means cotangents to the same null geodesic, parallel
transported along it, with the null tangent proportional to their metric
duals. It includes $x=x'$ with $k=k'\ne0$ null. It does not require
$x$ to be later than $x'$.
\end{theorem}

In four spacetime dimensions one can see exactly what is being prescribed
in position space. Work in a convex normal neighborhood, taking the pair
in a smaller causally convex neighborhood so no returning null ray leaves
the chart. Let $\gamma$ be the unique local geodesic from $x'$ to $x$,
affinely parametrized from zero to one. Write $\sigma=\sigma_g(x,x')$ for
Synge's world function \eqref{f:world-function}: it is half the signed
squared interval, $\tfrac12\int_0^1 g(\dot\gamma,\dot\gamma)\,\mathrm{d}s$.
For a smooth temporal function $T$ increasing to the future, in length units,
put $\sigma_\epsilon=\sigma+i\epsilon(T(x)-T(x'))+\epsilon^2/2$.
The local singular part for a canonically normalized scalar is
\begin{equation}
 H_\ell=\frac1{8\pi^2}\lim_{\epsilon\downarrow0}
 \left(\frac{U}{\sigma_\epsilon}
       +V\log\frac{\sigma_\epsilon}{\ell^2}\right),\qquad \ell>0 ,
 \label{f:wf-parametrix}
\end{equation}
with the logarithm real on the positive axis and $\ell$ a length. The
coefficients are fixed geometrically up to the smooth freedom of a
parametrix. For instance, off the cone use
$\nabla_a\sigma\nabla^a\sigma=2\sigma$ to calculate
\[
 P(U/\sigma)=\frac{PU}{\sigma}
  +\frac{(4-\Box_g\sigma)U-2\nabla^a\sigma\nabla_a U}{\sigma^2} .
\]
Canceling the second numerator transports $U$ along the geodesic, from
$U(x',x')=1$; the logarithmic term cancels the subsequent singular orders.
On smooth metrics this is a parametrix modulo smooth terms, not a claim
that a series for $V$ converges. At any specified finite differentiability
order one can use sufficiently many transport coefficients
\cite{WaldQFT,HWick}. In the massless flat vacuum $U=1,V=0$. Write
$\tau=t-t'$, $\mathbf z=\mathbf x-\mathbf x'$ and $r=|\mathbf z|$.
For $p=|\mathbf p|$, the angular average of $e^{i\mathbf p\cdot\mathbf z}$
is $\sin(pr)/(pr)$. Thus the damped positive-frequency integral is
\begin{align*}
 \int \frac{\mathrm{d}^3\mathbf p}{(2\pi)^3\,2p}
       e^{-(\epsilon+i\tau)p+i\mathbf p\cdot\mathbf z}
 &=\frac1{4\pi^2 r}\int_0^\infty
         e^{-(\epsilon+i\tau)p}\sin(rp)\,\mathrm{d}p\\
 &=\frac1{4\pi^2[r^2-(\tau-i\epsilon)^2]} .
\end{align*}
Taking $\epsilon\downarrow0$ checks the factor of two from the world
function. It also recovers the Fourier test: near a noncoincident null
pair use the normal coordinate $q=r^2-\tau^2$. For $\tau>0$ the boundary
value is $b_+(q)$ up to the prefactor, and for $\tau<0$ it is $b_-(q)$.
In either case a singular covector $a\,\mathrm{d}_x q$ in the first
argument has $k_t=-2a\tau<0$ by \eqref{f:wf-bv}.
The local Hadamard requirement is that $W-H_\ell$ be smooth, using the
full parametrix modulo smooth terms. A change of $\ell$ adds a smooth
multiple of $V$, not a new frequency direction. $H_\ell$ alone
need not be a positive bisolution. A true boundary may introduce waves
reflected along additional rays; the bulk theorem does not supply a
global wavefront condition for an unspecified boundary problem.

\Needspace{.47\textheight}
Two consequences separate measurement from subtraction. First, if $W_1$
and $W_2$ obey \eqref{f:wf-spectrum} and have the same commutator, their
difference is symmetric and has wavefront set contained in $C^+$. Its
transpose instead has past first covector. Since these cones are disjoint,
the difference is a smooth symmetric bisolution. Smooth does not mean
zero: adding $c\,v(x)v(x')$ with $c\ge0$ and $v$ a smooth real solution
preserves positivity and the commutator and can change a fluctuation.
Second, restricting both arguments of $W$ to a spacelike hypersurface is
permitted: its nonzero conormals are timelike and cannot be either null
entry in \eqref{f:wf-spectrum}. On a null hypersurface two arguments on
the same generator can instead have a wavefront pair normal in both
slots. Restricting to $x=x'$ is yet another
map: $(k,-k)$ is a diagonal conormal, so the ordinary pullback test
fails. By contrast $W^2$ as a distribution of \emph{two} arguments passes
the product test because the first covectors cannot be opposite.

\begin{exercise}
Restrict both arguments of the four-dimensional flat massless vacuum to
one inertial clock $\gamma(t)=(t,\mathbf0)$. Why is the pullback allowed
even at equal clock readings? As a distribution of $\tau=t-t'$, compute
the Fourier transform of the result and check its spectral sign.
\end{exercise}
\noindent\emph{Solution.} A nonzero null covector cannot annihilate the
clock's timelike tangent. The restricted distribution is
$-(4\pi^2)^{-1}(\tau-i0)^{-2}=(4\pi^2)^{-1}\partial_\tau b_-$, as one
also finds by doing the exponentially damped mode integral at $r=0$.
Equation~\eqref{f:wf-bv} gives
$-\xi\,\mathbf 1_{\xi<0}/(2\pi)$, a nonnegative spectral density on the
negative temporal Fourier ray. This is a distribution on two clock
readings, not a finite pointwise variance at one reading.
% !TeX root = ../main.tex
\chapter{Causal perturbations and extension of distributions}\label{f:extension}

Send in two weak sources and ask for the part of a detector reading that
vanishes when either source is switched off. For a classical nonlinear
wave this question requires a retarded inverse. For a quantum wave it
matters whether the reading is a change of the mean field or a correlation
of two insertions. The same retarded tree will occur in the change of
mean. A pair of quantum vertices, even when each vertex has already been
defined, will require a different operation.

\section{Differentiating a retarded solution}
A formal series in $\lambda,s,r$ with coefficients in an algebra $A$ is
a list of coefficients $a_{ijk}\in A$ for the monomials $\lambda^i s^j
r^k$, with nonnegative indices. In a product the coefficient at $(i,j,k)$
is the sum over decompositions of that triple between the factors, in the
specified order if $A$ is noncommutative. This is a finite sum even when
the series do not converge anywhere. Differentiation and setting the
parameters to zero extract coefficients, with the usual factorials:
$\partial_s\partial_r|_0$ extracts the $sr$ coefficient, whereas
$\partial_s^2|_0$ gives twice the $s^2$ coefficient (leaving $\lambda$
formal in both cases).

For a differentiable functional $F$ of a smooth field its $k$th variation,
when it exists, is
\[
 F^{(k)}[\phi](h_1,\ldots,h_k)=
 \left.\partial_{s_1}\cdots\partial_{s_k}
 F[\phi+s_1h_1+\cdots+s_kh_k]\right|_{s_1=\cdots=s_k=0} .
\]
The directions $h_i$ are smooth; compact supports or switching functions
keep all displayed integrals finite. For the local polynomial
$F[\phi]=\int\chi\phi^3\,\mathrm{d}V/3!$, the first variation is
$\tfrac12\int\chi\phi^2 h_1\,\mathrm{d}V$ and the second is
$\int\chi\phi h_1 h_2\,\mathrm{d}V$. The latter has kernel
$\chi(x)\phi(x)\delta(x,y)$, where the delta is relative to spacetime
volume. The two variations act at the same point because $F$ is local,
not because the functions $h_1,h_2$ were chosen to be identical.

Take $P=\Box_g-m^2-\xi R$ in a smooth globally hyperbolic spacetime
without boundary.
Use zero incoming data and smooth compactly supported $\chi,f,h$; an
initial slice before the supports suffices. Add
$-\lambda F-\int j\phi\,\mathrm{d}V$ to the free action with bulk
quadratic form $\tfrac12\int\phi P\phi\,\mathrm{d}V$. The equation and its
retarded form are
\begin{equation}
 P\phi=j+\tfrac\lambda2\chi\phi^2,\qquad
 \phi=G_Rj+\tfrac\lambda2 G_R(\chi\phi^2),\qquad j=sf+rh .
 \label{f:ext-retarded}
\end{equation}
Here $PG_R=1$ on tests. The background at $j=0$ is zero. For a nonzero
background solution $\bar\phi$ at source $j_0$, subtract its equation:
the linear operator on $\eta=\phi-\bar\phi$ is $P-\lambda\chi\bar\phi$,
and the new source is $j-j_0$. Its retarded inverse must then be used,
with fixed initial and exterior conditions. We keep the zero background
so that powers of $\lambda$ count only cubic vertices.

Write $u_f=G_R f$, $u_h=G_R h$. Substituting the leading solution on the
right of \eqref{f:ext-retarded} gives
\[
 \phi=su_f+ru_h+\tfrac\lambda2 G_R\!\left(\chi(su_f+ru_h)^2\right)
          +O(\lambda^2).
\]
The $\lambda sr$ coefficient is $G_R(\chi u_f u_h)$, not half this
quantity: there are two assignments of source labels to the square. A
smooth detector $k$ of compact support measures this coefficient as
\begin{equation}
 \int k\,G_R(\chi u_f u_h)\,\mathrm{d}V
       =\int (G_A k)\chi u_f u_h\,\mathrm{d}V .
 \label{f:ext-tree}
\end{equation}
This uses formal self-adjointness and retarded--advanced adjointness on
the same domain. It locates the signal at an overlap of the futures of
both pulses and the past of the detector. Drawing a line for each inverse
gives a rooted tree: two inputs meet, and their product propagates to the
output. Iterating substitutes a new branch at a field factor; it never
pairs two previously existing factors. No state has entered.

\begin{exercise}
In a flat massless region take $\chi=f=h=1$ in the portion of the causal
past of a detector between the initial slice and elapsed time $T$. Compute
the $\lambda sr$ and $\lambda s^2$ field coefficients there. The data
outside that domain may be cut off smoothly.
\end{exercise}
\noindent\emph{Hint.} In the spatially constant calculation $P=-\partial_t^2$
and $u_f=u_h=-t^2/2$. Propagate $t^4/4$ once more with kernel $-(T-t)$;
the mixed answer is $-T^6/120$, the unmixed one half as large.

\section{Ordered insertions: the tree and the two-line loop}
Choose a centered quasifree Hadamard state with two-point function $W$ for
the free field $\Phi$, and set $\hbar=1$. First use Wick ordering in this
reference state. The first subtraction is explicit:
$:\!\Phi(x)\Phi(y)\!:=\Phi(x)\Phi(y)-W(x,y)\mathbf1$.
In matrix elements between vectors generated from the quasifree cyclic
vector by finitely many smeared fields, the ordered pairing rule derived on
p.~\pageref{f:wick-pairings} now pairs each
unsmeared field only with a smeared one. Such a contraction is smooth:
a singular direction cannot have a zero covector in the smeared slot by
\eqref{f:wf-spectrum}. These matrix elements can be restricted to $x=y$;
the result is the Wick square as a distribution in the remaining point.
For higher Wick powers the colons likewise exclude all internal pairings.
Local covariant Wick powers subtract a geometric parametrix rather than
the full $W$; the resulting finite local normalizations must be specified
as well \cite{HWick}.
For a nonderivative potential the Hamiltonian insertion has the opposite
sign to the added Lagrangian term. Thus the finite-time Dyson convention
is $U=\mathcal T\exp[-i\int (j\Phi+\lambda\chi:\!\Phi^3\!:/3!)\,\mathrm{d}V]$,
with later times on the left. Here and below a field at a point is kernel
notation for a spacetime-smeared insertion, not a point operator.

For a change of the mean, the observable evolves as $U^*\Phi(x)U$ up to
the observation time. To first order in an insertion its change is $i$
times the ordered integral of its commutator with $\Phi(x)$. Since
$i[\Phi(y),\Phi(x)]=E(x,y)\mathbf1$, a linear source alone shifts the
field by $u=G_Rj$; further source commutators vanish. Write $U=U_j U_{\rm rel}$
with $U_j$ generated by the source alone. The cubic insertion in the
relative evolution is conjugated by $U_j$ up to its insertion time, so
its field is $\Phi(y)+u(y)$. Its commutator is then
\[
 i[:\!(\Phi(y)+u(y))^3\!:,\Phi(x)]
       =3E(x,y):\!(\Phi(y)+u(y))^2\!: .
\]
Retarded ordering supplies $G_R(x,y)$. Of the terms on the right only
$u(y)^2$ is quadratic in the source. The $3/3!$ from the commutator and
the $2$ from its $sr$ coefficient reproduce \eqref{f:ext-tree}. This
calculation explains how a retarded line arises from quantum evolution;
it is not a prescription to contract two fields with $G_R$.

For a time-ordered correlation the single contraction is instead
\begin{equation}
 G_F=W+iG_A=W^T+iG_R,\qquad P_xG_F(x,y)=i\delta(x,y) .
 \label{f:ext-feynman}
\end{equation}
Indeed $W-W^T=i(G_R-G_A)$, so the first expression equals $W$ to the
future of $y$ and the second equals $W^T$ to its past, off coincidence.
At spacelike separation they agree. This choice defines the free scalar
contraction on the diagonal as well, without multiplying steps at zero;
there is \emph{no} extra $i$ multiplying the time-ordered expectation.
The delta follows by applying $P_x$ to the formula, not to $W$ alone.

For bosonic local insertions the multilinear time-ordered kernels $T_n$
obey \emph{causal factorization}. Split the labels into two nonempty sets
$I,J$. On supports for which no point of group $I$ lies in the causal
past of a point of group $J$,
\begin{equation}
 T_n(F_1,\ldots,F_n)=T_{|I|}(F_I)\,T_{|J|}(F_J) .
 \label{f:ext-factor}
\end{equation}
The right side is the ordinary quantum product, not another ordering
across the groups. To see the reach of this prescription, take a tuple
not all at the same point. The finitely many distinct points have a
maximal element in their causal order. Put \emph{all} labels at this
element into $I$, the others into $J$. Closedness of the causal relation
in a globally hyperbolic spacetime gives open neighborhoods with the
same separation. Such neighborhoods cover the complement of the total
diagonal. Lower products already include their partial coincidences;
causal factorization at lower orders and commutativity at spacelike
separation give agreement on overlaps. This is the off-diagonal step of
the Epstein--Glaser induction \cite{BF,HWT,FR}. A spatial division does
not order the two groups in time.

In the product of two Wick cubes, contract $j$ fields at $x$ with $j$
fields at $y$. The choices of fields and their bijections number
$\binom3j^2j!$, respectively $1,9,18,6$ for $j=0,1,2,3$. Each cross
contraction contributes $W(x,y)$ in the indicated operator order, hence
$G_F(x,y)$ after time ordering away from coincidence. In particular the
two-line term of the quadratic Dyson coefficient is
\begin{equation}
 -\frac{\lambda^2}{4}\int \chi(x)\chi(y)\,G_F(x,y)^2
       :\!\Phi(x)\Phi(y)\!:\,\mathrm{d}V_x\,\mathrm{d}V_y ,
 \qquad x\ne y\ \text{initially} .
 \label{f:ext-fish}
\end{equation}
The factor is $18(-i\lambda)^2/[2!(3!)^2]$; the remaining fields have
not been paired. Two lines between the same two vertices make a loop.
It is not the power of $\lambda$ that identifies the loop: at order
$\lambda^2$ the classical recursion gives the tree
$\tfrac12 G_R(\chi u\,G_R(\chi u^2))$, with $u=G_Rj$.
At a null pair off the diagonal $G_F$ locally selects one of $W,W^T$, so its square
passes the product test of Chapter~\ref{f:microlocal}. At coincidence
\eqref{f:ext-feynman} implies that $\operatorname{WF}(G_F)$ contains all
conormals of the diagonal: differentiation cannot create directions and
$\delta(x,y)$ has them all. The product test therefore fails there. The
object to be constructed is an extension of the \emph{already defined}
separated-vertex distribution, not a point value of $G_F^2$.

\begin{exercise}
Replace $\chi$ in the cubic interaction by $a\chi_1+b\chi_2$. Find the
coefficient of $\lambda^2 ab$ in the two-line term. If an extension of
$G_F^2$ is changed by $c\delta(x,y)$, what change does this coefficient
undergo? Can its nonlocal part persist for disjoint supports?
\end{exercise}
\noindent\emph{Solution.} Exchange $x,y$ in one cross term. The two give
$-\tfrac12$ times the integral in \eqref{f:ext-fish} with switchings
$\chi_1(x)\chi_2(y)$ and no $\lambda^2$. The local change is
$-\tfrac c2\int\chi_1\chi_2:\!\Phi^2\!:\,\mathrm{d}V$. Only this latter
term must vanish for disjoint supports.

\section{Normal scaling and the missing Taylor moments}
At two coincident bulk points in dimension $D$, take coordinates
$(X,z)=(y,x-y)$ in a chart. The diagonal is $z=0$; approaching it means
scaling the $D$ relative coordinates, not scaling $X$. We first fix the
model normal space $\mathbb R^k$. Let $t^\circ$ be a distribution on its
punctured version. For a test $\psi$ supported off zero set
$\psi_a(z)=a^{-k}\psi(z/a)$ and $t^\circ_a(\psi)=t^\circ(\psi_a)$.
Pairing a regular function this way gives the function $t^\circ(az)$;
the factor $a^{-k}$ removes the volume Jacobian.

\begin{definition}
The \emph{scaling degree} at zero is
\[
 \operatorname{sd}_0(t^\circ)=\inf\{d\in\mathbb R:
      a^d t^\circ_a\longrightarrow0\ \text{in distributions as }a\downarrow0\}.
\]
For a distribution defined at zero as well the definition uses all tests,
not just punctured-space tests. A nonzero homogeneous distribution of
degree $-b$ has scaling degree $b$; logarithmic factors in its scaling
law do not increase the infimum.
\end{definition}

\begin{theorem}[Point extension with finite scaling degree, Brunetti--Fredenhagen \cite{BF}]
\label{f:ext-theorem}
A distribution on a punctured neighborhood in $\mathbb R^k$ of finite
scaling degree $\rho$ at zero has an extension of the same scaling degree.
With this degree requirement it is unique for $\rho<k$. For $\rho\ge k$
the freedom is exactly addition of
\begin{equation}
 \sum_{|\alpha|\le\lfloor\rho-k\rfloor}c_\alpha\,\partial^\alpha\delta_0 .
 \label{f:ext-jets}
\end{equation}
No symmetries or field identities have yet been imposed on the constants.
\end{theorem}

Two extensions agree on all tests supported off zero. By the structure
theorem for point-supported distributions \cite{Schwartz,HormI}, their
difference is a finite sum of delta derivatives. On a scaled test a term
$\partial^\alpha\delta_0$ gives
$(-1)^{|\alpha|}a^{-k-|\alpha|}\partial^\alpha\psi(0)$. Prescribing a
Taylor jet of $\psi$ can isolate any highest derivative. This proves the
ambiguity bound and explains why merely requiring \emph{some} extension
would not imply uniqueness even below threshold.

Here is how to use the existence statement. If $m=\lfloor\rho-k\rfloor
\ge0$, take a compactly supported smooth $w$ equal to one near zero and
subtract the jet from the test:
\[
 \psi^\flat(z)=\psi(z)-w(z)\sum_{|\alpha|\le m}
                   \frac{z^\alpha}{\alpha!}\partial^\alpha\psi(0) .
\]
On an annulus of radius $a$ this gains $m+1$ powers of $a$. To see why
that suffices, localize the annulus with a fixed smooth annular cutoff
$\beta(z/a)$. Its pairing with $t^\circ$ is
$a^k t^\circ_a(\beta(\cdot)\psi^\flat(a\,\cdot))$. For any $d>\rho$,
$a^d t^\circ_a$ is bounded on tests of bounded seminorms on that fixed
annulus; Taylor's formula bounds each needed derivative of
$\psi^\flat(a\,\cdot)$ there by $O(a^{m+1})$. The shell is therefore
$O(a^{k+m+1-d})$. Choose $\rho<d<k+m+1$; the bound is summable on a
smooth annular partition with radii $2^{-j}$. If $\rho<k$, no jet has to
be subtracted and the same argument works with exponent $k-d>0$.
Assigning values to the finitely many subtracted Taylor moments completes
the extension; the theorem controls its scaling on all tests. It is the
codimension $k$, not the dimension of the whole configuration space, that
pays for the gain from the volume element.

\section{A marginal subtraction and its scale residue}
In the flat massless four-dimensional theory the time-ordered contraction
is $(4\pi^2)^{-1}(z^2+i0)^{-1}$ with $z^2=|\mathbf z|^2-(z^0)^2$.
Its square on punctured relative space has scaling degree four; it is
still a boundary value \emph{on} the noncoincident light cone. To perform
a subtraction by an ordinary radial integral, first take a different,
Euclidean distribution with the same marginal degree:
\[
 t_E^\circ(z)=\frac{1}{16\pi^4|z|^4},\qquad z\in\mathbb R^4\setminus\{0\},
\]
where $|z|$ is Euclidean length. Its integral against a test constant near
zero would contain $\int_0^1\mathrm{d}r/r$. The theorem calls for one
Taylor moment. For a smooth radial cutoff profile $w$ equal to one near
zero and zero at large radius, and a length $\ell>0$, put
\begin{equation}
 R_\ell(\psi)=\frac1{16\pi^4}\int_{\mathbb R^4}
  \frac{\psi(z)-\psi(0)w(|z|/\ell)}{|z|^4}\,\mathrm{d}^4z .
 \label{f:ext-subtraction}
\end{equation}
The numerator is $O(|z|)$ near zero by the mean value bound. Thus the
radial integral is absolutely convergent, controlled by a first derivative
seminorm of the test there and its supremum away from zero. This defines
a distribution and leaves every punctured-space pairing unchanged. The
subtraction is needed for existence; choosing which constant to assign
to the removed moment is a separate finite choice.

Only the cutoff changes when $\ell$ changes. Since the area of the unit
three-sphere is $2\pi^2$, radial integration gives
\begin{align}
 (R_{\ell'}-R_\ell)(\psi)
 &=-\frac{\psi(0)}{8\pi^2}\int_0^\infty
       [w(r/\ell')-w(r/\ell)]\,\frac{\mathrm{d}r}{r}
   =-\frac{\log(\ell'/\ell)}{8\pi^2}\psi(0) .
 \label{f:ext-residue}
\end{align}
For the last equality differentiate the integral in $\log\ell'$: after
$u=r/\ell'$ the derivative is $-\int_0^\infty w'(u)\,\mathrm{d}u=1$.
The integral is zero at $\ell'=\ell$. Changing variables $z=a v$ in
\eqref{f:ext-subtraction} now yields the full scaling law
\[
 a^4(R_\ell)_a=R_{\ell/a}=R_\ell+\frac{\log a}{8\pi^2}\delta_0 .
\]
Multiplication by any additional positive power of $a$ gives zero in the
limit, and off zero the scaling degree was already four. Thus the
extension preserves scaling degree but not exact homogeneity. Adding a
finite multiple of $\delta_0$ cannot cancel the logarithm because the
delta itself scales homogeneously. The residue just computed uses a
Euclidean radial measure; it is not the complex scale change of the
Lorentzian time-ordered contraction.

\begin{exercise}
Recover the Lorentzian scale change without using a Euclidean sphere on a
light cone. Write $q=z^2$, $q_+=q+i0$, and extend the bracket below
without increasing its scaling degree. Why does the result
\[
 L_\ell=-\frac1{64\pi^4}\Box\!\left[\frac{\log(q_+/\ell^2)}{q_+}\right]
\]
agree off zero with the scalar kernel in \eqref{f:ext-fish}? Determine
$L_{\ell'}-L_\ell$.
\end{exercise}
\noindent\emph{Solution.} Off the cone use
$\Box h(q)=4q h''(q)+8h'(q)$: for $h(q)=\log(q/\ell^2)/q$ the result
is $-4/q^2$. At any nonzero point of the cone $\mathrm{d}q\ne0$, so the
one-variable upper boundary value in $q$ gives the same identity there.
The bracket has scaling degree two (with a logarithm); extend it uniquely
at the tip before differentiating. Two derivatives raise scaling degree
by at most two, and the punctured square already has degree four.
Finally \eqref{f:ext-feynman} gives $\Box(1/q_+)=4\pi^2 i\delta_0$.
The difference is $i\log(\ell'/\ell)\delta_0/(8\pi^2)$. The factor $i$
and its sign depend on the Lorentzian source convention.

\section{Where a contact can live, and which identities it must obey}
For the total diagonal of $n$ bulk vertices in spacetime dimension $D$,
there are $k=(n-1)D$ relative coordinates and $D$ coordinates along the
diagonal. Near a smooth collision stratum use a product chart $(X,z)$ with
$z=0$ on the stratum. Dilate only $z$, using tests
$\psi_a(X,z)=a^{-k}\psi(X,z/a)$. With a finite normal scaling bound
$\rho$ uniform on compact tangential test families, two extensions
obeying that bound can differ locally only by terms of the form
\eqref{f:ext-jets} in $z$, with distributional coefficients in $X$.
This follows from the local structure of distributions supported on a
submanifold and the same scaled-jet test. Smoothness or a prescribed
wavefront set for the coefficients is an additional condition. If they are
singular at the edge of a face, that smaller stratum needs separate
analysis. Power counting classifies possible contacts once the distributions
away from the stratum and their normal estimates are known.

If all $n$ bulk vertices collide \emph{on} a face of codimension $q$,
the common point must move to the face as well: $k=(n-1)D+q$. Vertices
constrained from the start to that face instead have total diagonal of
codimension $(n-1)(D-q)$ in their own configuration space. For example,
draw a plane in the flat four-dimensional fish calculation. The uncut
kernel has no dependence on the normal position $X_\perp$ of the common
point. Under joint scaling of $(z,X_\perp)$ it still has degree four,
whereas $\delta^{(4)}(z)\delta(X_\perp)$ has degree five. With that bound
there is no extra sheet mass term. More generally an artificial division
supplies no new physical coupling: split a switching function smoothly and
sum its terms using the same whole-space extension. A true wall with its
own action, boundary condition or fields supplies new local data. Its
face and corner counterterms must be tested against those data, not
inferred by declaring the wall transparent. A bulk existence theorem does
not solve a theory on an arbitrary boundary.

Symmetries test the \emph{extended} products. At zero interaction the
massless flat scalar has the shift current $j^\mu=-\partial^\mu\Phi$,
with $j^0=\partial_t\Phi$. It is conserved in the bulk field equation;
its time-ordered insertion with a field is not divergence-free at
coincidence. With the derivative normalization of the free contraction,
\begin{equation}
 \partial_\mu^x\langle\mathcal T j^\mu(x)\Phi(y)\rangle
   =-\Box_x G_F(x,y)=-i\delta(x,y) .
 \label{f:ext-ward}
\end{equation}
The first derivative of the ordering has no contact proportional to the
equal-time field--field commutator, which vanishes. The divergence then
detects $[j^0(t,\mathbf x),\Phi(t,\mathbf y)]=-i\delta^{(3)}(\mathbf
x-\mathbf y)$. A test equal to one near $y$ reads the coefficient $-i$.
Thus extending a classically conserved expression by setting its ordered
divergence to zero would give the wrong generator on the inserted field.
The cubic interaction explicitly breaks this shift; such breaking is not
a quantum anomaly. More generally, locality, covariance, power counting,
unitarity and Ward identities restrict the finite terms allowed in an
extension \cite{HWick,HWT,FR}.

In a gauge theory the obstruction has the cochain form of
Chapter~\ref{f:gauge-complex}. Let $s$ be the nilpotent linearized
symmetry differential on admissible local densities, raising ghost number
by one. Suppose all Ward identities through order $n-1$ hold. The first
possible breaking $\mathcal A_n$ is a local contact of ghost number one:
applying the symmetry consistency identity a second time leaves
$s\mathcal A_n=0$ at this order, modulo permitted total derivatives,
since the terms containing lower-order breakings vanish. Adding a finite
local counterterm $B$ of ghost number zero at order $n$ changes the
breaking to $\mathcal A_n+sB$ in this sign convention. Thus the relevant
cohomology is closed local breakings modulo such exact ones, not the
space of arbitrary distributions on the diagonal.

For example take an Abelian potential $A$ and odd ghost $c$ with
$sA_\mu=\partial_\mu c$, $sc=0$. Suppose the dimension and other
normalizations permit the counterterm $B=\tfrac\kappa2\int A_\mu A^\mu\,\mathrm{d}V$.
Then
\[
 sB=\kappa\int A^\mu\partial_\mu c\,\mathrm{d}V
   =-\kappa\int c\,\nabla_\mu A^\mu\,\mathrm{d}V
      +\kappa\int_{\partial M} c A^\mu\,\mathrm{d}\Sigma_\mu ,
\]
with the directed surface element in the last integral. A candidate
breaking $+\kappa\int c\,\nabla_\mu A^\mu\,\mathrm{d}V$ is removable in
the compactly supported bulk problem. On a true wall the flux must vanish
by the physical conditions or be canceled in the combined boundary
identity before the same total derivative can be discarded. The example
identifies a coboundary under stated assumptions; it neither produces an
anomaly nor excludes other cohomology classes.

Finally all these equations are statements in a specified theory and, if
several global sectors exist, in a specified sector. Even if two sectors
have indistinguishable local singularities, no step in the Taylor
subtraction assigns a relative probability or phase to them. To combine
sector answers one needs a common observable and an independently chosen
measure: mixing normalized states needs nonnegative weights summing to
one; adding amplitudes is a different prescription. Conservation and the
extension theorem dictate neither. Equality after a weighted sum would
also be weaker than equality in each sector, since two discrepancies can
cancel.

\begin{exercise}
As a power-counting test consider $|z|^{-6}$ off zero in Euclidean four
space. Which normal contacts can change a degree-preserving extension?
If the result must also be a rotation-invariant scalar, which survive?
Does rotation invariance determine their coefficients?
\end{exercise}
\noindent\emph{Solution.} The excess degree is $6-4=2$. Before symmetry,
delta derivatives through order two are allowed. There is no invariant
vector for the order-one coefficient and an invariant symmetric
rank-two tensor is proportional to the Euclidean metric. Thus
$c_0\delta_0+c_2\Delta\delta_0$ remains. Neither constant is fixed by
this count, which imposed no scaling law beyond the degree bound.

\part{An application}
% !TeX root = ../main.tex
\chapter{A field on a compact spatial circle}\label{f:capstone}
Consider a small static source placed at a point of a compact one-dimensional space. Does the static field produced at the source depend on how field values on that space are identified? And if the static response at that point is finite, is the vacuum fluctuation there finite too? We shall obtain three different kernels for the same system. Each answers exactly one of these questions: the inverse of a positive spatial operator; the difference of causal inverses of a spacetime equation; and a positive two-point function. They involve the same frequencies with different weights.

\section{The domain comes before the operator}
Take two intervals of lengths $d_1,d_2>0$ at each real time. Pair their ends cyclically at equal time, using the common flat metric $-dt^2+dx_i^2$, increasing future time, and no sheet action. The base is $\mathbb R\times S^1_L$, where $L=d_1+d_2$: use $X=x_1$ on the first interval and $X=d_1+x_2$ on the second, modulo $L$. At each paired endpoint the signed chart of Chapter~\ref{f:submanifolds} proves that these are smooth coordinates across the join, not just a parametrization by equivalence classes. Every constant-time circle is Cauchy. Indeed a causal curve parametrized by $t$ has spatial speed at most one. If the range of $t$ had a finite endpoint the spatial curve would have a limit on the compact circle, and the curve could be continued. Thus a future-directed inextendible curve passes each time exactly once.

Use a real Euclidean line as field fiber. Choose signs $a,b$ for the transitions, in cyclic order, and no fixed frames on the two pieces. Changing the second frame by $a$ makes the first transition $+1$ and the second $H=ba$. In a lifted real coordinate the field component then satisfies
\begin{equation}
 u(t,X+L)=H u(t,X),\qquad H=(-1)^\epsilon,\quad \epsilon\in\{0,1\}.\label{f:twist-domain}
\end{equation}
The coefficient of the derivative in this common frame has the same transition. In fact if at first only values are matched, varying the homogeneous scalar action gives at the first join the coefficient $-u_{1,x}(d_1)+a\,u_{2,x}(0)$ multiplying the common variation on side 1. Its vanishing matches both slopes in the common frame, and the return join works the same way. This derivation is why the domain of the operator below involves \emph{two} endpoint conditions; it is not obtained by varying each interval with a free, unrelated endpoint.

For $m>0$ let $A=-d^2/dX^2+m^2$ on the complexification of the space of spatial sections. Its Hilbert space in this description is $L^2(0,L)$ with inner product $\int_0^L\bar f g\,dX$. Set $D(A)=\{f\in H^2(0,L): f(L)=H f(0),\ f'(L)=H f'(0)\}$. These are traces of $H^2$ functions, not values assigned to arbitrary $L^2$ representatives. Integration by parts twice gives the boundary form $[\bar f'g-\bar f g']_0^L=0$ for $f,g\in D(A)$, since $H^2=1$ as a sign. Also $\langle f,Af\rangle=\|f'\|^2+m^2\|f\|^2$. To see self-adjointness rather than merely symmetry, it is useful to exhibit a complete diagonalization.

\subsection{Why the modes are complete}\label{f:circle-fourier}
Set $k_n=(2\pi n+\pi\epsilon)/L$ and $e_n(X)=L^{-1/2}e^{ik_n X}$ for $n\in\mathbb Z$. Multiplication by $e^{-i\pi\epsilon X/L}$ identifies the twisted Fourier family with the ordinary periodic one and is unitary on $L^2$. For completeness of the periodic family, in the $2\pi$ angle coordinate use the Fej\'er kernels
\[
 F_N(\theta)=\frac1N\left|\sum_{j=0}^{N-1}e^{ij\theta}\right|^2 .
\]
They are nonnegative trigonometric polynomials of integral $2\pi$. Outside any arc about zero the geometric-sum formula bounds them by a constant divided by $N$. Convolution with $F_N/(2\pi)$ therefore converges uniformly to every continuous periodic function by uniform continuity. Continuous functions are dense in $L^2$ (already smooth functions supported in the open interval are dense there). This proves completeness. On the domain, integration by parts gives $(Af)_n=\omega_n^2f_n$, with $\omega_n=\sqrt{k_n^2+m^2}$. Conversely a sequence with $\sum(1+k_n^2)^2|f_n|^2<\infty$ has partial sums Cauchy in $H^2$, which converges on the interval and retains both trace conditions. The two descriptions of the domain thus agree. Multiplication by the real numbers $\omega_n^2$ on this maximal square-summability domain is self-adjoint: testing an adjoint candidate against individual basis vectors forces exactly the same coefficient condition.

\section{A point source and a static measurement}\label{f:static-circle}
A unit point source at $X=0$ is specified relative to a chosen unit vector in that fiber. Transport all reported values to the frame on the lift. The static response is $A^{-1}\delta_0$ in that frame, with value at the source given by the kernel $K_H(0)$. In one spatial dimension $\delta_0$ is not in $L^2$, but the pairing with an $H^1$ section is continuous: use its Fourier series and Cauchy--Schwarz with $\sum(1+k_n^2)^{-1}<\infty$. Despite the point source, the inverse has enough smoothing for its value to exist. Indeed its Fourier kernel with second argument zero is
\[
 K_H(X)=\frac1L\sum_{n\in\mathbb Z}\frac{e^{ik_n X}}{k_n^2+m^2} .
\]
The sum converges absolutely and uniformly. Testing against a smooth section shows distributionally $AK_H=\delta_0$, since the twist-Fourier delta evaluates the test at zero. Changing the frame at the source changes both the components of the unit source and the returned field. The diagonal kernel, as the scalar of a linear map on the same fiber, does not depend on that choice.

There is a second computation which does not need a summation identity. On the real line, $K_\infty(X)=e^{-m|X|}/(2m)$ solves the homogeneous equation off zero and has derivative jump $-1$, exactly what the leading $-d^2/dX^2$ requires. Add its translates with factors $H^j$. The series $\sum_{j\in\mathbb Z}H^j K_\infty(X+jL)$ converges locally uniformly, has precisely this jump per fundamental interval, and obeys the required transition. It agrees with the spectral answer by uniqueness (a homogeneous difference has zero positive quadratic form). On $0\le X\le L$, summing the two geometric series gives
\begin{align}
 K_H(X)&=\frac{e^{-mX}+H e^{-m(L-X)}}{2m(1-He^{-mL})},\nonumber\\
 K_+(0)&=\frac{1}{2m}\coth\frac{mL}{2},\qquad
 K_-(0)=\frac{1}{2m}\tanh\frac{mL}{2}.\label{f:static-measurement}
\end{align}
For example $K_H(L)=H K_H(0)$ and $K'_H(0^+)-H K'_H(L^-)=-1$. These are checks at the source, where matching equal slopes would have been incorrect. Far from the massless limit both sectors tend to the infinite-line response as $mL\to\infty$ at fixed $m$. As $m\downarrow0$ at fixed $L$, $K_+(0)=1/(Lm^2)+L/12+O(m^2L^3)$, while $K_-(0)\to L/4$. In the periodic sector a constant mode obstructs inversion at zero mass for a source with nonzero total integral. In the twisted sector no such mode exists. A local detector distinguishes two global domains without any force added at either artificial join.

\section{The causal inverse and the commutator}
Return to $P=\Box-m^2=-(\partial_t^2+A)$. In each spatial mode, a unit time impulse forces the retarded response
\[
 g^R_n(\tau)=-\mathbf1_{\tau>0}\frac{\sin(\omega_n\tau)}{\omega_n},\qquad
 g^A_n(\tau)=\mathbf1_{\tau<0}\frac{\sin(\omega_n\tau)}{\omega_n},\quad \tau=t-t'.
\]
Both are continuous across zero with derivative jump $-1$. This proves $-(\partial_\tau^2+\omega_n^2)g_n^{R/A}=\delta$ as a distribution; the support distinguishes the two. For a smooth compact-in-time test source $f(t,X)$, spatial Fourier coefficients decay faster than any power locally uniformly in $t$ with all derivatives, by integration by parts and the twist conditions. Convolving them in time with $g_n^R$ and reconstructing the Fourier series gives a smooth solution with zero past data. Homogeneous energy uniqueness and finite propagation as in Chapter~\ref{f:waves} identify this solution as the retarded inverse, not just an operator whose output happens to vanish before its input starts. The advanced inverse is obtained with zero future data. Subtracting gives the causal distribution kernel
\[
 E(t,X;t',X')=-\sum_n e_n(X)\overline{e_n(X')}\frac{\sin(\omega_n\tau)}{\omega_n} .
\]
The sum is used after smearing. Replacing $(t,X)$ by $(t',X')$ changes its sign since conjugation interchanges $n$ with $-n-\epsilon$. Unlike $K_H$, it solves the \emph{homogeneous} spacetime equation in both variables: its derivative jumps in the two inverses cancel. The field commutator is $iE$ in units $\hbar=1$. Nothing in the subtraction asks for an inner product on a space of particles or a ground state.

\section{Which kernel may be evaluated at a point?}
There is time-translation symmetry on this cylinder. Selecting modes $e^{-i\omega_n t}$ of positive energy gives the candidate two-point kernel
\begin{equation}
 W_H(t,X;t',X')=\sum_n\frac{e_n(X)\overline{e_n(X')}}{2\omega_n}\,e^{-i\omega_n(t-t')} .\label{f:circle-two-point}
\end{equation}
For a complex test $f$ put $c_n(f)=\int e^{i\omega_n t}\overline{e_n(X)} f(t,X)\,dt\,dX$. Smearing $\bar f$ in the first slot and $f$ in the second gives $\sum_n |c_n(f)|^2/(2\omega_n)\ge0$. The mode equation makes the kernel a bisolution, and $W_H-W_H^{\mathrm{op}}=iE$ by the sine identity. It defines a positive quasifree state by Chapter~\ref{f:quantum}. These properties still do not say the kernel is a function at coincidence. In fact a formal equal-point sum is $(2L)^{-1}\sum_n1/\omega_n$, divergent like the harmonic series. The difference between this and the convergent static sum is one power of frequency, not a change of geometry. Nor does replacing a vacuum by another admissible state remove the local singularity: adding $c\,u(x)u(x')$ with $c\ge0$ and $u$ a real smooth solution changes the positive symmetric part while leaving both the commutator and the singular directions unchanged.

One can also see the local singularity directly. On the line at equal times the positive-frequency kernel off its diagonal is
\[
 W_\infty(0,X)=\frac1{2\pi}\int_{\mathbb R}\frac{e^{ikX}}{2\sqrt{k^2+m^2}}\,dk
   =\frac1{4\pi}\int_0^\infty e^{-m^2s-X^2/(4s)}\frac{ds}{s} .
\]
The first integral is a Fourier transform in the distributional sense; insert a Gaussian cutoff or smear first. The second follows by the Laplace representation of the inverse square root and the Gaussian Fourier integral. At $X\ne0$ both endpoints of the $s$ integral converge, and near $X=0$ the singularity is logarithmic. For instance, $X^2\ll s\ll m^{-2}$ contributes the logarithm, and substituting $s=X^2 v$ at the small endpoint bounds the remainder from that endpoint. The kernel is $K_0(m|X|)/(2\pi)$ in terms of the modified Bessel function. At spacelike separation replace $|X|$ by $\sqrt{X^2-\tau^2}$: writing $k=m\sinh\theta$ shows $dk/\sqrt{k^2+m^2}=d\theta$ is boost invariant, so the time difference can be transformed to zero.

To find the singularity also at unequal times, first damp the frequency integral by $e^{-a\omega}$ for real $a>0$, $\omega=\sqrt{k^2+m^2}$. The Gaussian identity
\[
 \frac{e^{-a\omega}}{\omega}=\frac1{\sqrt\pi}\int_0^\infty s^{-1/2}e^{-s\omega^2-a^2/(4s)}\,ds
\]
adds $a^2$ to $X^2$ in the integral above. Thus the damped kernel is $K_0(m\sqrt{X^2+a^2})/(2\pi)$. Both the damped mode integral and this Bessel expression are analytic for $\operatorname{Re}a>0$ with the square root of positive real part. Set $a=\varepsilon+i\tau$ there and take its distributional boundary value as $\varepsilon\downarrow0$. Write $Q=X^2-(\tau-i0)^2$ for this prescription. The expansion $K_0(z)=-I_0(z)\log(z/2)+b(z^2)$, with $b$ analytic near zero and $I_0(z)=\sum_{j\ge0}(z^2/4)^j/(j!)^2$ \cite[Eq.~10.31.2]{DLMF}, shows that the local singular part modulo smooth terms is
\begin{equation}
 -\frac{1}{4\pi}I_0(m\sqrt Q)\log(Q/\ell^2),\qquad \ell>0 .\label{f:circle-subtraction}
\end{equation}
Changing the reference length $\ell$ changes a smooth function. This is the logarithmic subtraction in \emph{two} spacetime dimensions; inserting the four-dimensional $1/\sigma$ term of \eqref{f:wf-parametrix} here would subtract the wrong singularity.

Periodizing this flat kernel in the spatial difference with factors $H^j$ gives \eqref{f:circle-two-point}. Indeed taking the $n$th Fourier coefficient over an interval unfolds the translated intervals to the full line: $e^{ik_n jL}=H^j$ cancels the chosen image sign. This calculation can first be performed with a positive imaginary-time regulator $e^{-\varepsilon\omega_n}$ and then against smooth tests. Near any point of the diagonal choose the coordinate differences small enough that all $j\ne0$ images are spacelike with a fixed positive separation. Their sum is smooth there, with all derivatives locally uniformly convergent: for large spacelike distance $r$, the representation $K_0(mr)=\int_0^\infty e^{-mr\cosh v}dv$ and its differentiated versions give exponential bounds when $m>0$. Finitely many nearby nonzero images are individually smooth. Thus only the unshifted Minkowski singularity remains locally: the same subtraction \eqref{f:circle-subtraction} works in both sectors.

There is now a finite fluctuation that can be compared. Define $\langle\Phi^2\rangle_{H,\ell}$ by subtracting \eqref{f:circle-subtraction} from the positive-frequency two-point function and then taking coincidence in a common local frame. The scalar result does not depend on a reversal of that frame. At fixed $m,L$ the unshifted term, including its smooth remainder, cancels in the difference. The even images cancel as well; the images at $jL$ and $-jL$ have equal values. Hence
\begin{equation}
 \langle\Phi^2\rangle_{+,\ell}-\langle\Phi^2\rangle_{-,\ell}
   =\frac{2}{\pi}\sum_{j=0}^\infty K_0\big(m(2j+1)L\big)>0 .\label{f:relative-fluctuation}
\end{equation}
The integral representation above makes every term positive and gives exponential convergence. The difference is independent of $\ell$, not because either unrenormalized variance was finite, but because the two vacua have the same local subtraction. Only paths returning after an odd winding distinguish the two transports. The conclusion compares these particular vacua, not arbitrary states sharing their commutation relations.

\Needspace{.52\textheight}
\section*{Exercises}
\begin{exercise}In the twisted sector put the detector half a circumference from the source. Evaluate $K_-(L/2)$ and interpret the result without invoking energy-level degeneracy.\end{exercise}
\noindent\emph{Solution.} In the numerator of the image formula both exponentials are $e^{-mL/2}$, with opposite signs. The value is zero. Paths around the two sides have equal attenuation and opposite transport signs; in the image series they cancel in pairs.
\begin{exercise}For a source with zero spatial average in the periodic massless theory, produce a static solution and describe exactly what remains undetermined.\end{exercise}
\noindent\emph{Solution.} Divide the nonconstant Fourier components by $k_n^2$ and choose any constant component. The zero-average condition is the range condition; fixing the average of the solution selects a representative. At the level of the time equation the massless constant mode has two real initial data, not an oscillator vacuum of frequency zero.
\begin{exercise}Restrict \eqref{f:circle-two-point} to a spacelike time slice in both variables. Why does the singular-direction criterion permit this restriction but not an ordinary scalar value at a coincident point?\end{exercise}
\noindent\emph{Hint.} Locally the kernel depends on the two coordinate differences, so a singular covector pair must have the form $(k,-k)$. Apply Theorem~\ref{f:wf-propagation} to the homogeneous equation in the difference variables: $k$ must be null. At a spacelike slice a nonzero conormal is timelike. On a point all covectors are conormal; the sufficient criterion no longer applies there. In this example the positive divergent sum gives a separate direct obstruction to assigning the unsubtracted variance at a point. Not all pairs of points need be coincident to be singular in spacetime.

\backmatter% !TeX root = main.tex
\chapter*{Reading alongside Lorentzian Sewing}\addcontentsline{toc}{chapter}{Reading alongside Lorentzian Sewing}\markboth{Companion guide}{}
The bold numbers below refer to the chapters or appendices of \emph{Lorentzian Sewing: Geometry, Dynamics, and Observables}, the volume this text accompanies. Chapter, section, equation and page references after them refer to the present book. This distinction matters: a theorem giving the conditions for a construction is not the assertion that those conditions hold in every proposed field theory.

\begin{description}[style=nextline,leftmargin=0pt,labelindent=0pt,labelwidth=0pt,labelsep=0pt,font=\normalfont\bfseries]
\item[1. Sewing spacetime regions.] For the quotient topology use Theorem~\ref{f:q-descent}; for an actual smooth atlas use the sector and overlap test on p.~\pageref{f:sector-theorem}. Regular levels, normal quotients and surviving exterior walls are built in Chapter~\ref{f:submanifolds}. A causal test is separate: apply the shifted cylinder theorem on p.~\pageref{f:shift-cylinder}.
\item[2. Bundles and connections at a cut.] The bundle-gluing theorem on p.~\pageref{f:bundle-gluing} constructs fibers, not merely their coordinate values. Check a proposed lift by \eqref{f:morphism-overlap}; compare connection forms only after pulling to the same base, as in \eqref{f:connection-change}. The two sign orbits and their stabilizers appear beside \eqref{f:sign-groupoid}.
\item[3. Smooth fields and gauge transformations.] Start with the affine space of first jets in \S\ref{f:jet}. Proposition~\ref{f:jet-pasting} gives necessity and sufficiency for \emph{actual} half-sections, including mixed orders and the condition at a remaining exterior wall. The same test is applied to the arrows of the gauge groupoid, not to a set from which the arrows have already been removed.
\item[4. Varying a region.] Localize and integrate an actual top form using Theorem~\ref{f:partition}; Stokes is proved on p.~\pageref{f:stokes}. In \S\ref{f:first-variation} the action differentiates on vertical field variations and yields both an Euler operator and a boundary covector. Fixed and released trace problems therefore have different stationarity conditions.
\item[5. Covariant phase space and boundary charges.] Chapter~\ref{f:variation} keeps the wall action and its corner potential in the integrated slice potential. Chapter~\ref{f:symplectic} derives the Hamiltonian convention and the Noether charge, and Theorem~\ref{f:regular-reduction} states the regular hypotheses for reduction. The gauge-complex boundary exercise on p.~\pageref{f:gauge-boundary-ex} tests whether a parameter is proper.
\item[6. Causal response and the Peierls bracket.] Pair the Hessian with dual densities by \eqref{f:green-variation}; check both support and adjoint domains in \S\ref{f:green-domains}, including the one-sided extension on p.~\pageref{f:one-sided-extension}. Compact source representatives and their Hamiltonian solutions are related by Theorem~\ref{f:source-exact}. For a gauge equation the relevant inverse is instead the degree-lowering homotopy of Proposition~\ref{f:green-homotopy}.
\item[7. Transmission and reconstruction.] Chapter~\ref{f:variation} recovers normal derivatives from value and conjugate momentum on one inverse branch, and gives a counterexample with two regular branches. Compare that local statement with the symbol and the evolution estimates in \S\ref{f:wave-symbol} and \S\ref{f:wave-energy}. The last exercise of Chapter~\ref{f:variation} distinguishes choosing a gauge from transmitting an observable derivative.
\item[8. Sewing a causal response.] The trace and its right inverse on a flat face are in Theorem~\ref{f:plane-trace}, and the weak flux is in Proposition~\ref{f:weak-flux}. The joined inverse on p.~\pageref{f:sewn-inverse} is conditional on an inverse on the produced mismatches and on causal uniqueness. The impedance example realizes the required history map rather than assuming invertibility of an abstract matrix.
\item[9. Presymplectic sewing and observables.] A kernel is taken \emph{after} restricting the form (Example~\ref{f:two-rotor}). In Chapter~\ref{f:quantum} the invariant algebra of two open transports gives a closed loop observable; preserving brackets among already known generators need not generate this observable.
\item[10. Junctions, elimination and refinement.] Use \S\ref{f:glue-space} for simultaneous sectors, then \S\ref{f:delayed-schur} for Schur elimination on histories. The latter keeps the unsolved boundary data as variables. Geometric compatibility alone says nothing about whether an inverse pivot preserves the direction of travel in time.
\item[11. The linear commutator.] Apply \eqref{f:label-sign} on source labels, not on unreduced field components. Chapter~\ref{f:quantum} explicitly constructs a CCR algebra, a state and its GNS representation as three operations. Proposition~\ref{f:gaussian-positivity} tests candidate mixed covariances; Corollary~\ref{f:wick-pairings} derives all ordered moments of a centered quasifree state. The operator-algebra meaning of a sharp cut is discussed on p.~\pageref{f:local-factors}.
\item[12. Perturbations about a sewn solution.] Linearization and domains come from Chapters~\ref{f:variation} and \ref{f:waves}. Equation~\eqref{f:ext-retarded} separates the background, the free inverse and the nonlinear source. Equation~\eqref{f:ext-tree} then calculates a particular mixed response; the next section compares it with ordered insertions. Formal orders are algebraic coefficients, not a statement of convergence at finite coupling.
\item[13. Time-ordered products.] Theorems~\ref{f:wf-pullback} and \ref{f:wf-hadamard} specify singular directions and permissible pullbacks; \eqref{f:ext-fish} is a first product off the diagonal. Theorem~\ref{f:ext-theorem} bounds the freedom at coincidence. The scale residue \eqref{f:ext-residue} and Ward contact \eqref{f:ext-ward} are distinct constraints; neither assigns physical weights to global sectors.
\end{description}

\paragraph{The earlier appendices.} For A (orientations and flux), use Chapters~\ref{f:forms} and \ref{f:variation}. For B (traces and hyperbolic domains), use Chapters~\ref{f:distributions}--\ref{f:waves}. For C (Hadamard tests), the geodesic input is constructed in Chapter~\ref{f:connections} before the theorem on p.~\pageref{f:wf-hadamard}; the singular-product and extension arguments are in Chapters~\ref{f:microlocal} and \ref{f:extension}. The two-strip construction of D reappears in Chapter~\ref{f:capstone}, where it determines three different kernels. Compare the static response \eqref{f:static-measurement} with the difference of subtracted vacuum fluctuations \eqref{f:relative-fluctuation}.

\paragraph{What remains an input.} Among the results for which we supply hypotheses and an application rather than a full general proof are the smooth collar and convex-neighborhood existence theorems, the global spacetime splitting and normally hyperbolic Cauchy theorems, Plancherel's theorem, the free proper action quotient theorem, and the microlocal pullback, propagation, Hadamard and extension theorems. The text does prove, for example, the flat sharp trace bound, finite-energy uniqueness in the smooth flat model, the compact-source exactness argument, GNS, and the local counterterm ambiguity once the corresponding inputs are available. No general boundary Green theorem, proof of quantization of interacting gauge fields, or classification of physical defects is assumed to follow from those results.

% !TeX root = main.tex
\chapter*{Index of concepts}\addcontentsline{toc}{chapter}{Index of concepts}\markboth{Index of concepts}{}
The page given is a place to learn or use the concept, not merely its first occurrence. The PDF page links lead to the corresponding construction.
\newcommand{\fentry}[2]{\par\noindent\hangindent=1em\hangafter=1 #1, \pageref{#2}\par}
\newcommand{\fcontinued}[3]{\par\noindent\hangindent=1em\hangafter=1 #1, \pageref{#2}, \pageref{#3}\par}
\begingroup\setlength{\columnsep}{1.7em}\raggedright
\subsection*{A--G}
\begin{multicols}{2}
\fentry{Adjoint boundary domain}{f:green-domains}
\fentry{Affine space of derivatives}{f:jet}
\fentry{Algebra by generators and relations}{f:polynomial-ccr}
\fentry{Associated bundle}{f:associated-bundle}
\fentry{Atlas and chart changes}{f:stereo}
\fentry{Basis of a topology}{f:top-basis}
\fcontinued{Boundary action and charge}{f:first-variation}{f:gauge-boundary-ex}
\fentry{Boundary value, distributional}{f:wf-bv}
\fentry{Bundle map over a base map}{f:morphism-overlap}
\fentry{Bump function}{f:partition}
\fentry{Canonical commutation relations}{f:polynomial-ccr}
\fentry{Causal factorization}{f:ext-factor}
\fentry{Causal propagator}{f:causal-propagator}
\fentry{Cauchy temporal function}{f:cauchy-splitting}
\fentry{Characteristic hypersurface}{f:wave-symbol}
\fentry{Clutching and cocycles}{f:bundle-gluing}
\fentry{Cohomology and chain homotopy}{f:gauge-complex}
\fentry{Collar and corner}{f:glue-space}
\fentry{Compactness and separation}{f:q-compact-hausdorff}
\fentry{Completion and Hilbert dual}{f:hilbert-representation}
\fentry{Connected component}{f:q-components}
\fentry{Connection, change of frame}{f:connection-change}
\fcontinued{Conormal space, conormal test}{f:normal}{f:wf-pullback}
\fentry{Constraint, regular reduction}{f:regular-reduction}
\fentry{Continuity and homeomorphism}{f:top-continuity}
\fentry{Convex normal neighborhood}{f:convex-normal}
\fentry{Counterterms}{f:ext-jets}
\fentry{Covariance, positivity}{f:gaussian-positivity}
\fentry{Curvature}{f:curvature}
\fentry{Deformation complex}{f:deformation-complex}
\fentry{Density and orientation}{f:orientation-density}
\fentry{Derivation at a point}{f:tangent-coords}
\fentry{Differential forms}{f:forms}
\fentry{Distribution, order and support}{f:distribution-estimate}
\fentry{Energy and finite propagation}{f:energy-estimate}
\fentry{Euler expression}{f:first-variation}
\fentry{Exact sequence of source labels}{f:source-exact}
\fentry{Exponential map}{f:convex-normal}
\fentry{Extension across a collision}{f:ext-theorem}
\fentry{Exterior derivative and Lie derivative}{f:exterior}
\fentry{Factor and split inclusion}{f:local-factors}
\fentry{Fej\'er kernel}{f:circle-fourier}
\fentry{Feynman contraction}{f:ext-feynman}
\fentry{Fluctuation, subtracted difference}{f:relative-fluctuation}
\fentry{Fourier transform, convention}{f:fourier-pair}
\fentry{Gauge fixing}{f:green-homotopy}
\fentry{Gauss flux at a boundary}{f:gauge-boundary-ex}
\fentry{Geodesic and world function}{f:world-function}
\fentry{Germs and tangents}{f:tangent-coords}
\fentry{Global hyperbolicity}{f:cauchy-splitting}
\fentry{GNS representation}{f:gns}
\fentry{Green operators, compact and one-sided domains}{f:green-identities}
\fentry{Groupoid and stabilizer}{f:sign-groupoid}
\end{multicols}\clearpage
\subsection*{H--W}
\begin{multicols}{2}
\fcontinued{Hadamard state and parametrix}{f:wf-hadamard}{f:circle-subtraction}
\fentry{Hamiltonian vector and bracket sign}{f:poisson-convention}
\fentry{Hausdorff failure after identification}{f:q-double-origin}
\fentry{Hodge star and formal adjoint}{f:hodge}
\fentry{Holonomy}{f:parallel-ode}
\fentry{Hypersurface, regular value}{f:regular-value}
\fentry{Jet bundle and holonomic jet}{f:jet}
\fentry{Levi--Civita connection}{f:levi-civita}
\fentry{Lie algebra and representation}{f:group-action}
\fentry{Lindel\"of property and second countability}{f:q-lindelof}
\fentry{Momentum map}{f:regular-reduction}
\fentry{Normal Legendre branch}{f:normal-legendre}
\fentry{Null screen}{f:null-screen}
\fentry{Open set and neighborhood}{f:top-open}
\fentry{Orbit, stabilizer, transition signs}{f:sign-groupoid}
\fentry{Parallel transport}{f:parallel-ode}
\fentry{Partition of unity}{f:partition}
\fentry{Peierls bracket}{f:peierls-labels}
\fentry{Presymplectic kernel}{f:symplectic}
\fentry{Principal symbol}{f:wave-symbol}
\fentry{Propagation of singularities}{f:wf-propagation}
\fentry{Pullback of a distribution}{f:wf-pullback}
\fentry{Quasifree state}{f:gaussian-positivity}
\fentry{Quotient topology and descent}{f:q-descent}
\fentry{Radical and physical observables}{f:alternating-quotient}
\fentry{Scaling degree and ambiguity}{f:ext-theorem}
\fentry{Schur complement on histories}{f:delayed-schur}
\fentry{Sobolev spaces and duality}{f:sobolev-scale}
\fentry{Static response of a twisted field}{f:static-measurement}
\fentry{Stokes' theorem}{f:stokes}
\fentry{Tensor product and complexification}{f:tensor-product}
\fentry{Test functions}{f:test-topology}
\fcontinued{Trace, strong and weak}{f:plane-trace}{f:weak-flux}
\fentry{Transversality and clean intersection}{f:normal}
\fentry{Tree response}{f:ext-tree}
\fentry{Two-point function, mode construction}{f:circle-two-point}
\fentry{Van Vleck determinant}{f:van-vleck}
\fentry{Ward contact and consistency}{f:ext-ward}
\fentry{Wavefront set}{f:wf-test}
\fentry{Weyl algebra}{f:weyl-relations}
\fcontinued{Wick pairings and a two-line loop}{f:wick-pairings}{f:ext-fish}
\fentry{Wilson loop}{f:two-rotor}
\end{multicols}\endgroup

% !TeX root = main.tex
\renewcommand{\bibname}{References}
\begin{thebibliography}{99}\addcontentsline{toc}{chapter}{References}\markboth{References}{}
\bibitem{Munkres} J. R. Munkres, \emph{Topology}, 2nd ed., Prentice Hall (2000).
\bibitem{LeeSM} J. M. Lee, \emph{Introduction to Smooth Manifolds}, 2nd ed., Springer (2013).
\bibitem{LeeRM} J. M. Lee, \emph{Introduction to Riemannian Manifolds}, 2nd ed., Springer (2018).
\bibitem{KN} S. Kobayashi and K. Nomizu, \emph{Foundations of Differential Geometry}, vol. I, Interscience (1963).
\bibitem{ON} B. O'Neill, \emph{Semi-Riemannian Geometry with Applications to Relativity}, Academic Press (1983).
\bibitem{BS} A. N. Bernal and M. S\'anchez, ``Smoothness of time functions and the metric splitting of globally hyperbolic spacetimes,'' \emph{Commun. Math. Phys.} \textbf{257} (2005), 43--50.
\bibitem{GMS} G. Giachetta, L. Mangiarotti and G. Sardanashvily, \emph{Advanced Classical Field Theory}, World Scientific (2009).
\bibitem{Evans} L. C. Evans, \emph{Partial Differential Equations}, 2nd ed., American Mathematical Society (2010).
\bibitem{Taylor} M. E. Taylor, \emph{Partial Differential Equations I: Basic Theory}, 2nd ed., Springer (2011).
\bibitem{RS} M. Reed and B. Simon, \emph{Methods of Modern Mathematical Physics I: Functional Analysis}, revised ed., Academic Press (1980).
\bibitem{Schwartz} L. Schwartz, \emph{Th\'eorie des distributions}, Hermann (1966).
\bibitem{HormI} L. H\"ormander, \emph{The Analysis of Linear Partial Differential Operators I: Distribution Theory and Fourier Analysis}, 2nd ed., Springer (1990).
\bibitem{HormIII} L. H\"ormander, \emph{The Analysis of Linear Partial Differential Operators III: Pseudo-Differential Operators}, Springer (1985).
\bibitem{BGP} C. B\"ar, N. Ginoux and F. Pf\"affle, \emph{Wave Equations on Lorentzian Manifolds and Quantization}, European Mathematical Society (2007).
\bibitem{BW} C. B\"ar and R. Tagne Wafo, ``Initial value problems for wave equations on manifolds,'' \emph{Math. Phys. Anal. Geom.} \textbf{18} (2015), article 7.
\bibitem{Arnold} V. I. Arnold, \emph{Mathematical Methods of Classical Mechanics}, 2nd ed., Springer (1989).
\bibitem{Marsden} J. E. Marsden and T. S. Ratiu, \emph{Introduction to Mechanics and Symmetry}, 2nd ed., Springer (1999).
\bibitem{HW} D. Harlow and J.-q. Wu, ``Covariant phase space with boundaries,'' \emph{JHEP} \textbf{10} (2020), 146.
\bibitem{CMR} A. S. Cattaneo, P. Mnev and N. Reshetikhin, ``Classical BV theories on manifolds with boundary,'' \emph{Commun. Math. Phys.} \textbf{332} (2014), 535--603.
\bibitem{Khavkine} I. Khavkine, ``Covariant phase space, constraints, gauge and the Peierls formula,'' \emph{Int. J. Mod. Phys. A} \textbf{29} (2014), 1430009.
\bibitem{WaldQFT} R. M. Wald, \emph{Quantum Field Theory in Curved Spacetime and Black Hole Thermodynamics}, University of Chicago Press (1994).
\bibitem{Haag} R. Haag, \emph{Local Quantum Physics: Fields, Particles, Algebras}, 2nd ed., Springer (1996).
\bibitem{Yngvason} J. Yngvason, ``The role of type III factors in quantum field theory,'' \emph{Rep. Math. Phys.} \textbf{55} (2005), 135--147. \href{https://arxiv.org/abs/math-ph/0411058}{arXiv:math-ph/0411058}.
\bibitem{FewSplit} C. J. Fewster, ``The split property for quantum field theories in flat and curved spacetimes,'' \href{https://arxiv.org/abs/1601.06936}{arXiv:1601.06936} (2016).
\bibitem{BR} O. Bratteli and D. W. Robinson, \emph{Operator Algebras and Quantum Statistical Mechanics}, vols. 1 and 2, 2nd eds., Springer (1987, 1997).
\bibitem{Rad} M. J. Radzikowski, ``Micro-local approach to the Hadamard condition in quantum field theory on curved space-time,'' \emph{Commun. Math. Phys.} \textbf{179} (1996), 529--553.
\bibitem{SV} H. Sahlmann and R. Verch, ``Microlocal spectrum condition and Hadamard form for vector-valued quantum fields in curved spacetime,'' \emph{Rev. Math. Phys.} \textbf{13} (2001), 1203--1246. \href{https://arxiv.org/abs/math-ph/0008029}{arXiv:math-ph/0008029}.
\bibitem{BF} R. Brunetti and K. Fredenhagen, ``Microlocal analysis and interacting quantum field theories: renormalization on physical backgrounds,'' \emph{Commun. Math. Phys.} \textbf{208} (2000), 623--661.
\bibitem{HWick} S. Hollands and R. M. Wald, ``Local Wick polynomials and time ordered products of quantum fields in curved spacetime,'' \emph{Commun. Math. Phys.} \textbf{223} (2001), 289--326.
\bibitem{HWT} S. Hollands and R. M. Wald, ``Existence of local covariant time ordered products of quantum fields in curved spacetime,'' \emph{Commun. Math. Phys.} \textbf{231} (2002), 309--345.
\bibitem{FR} K. Rejzner, \emph{Perturbative Algebraic Quantum Field Theory}, Springer (2016).
\bibitem{DLMF} NIST, \emph{Digital Library of Mathematical Functions}, Ch.~10, ``Bessel Functions,'' \href{https://dlmf.nist.gov/10.31}{\S10.31: Power Series}.
\end{thebibliography}

\end{document}
